% La situation de la femme (en ville et au village) — Allemand Tle A — Cours signé OnBuch+
% Programme officiel MINESEC. Compiler avec : tectonic AL01-la-situation-de-la-femme-en-ville-et-au-village.tex
\newcommand{\DOCMATIERE}{Allemand}
\newcommand{\DOCNIVEAU}{Tle A}
\newcommand{\DOCMODULE}{Chapitre 1 : Vie familiale et sociale}
\newcommand{\DOCLECON}{AL01}
\newcommand{\DOCTITRE}{La situation de la femme (en ville et au village)}
% OnBuch+ — préambule « cours signés » ENRICHI (compilé avec tectonic / XeLaTeX)
% Système visuel « Pop » d'OnBuch+ : fond crème (jamais blanc), bordures encre
% épaisses, ombres « dures » décalées, titres Archivo Black en capitales.
% Version MATHÉMATIQUES : graphes (pgfplots/tikz), boîtes pédagogiques
% (définition, propriété, méthode, attention, le savais-tu, exercice résolu,
% exercice, TP, bilan), page de garde et signature OnBuch+ ancrée en bas.
%
% Macros attendues dans main.tex AVANT \input{preamble} :
%   \DOCMATIERE  \DOCNIVEAU  \DOCMODULE  \DOCLECON (numéro)  \DOCTITRE
\documentclass[11pt]{article}

\usepackage{fontspec}
\usepackage[ngerman,french]{babel} % français = langue principale ; l'allemand via \all{..} / textebox
\usepackage[a4paper,margin=2cm,top=3.4cm,bottom=2.8cm,headheight=1.4cm,footskip=1.3cm]{geometry}
\usepackage{amsmath,amssymb,amsfonts}
\usepackage{mathtools} % \xmapsto, \coloneqq... souvent utilisés par les modèles
\usepackage{cancel} % simplifications barrées (bilans de forces)
% \abs{z} (module, valeur absolue) et \norm{u} (norme), utilisés par les modèles
\providecommand{\abs}{}\let\abs\relax\DeclarePairedDelimiter{\abs}{\lvert}{\rvert}
\providecommand{\norm}{}\let\norm\relax\DeclarePairedDelimiter{\norm}{\lVert}{\rVert}
\usepackage[table]{xcolor} % [table] : fournit aussi \rowcolors (alternance de lignes)
\usepackage{pagecolor}
\usepackage{tikz}
\usetikzlibrary{arrows.meta,positioning,calc,shapes.geometric,shapes.misc,shapes.arrows,decorations.pathmorphing,decorations.pathreplacing,decorations.markings,patterns,shadows,mindmap,trees,fit,backgrounds,matrix,intersections,angles,quotes}
\usepackage{pgfplots}
\usepackage[version=4]{mhchem} % équations nucléaires \ce{^{235}_{92}U}
\usepackage[european,siunitx]{circuitikz} % schémas électriques
\pgfplotsset{compat=1.18}
\usepgfplotslibrary{fillbetween,groupplots} % aires entre courbes (calcul intégral)
\usepackage{enumitem}
\usepackage{fancyhdr}
% [most] charge toutes les bibliothèques tcolorbox (listings, minted, raster,
% documentation...) et gonfle inutilement la mémoire TeX sur un document long
% (~300 boîtes) : on ne charge que ce qui est réellement utilisé.
\usepackage[skins,breakable]{tcolorbox}
\usepackage{graphicx}
\usepackage{colortbl}
\usepackage{booktabs}
\usepackage{tabularx}
\usepackage{diagbox} % case d'en-tête diagonale des tableaux à double entrée
% Colonnes extensibles centrée / gauche / droite (C, L, R), souvent utilisées par les modèles
\newcolumntype{C}{>{\centering\arraybackslash}X}
\newcolumntype{L}{>{\raggedright\arraybackslash}X}
\newcolumntype{R}{>{\raggedleft\arraybackslash}X}
\usepackage{array}
\usepackage{multicol}
\usepackage{multirow}
\usepackage{siunitx}
% parse-numbers=false : accepte aussi \SI{2,5 \times 10^{-3}}{...}
\sisetup{locale=FR,output-decimal-marker={,},group-minimum-digits=4,parse-numbers=false}
\DeclareSIUnit\ohm{\ensuremath{\symup{\Omega}}} % U+2126 absent de la police
\DeclareSIUnit\division{div} % réglages d'oscilloscope (ms/div, V/div)
\DeclareSIUnit\tr{tr} % tours (vitesse de rotation, ex. \SI{1500}{\tr\per\minute})
\DeclareSIUnit\molar{M} % concentration molaire (ex. \SI{3}{\molar}, \SI{1}{\micro\molar})
\DeclareSIUnit\an{an} % année (le modèle confond parfois avec \year, un compteur TeX) — ex. \SI{5}{\kilo\meter\per\an}
\DeclareSIUnit\FCFA{FCFA} % franc CFA (ex. \SI{500}{\FCFA\per\kilo\gram})
\DeclareSIUnit\degreeBrix{\textdegree Brix} % teneur en sucre d'un sirop/jus (réfractométrie)
\DeclareSIUnit\month{mois} % le modèle utilise parfois \month, un compteur TeX, comme unité de durée
\usepackage{qrcode}
% \degree : commande fréquemment utilisée par les modèles pour un angle ou une
% température en dehors de \SI{}{} (non fournie par les paquets chargés).
\providecommand{\degree}{^{\circ}}
% \qed (fin de démonstration), utilisé par les modèles sans amsthm
\providecommand{\qed}{\hfill$\square$}
% \barre : double barre des valeurs interdites dans un tableau de variations (tkz-tab)
\providecommand{\barre}{\ensuremath{\|}}
% environnement proof (amsthm non chargé) utilisé par les modèles
\ifdefined\proof\else\newenvironment{proof}[1][Démonstration]{\par\noindent\textit{#1.}\ }{\qed\par}\fi
\usepackage{lastpage}
\usepackage{hyperref}
\hypersetup{hidelinks}

% ============================================================
% Palette « Pop » OnBuch+ — mêmes hex que la classe P (plus_ui.dart)
% ============================================================
\definecolor{poporange}{HTML}{F59321}
\definecolor{popdark}{HTML}{E07A0C}
\definecolor{popcream}{HTML}{FFF6E8}
\definecolor{popink}{HTML}{1C1714}
\definecolor{poppurple}{HTML}{7A5AE0}
\definecolor{popgreen}{HTML}{2E9E6A}
\definecolor{poppink}{HTML}{E0457B}
\definecolor{popblue}{HTML}{2F7DE1}
\definecolor{popgold}{HTML}{C98A12}
\definecolor{popmuted}{HTML}{8A7F76}
\definecolor{popline}{HTML}{EADFD2}
\definecolor{popgray}{HTML}{8A7F76}
\colorlet{popgrey}{popgray}
% Alias tolérants (le modèle nomme parfois les couleurs différemment)
\colorlet{popwhite}{white}
\colorlet{popblack}{popink}
\colorlet{popred}{poppink}
\colorlet{popyellow}{popgold}
% Teintes claires (fonds de tableaux/encadrés) — le modèle nomme parfois
% les couleurs avec un suffixe L (ex. popblueL) sans les avoir définies.
\colorlet{poporangeL}{poporange!20}
\colorlet{popdarkL}{popdark!20}
\colorlet{poppurpleL}{poppurple!20}
\colorlet{popgreenL}{popgreen!20}
\colorlet{poppinkL}{poppink!20}
\colorlet{popblueL}{popblue!20}
\colorlet{popgoldL}{popgold!20}
\colorlet{popredL}{popred!20}
\colorlet{popgrayL}{popgray!20}
% Teintes claires pour fonds de tableaux / zones
\colorlet{popblueL}{popblue!10!white}
\colorlet{poporangeL}{poporange!14!white}
\colorlet{popgreenL}{popgreen!12!white}
\colorlet{poppurpleL}{poppurple!12!white}
\colorlet{poppinkL}{poppink!12!white}
\colorlet{popgoldL}{popgold!14!white}

\pagecolor{popcream}

% ============================================================
% Typographie
% ============================================================
\defaultfontfeatures{Ligatures=TeX}
\setmainfont{PlusJakartaSans-Medium.ttf}[Path=../../../fonts/,BoldFont=PlusJakartaSans-Bold.ttf]
% Maths en sans-serif (Fira Math) avec chiffres et lettres droites de Plus
% Jakarta, pour que les formules se fondent dans le texte.
\usepackage[math-style=ISO,bold-style=ISO]{unicode-math}
\setmathfont{FiraMath-Regular.otf}[Path=../../../fonts/]
\setmathfont{PlusJakartaSans-Medium.ttf}[Path=../../../fonts/,range={up/num,up/{Latin,latin},\mathup}]
\usepackage{icomma} % virgule décimale sans espace en mode math
% Caractères mathématiques Unicode absents des polices de texte (Plus Jakarta,
% Archivo Black) : rendus par leur équivalent mathématique, en texte comme en
% formule — sinon ils s'affichent en carré vide (ex. « Arithmétique dans ℤ »).
\usepackage{newunicodechar}
\newunicodechar{ℝ}{\ensuremath{\mathbb{R}}}
\newunicodechar{ℤ}{\ensuremath{\mathbb{Z}}}
\newunicodechar{ℂ}{\ensuremath{\mathbb{C}}}
\newunicodechar{ℕ}{\ensuremath{\mathbb{N}}}
\newunicodechar{ℚ}{\ensuremath{\mathbb{Q}}}
\newunicodechar{𝔻}{\ensuremath{\mathbb{D}}}
\newunicodechar{α}{\ensuremath{\alpha}}
\newunicodechar{β}{\ensuremath{\beta}}
\newunicodechar{γ}{\ensuremath{\gamma}}
\newunicodechar{δ}{\ensuremath{\delta}}
\newunicodechar{ε}{\ensuremath{\varepsilon}}
\newunicodechar{θ}{\ensuremath{\theta}}
\newunicodechar{λ}{\ensuremath{\lambda}}
\newunicodechar{μ}{\ensuremath{\mu}}
\newunicodechar{σ}{\ensuremath{\sigma}}
\newunicodechar{τ}{\ensuremath{\tau}}
\newunicodechar{φ}{\ensuremath{\varphi}}
\newunicodechar{ψ}{\ensuremath{\psi}}
\newunicodechar{ω}{\ensuremath{\omega}}
\newunicodechar{Σ}{\ensuremath{\Sigma}}
% majuscules produites par \MakeUppercase dans les titres (α-aminés -> Α)
\newunicodechar{Α}{\ensuremath{\alpha}}
\newunicodechar{Β}{\ensuremath{\beta}}
\newunicodechar{Γ}{\ensuremath{\Gamma}}
\newunicodechar{Δ}{\ensuremath{\Delta}}
\newunicodechar{Ε}{\ensuremath{\varepsilon}}
\newunicodechar{Θ}{\ensuremath{\Theta}}
\newunicodechar{Λ}{\ensuremath{\Lambda}}
\newunicodechar{Μ}{\ensuremath{\mu}}
\newunicodechar{Τ}{\ensuremath{\tau}}
\newunicodechar{Φ}{\ensuremath{\Phi}}
\newunicodechar{Ψ}{\ensuremath{\Psi}}
\newunicodechar{Ω}{\ensuremath{\Omega}}
\newunicodechar{↦}{\ensuremath{\mapsto}}
\newunicodechar{∈}{\ensuremath{\in}}
\newunicodechar{∉}{\ensuremath{\notin}}
\newunicodechar{∧}{\ensuremath{\wedge}}
\newunicodechar{⇔}{\ensuremath{\Leftrightarrow}}
\newunicodechar{⟺}{\ensuremath{\Longleftrightarrow}}
\newunicodechar{⇒}{\ensuremath{\Rightarrow}}
\newunicodechar{⟹}{\ensuremath{\Longrightarrow}}
\newunicodechar{∘}{\ensuremath{\circ}}
\newunicodechar{∪}{\ensuremath{\cup}}
\newunicodechar{∩}{\ensuremath{\cap}}
\newunicodechar{≡}{\ensuremath{\equiv}}
\newunicodechar{⊂}{\ensuremath{\subset}}
\newunicodechar{⊥}{\ensuremath{\perp}}
\newunicodechar{∃}{\ensuremath{\exists}}
\newunicodechar{∀}{\ensuremath{\forall}}
\newunicodechar{∛}{\ensuremath{\sqrt[3]{\,}}}
\newunicodechar{∜}{\ensuremath{\sqrt[4]{\,}}}
\newunicodechar{⃗}{\ensuremath{{}^{\to}}}
\newunicodechar{ⁿ}{\ifmmode^{n}\else\textsuperscript{\ensuremath{n}}\fi}
\newunicodechar{ˣ}{\ifmmode^{x}\else\textsuperscript{\ensuremath{x}}\fi}
\newunicodechar{ᵇ}{\ifmmode^{b}\else\textsuperscript{\ensuremath{b}}\fi}
\newunicodechar{ʳ}{\ifmmode^{r}\else\textsuperscript{\ensuremath{r}}\fi}
\newunicodechar{ᵘ}{\ifmmode^{u}\else\textsuperscript{\ensuremath{u}}\fi}
\newunicodechar{ʸ}{\ifmmode^{y}\else\textsuperscript{\ensuremath{y}}\fi}
\newunicodechar{ᵃ}{\ifmmode^{a}\else\textsuperscript{\ensuremath{a}}\fi}
\newunicodechar{ᵉ}{\ifmmode^{e}\else\textsuperscript{\ensuremath{e}}\fi}
\newunicodechar{ᵏ}{\ifmmode^{k}\else\textsuperscript{\ensuremath{k}}\fi}
\newunicodechar{ᵖ}{\ifmmode^{p}\else\textsuperscript{\ensuremath{p}}\fi}
\newunicodechar{ᴺ}{\ifmmode^{N}\else\textsuperscript{\ensuremath{N}}\fi}
\newunicodechar{ᶜ}{\ifmmode^{c}\else\textsuperscript{\ensuremath{c}}\fi}
\newunicodechar{ʰ}{\ifmmode^{h}\else\textsuperscript{\ensuremath{h}}\fi}
\newunicodechar{ᵠ}{\ifmmode^{q}\else\textsuperscript{\ensuremath{q}}\fi}
\newunicodechar{ₙ}{\ifmmode_{n}\else\textsubscript{\ensuremath{n}}\fi}
\newunicodechar{ₐ}{\ifmmode_{a}\else\textsubscript{\ensuremath{a}}\fi}
\newunicodechar{ᵢ}{\ifmmode_{i}\else\textsubscript{\ensuremath{i}}\fi}
\newunicodechar{ᵣ}{\ifmmode_{r}\else\textsubscript{\ensuremath{r}}\fi}
\newunicodechar{ₖ}{\ifmmode_{k}\else\textsubscript{\ensuremath{k}}\fi}
\newunicodechar{ᵦ}{\ifmmode_{\beta}\else\textsubscript{\ensuremath{\beta}}\fi}
\newunicodechar{ₓ}{\ifmmode_{x}\else\textsubscript{\ensuremath{x}}\fi}
\newfontfamily\popdisplay{ArchivoBlack-Regular.ttf}[Path=../../../fonts/]
\newfontfamily\poplogo{SpaceGrotesk-Bold.ttf}[Path=../../../fonts/]
\newfontfamily\popsemi{PlusJakartaSans-SemiBold.ttf}[Path=../../../fonts/]
\newfontfamily\popxbold{PlusJakartaSans-ExtraBold.ttf}[Path=../../../fonts/]
\newfontfamily\popxboldls{PlusJakartaSans-ExtraBold.ttf}[Path=../../../fonts/,LetterSpace=3.0]

\setlist[enumerate]{leftmargin=1.7em,itemsep=5pt,topsep=4pt}
\setlist[itemize]{leftmargin=1.7em,itemsep=4pt,topsep=4pt,label={\color{poporange}\textbullet}}
\setlength{\parindent}{0pt}
\setlength{\parskip}{5pt}
\renewcommand{\arraystretch}{1.25}

% pgfplots : style Pop par défaut
\pgfplotsset{
  every axis/.append style={axis line style={popink,line width=1pt},tick style={popink},
    grid style={popline,line width=0.5pt},label style={font=\small},tick label style={font=\footnotesize}},
  popcurve/.style={poporange,line width=1.8pt,no markers,smooth},
}
% Même style disponible dans un \draw TikZ simple (hors pgfplots)
\tikzset{popcurve/.style={poporange,line width=1.8pt}}
\tikzset{popgrid/.style={popline,very thin}}
\colorlet{popcurve}{poporange} % le nom du style est parfois utilisé comme couleur

% ============================================================
% Pill / squiggle
% ============================================================
\newcommand{\poppill}[3]{%
  \tikz[baseline=-0.55ex]\node[draw=popink,line width=1.2pt,rounded corners=2.6pt,%
    fill=#2,inner xsep=6.5pt,inner ysep=3pt]{%
    \popxboldls\fontsize{7}{7}\selectfont\color{#3}\MakeUppercase{#1}};%
}
\newcommand{\popsquiggle}[1]{%
  \tikz[baseline=-0.4ex]{
    \draw[#1,line width=2.6pt,line cap=round]
      plot [smooth] coordinates {(0,0.09) (0.34,0.01) (0.68,0.21) (1.02,0.01) (1.36,0.10)};
  }%
}
% Mot-clé surligné (marqueur orange sous le texte)
\newcommand{\cle}[1]{{\popxbold\color{popdark}#1}}

% ============================================================
% En-tête / pied
% ============================================================
\pagestyle{fancy}
\fancyhf{}
\renewcommand{\headrulewidth}{0pt}
\renewcommand{\footrulewidth}{0pt}
\lhead{\raisebox{-0.15ex}{\poplogo\fontsize{13}{13}\selectfont\color{popink}OnBuch\color{poporange}+}}
\rhead{\poppill{\DOCMATIERE\ \textbullet\ \DOCNIVEAU\ \textbullet\ Leçon \DOCLECON}{white}{popink}}
\lfoot{\popxbold\fontsize{7.6}{7.6}\selectfont\color{popmuted}\MakeUppercase{Cours signé OnBuch+}\ \textbullet\ {\popsemi\DOCTITRE}}
\rfoot{\tikz[baseline=-0.6ex]\node[rounded corners=8.5pt,fill=popink,minimum height=17pt,minimum width=17pt,inner xsep=5pt,inner sep=0]{\popxbold\fontsize{8}{8}\selectfont\color{popcream}\thepage\,/\,\pageref*{LastPage}};}

% ============================================================
% Titres
% ============================================================
\newcommand{\chaptitre}[2]{%
  \begin{tcolorbox}[enhanced,colback=poporange,colframe=popink,boxrule=2.6pt,arc=11pt,
      drop shadow=popink,left=15pt,right=15pt,top=12pt,bottom=11pt,before skip=3pt,after skip=18pt]
    \poppill{#2}{popink}{popcream}\par\vspace{9pt}
    {\raggedright\hyphenpenalty=10000\exhyphenpenalty=10000\popdisplay\fontsize{22}{24}\selectfont\color{popcream}\MakeUppercase{#1}\par}
  \end{tcolorbox}
}
% Section numérotée : pastille orange avec le numéro + titre Archivo
\newcounter{popsec}
\newcommand{\coursec}[1]{%
  \stepcounter{popsec}\setcounter{popsub}{0}%
  \par\vspace{16pt}\noindent
  \begin{minipage}{\linewidth}
  \tikz[baseline=-0.7ex]\node[draw=popink,line width=1.6pt,fill=poporange,rounded corners=4pt,minimum size=22pt,inner sep=0]{\popdisplay\fontsize{12}{12}\selectfont\color{popcream}\thepopsec};%
  \hspace{8pt}{\popdisplay\fontsize{15}{17}\selectfont\color{popink}\MakeUppercase{#1}}\par
  \vspace{4pt}\noindent{\color{poporange}\rule{36pt}{3.6pt}}
  \end{minipage}\par\nopagebreak\vspace{8pt}
}
\newcounter{popsub}
\newcommand{\courssub}[1]{%
  \stepcounter{popsub}%
  \par\vspace{9pt}\noindent{\popxbold\fontsize{11.5}{13}\selectfont\color{popink}{\color{poporange}\thepopsec.\thepopsub}\hspace{6pt}#1}\par\nopagebreak\vspace{4pt}
}

% ============================================================
% Boîtes pédagogiques — même recette : fond blanc, bordure épaisse colorée,
% ombre dure encre, badge plein. Titre optionnel [#1].
% ============================================================
\tcbset{popbase/.style={breakable,enhanced,colback=white,boxrule=2.3pt,arc=9pt,
  shadow={3pt}{-3pt}{0pt}{popink},left=14pt,right=14pt,top=11pt,bottom=11pt,before skip=11pt,after skip=15pt,
  fontupper=\popsemi\fontsize{10.6}{14.5}\selectfont\color{popink},
  fontlower=\popsemi\fontsize{10.6}{14.5}\selectfont\color{popink}}}
% Coche dessinée (le glyphe ✓ n'existe pas dans les polices du document)
\newcommand{\popcheck}[1]{\tikz[baseline=-0.5ex]\draw[#1,line width=1.6pt,line cap=round,line join=round](-0.13,0.0)--(-0.04,-0.09)--(0.13,0.1);}
\providecommand{\coche}{\popcheck{popgreen}}
\providecommand{\resultat}[1]{\par\smallskip\noindent\fcolorbox{popgreen}{popgreenL}{\parbox{\dimexpr\linewidth-2\fboxsep-2\fboxrule\relax}{\textbf{Résultat.} #1}}\par\smallskip}
\newcommand{\popboxhead}[3]{\poppill{#1}{#2}{white}\ifx\relax#3\relax\else\hspace{6pt}{\popxbold\fontsize{10.6}{12}\selectfont\color{popink}#3}\fi\par\vspace{7pt}}

\newtcolorbox{aretenir}[1][]{popbase,colframe=popgreen,before upper={\popboxhead{À retenir}{popgreen}{#1}}}
% Texte en allemand (document de lecture, dialogue, extrait) : cadre
% d'étude. Un titre optionnel (source, auteur) s'affiche dans l'en-tête.
\newtcolorbox{textebox}[1][]{popbase,colframe=popblue,before upper={\popboxhead{Text}{popblue}{#1}\begin{otherlanguage}{ngerman}},after upper={\end{otherlanguage}}}
% Marques ✓ / ✗ (forme correcte / fautive) souvent utilisées par les modèles
\providecommand{\cross}{\textcolor{poppink}{\ensuremath{\times}}}
\providecommand{\xmark}{\cross}\providecommand{\crossmark}{\cross}\providecommand{\cmark}{\ensuremath{\checkmark}}
% Ligne de réponse / justification à compléter (inventée par les modèles)
\providecommand{\justification}[1]{\par\noindent\textit{Justification~:} \dotfill\par}
% \textipa (package tipa non chargé) : la police ne couvre pas l'API -> simple passage
\providecommand{\textipa}[1]{#1}
% Soulignement (ulem non chargé) : \uline, \ul
\providecommand{\uline}[1]{\underline{#1}}\providecommand{\ul}[1]{\underline{#1}}
% \alert (beamer) inventée par les modèles : texte mis en évidence
\providecommand{\alert}[1]{\textbf{\textcolor{poppink}{#1}}}
% variantes ASCII d'ordinaux écrites en macro
\providecommand{\iere}{\textsuperscript{re}}\providecommand{\ieme}{\textsuperscript{e}}\providecommand{\ere}{\textsuperscript{re}}\providecommand{\eme}{\textsuperscript{e}}\providecommand{\er}{\textsuperscript{er}}\providecommand{\ie}{\textsuperscript{e}}
% Ordinaux français écrits en macro par les modèles : 1\ière, 3\ième
\newcommand{\ière}{\textsuperscript{re}}\newcommand{\ième}{\textsuperscript{e}}
% \exemple (inventée par les modèles) : étiquette « Exemple : »
\providecommand{\exemple}{\textbf{Exemple~:}\ }
% \square utilisable aussi hors mode math (cases à cocher)
\AtBeginDocument{\let\squareMath\square\renewcommand{\square}{\ensuremath{\squareMath}}}
% Mot ou phrase en allemand à l'intérieur d'un paragraphe français
\newcommand{\all}[1]{\ifmmode\text{\textit{#1}}\else\begingroup\selectlanguage{ngerman}\itshape #1\endgroup\fi}
\let\esp\all % alias toléré
\newtcolorbox{exemplebox}[1][]{popbase,colframe=poporange,before upper={\popboxhead{Exemple}{poporange}{#1}}}
\newtcolorbox{definition}[1][]{popbase,colframe=popblue,before upper={\popboxhead{Définition}{popblue}{#1}}}
\newtcolorbox{propriete}[1][]{popbase,colframe=poppurple,before upper={\popboxhead{Propriété}{poppurple}{#1}}}
\newtcolorbox{methode}[1][]{popbase,colframe=popgold,before upper={\popboxhead{Méthode}{popgold}{#1}}}
\newtcolorbox{attention}[1][]{popbase,colframe=poppink,before upper={\popboxhead{Attention}{poppink}{#1}}}
\newtcolorbox{experience}[1][]{popbase,colframe=popblue,colback=popblueL,before upper={\popboxhead{Expérience}{popblue}{#1}}}
% « Le savais-tu ? » : carte violette pleine (contexte, culture, Cameroun)
\newtcolorbox{savaistu}[1][]{popbase,colback=poppurple,colframe=popink,
  fontupper=\popsemi\fontsize{10.4}{14.2}\selectfont\color{white},
  before upper={\poppill{Le savais-tu\,?}{white}{poppurple}\ifx\relax#1\relax\else\hspace{6pt}{\popxbold\color{white}#1}\fi\par\vspace{7pt}}}
% Exercice résolu : énoncé en haut, \tcblower puis la solution sur fond crème
\newcounter{popexo}\newcounter{popexr}
\newtcolorbox{exoresolu}[1][]{popbase,colframe=poporange,colbacklower=popcream,
  segmentation style={popink,line width=1.2pt,dashed},
  before upper={\stepcounter{popexr}\popboxhead{Exercice résolu \thepopexr}{poporange}{#1}},
  before lower={\poppill{Solution}{popink}{popcream}\par\vspace{6pt}}}
% Exercice (non résolu en place) — la correction est dans la partie Corrigés
\newtcolorbox{exercice}[1][]{popbase,colframe=popink,
  before upper={\stepcounter{popexo}\popboxhead{Exercice \thepopexo}{popink}{#1}}}
% Corrigé
\newtcolorbox{corrige}[1][]{popbase,colframe=popgreen,colback=popgreenL,
  before upper={\popboxhead{Corrigé}{popgreen}{#1}}}
% Objectifs / prérequis / bilan
\newtcolorbox{objectifs}{popbase,colframe=popink,colback=poporangeL,
  before upper={\popboxhead{Objectifs de la leçon}{popink}{}}}
\newtcolorbox{prerequis}{popbase,colframe=popink,colback=popblueL,
  before upper={\popboxhead{Prérequis}{popblue}{}}}
\newtcolorbox{bilan}{popbase,colframe=popink,colback=white,boxrule=2.8pt,
  before upper={\popboxhead{Fiche bilan}{poporange}{L'essentiel à connaître pour l'examen}}}
% Figure encadrée avec légende
\newtcolorbox{popfigure}[1][]{enhanced,breakable=false,colback=white,colframe=popline,boxrule=1.4pt,arc=8pt,
  left=10pt,right=10pt,top=10pt,bottom=8pt,before skip=10pt,after skip=12pt,halign=center,
  lower separated=false,halign lower=center,
  fontlower=\popsemi\fontsize{9}{11.5}\selectfont\color{popmuted},
  #1}
\newcommand{\legende}[1]{\par\vspace{6pt}{\popsemi\fontsize{9}{11.5}\selectfont\color{popmuted}#1}}

% ============================================================
% Signature OnBuch+
% ============================================================
\newcommand{\poplogoinline}[1]{{\poplogo\fontsize{#1}{#1}\selectfont\color{popink}OnBuch\color{poporange}+}}
\newcommand{\signatureonbuch}{%
  \begin{tcolorbox}[enhanced,colback=popink,colframe=popink,boxrule=2.2pt,arc=10pt,
      drop shadow=poporange,left=14pt,right=14pt,top=10pt,bottom=10pt,before skip=0pt,after skip=0pt]
    \tikz[baseline=-0.7ex]\node[circle,fill=popgreen,draw=popcream,line width=1.3pt,minimum size=22pt,inner sep=0]{\popcheck{white}};%
    \hspace{9pt}\begin{minipage}[c]{0.62\linewidth}
      {\popxbold\fontsize{12}{13}\selectfont\color{popcream}Signé\ \textbullet\ {\poplogo OnBuch\color{poporange}+}}\par\vspace{2pt}
      {\raggedright\popsemi\fontsize{8}{10.5}\selectfont\color{popline}\MakeUppercase{Équipe pédagogique OnBuch+}\\\MakeUppercase{Conforme au programme officiel MINESEC}\par}
    \end{minipage}\hfill
    {\poplogo\fontsize{22}{22}\selectfont\color{popcream}OnBuch\color{poporange}+}
  \end{tcolorbox}%
}

% ============================================================
% Page de garde
% #1 titre · #2 matière · #3 niveau · #4 module · #5 n° leçon · #6 durée ·
% #7 description / objectifs (texte ou itemize)
% ============================================================
\newcommand{\pagegarde}[7]{%
  \thispagestyle{empty}
  \newgeometry{a4paper,margin=1.7cm,top=1.6cm,bottom=1.4cm}%
  \begin{tikzpicture}[remember picture,overlay]
    \fill[poporange!18] ([xshift=-3.2cm,yshift=-3.6cm]current page.north east) circle (4.6cm);
    \fill[poppurple!14] ([xshift=2.4cm,yshift=7.6cm]current page.south west) circle (3.8cm);
  \end{tikzpicture}%
  {\poplogo\fontsize{30}{30}\selectfont\color{popink}OnBuch\color{poporange}+}\hfill
  \raisebox{0.6ex}{\poppill{Cours signé \textbullet\ Premium}{popink}{popcream}}\par
  \vspace{1pt}\popsquiggle{poppurple}\par
  \vspace{8pt}
  \begin{tcolorbox}[enhanced,colback=poporange,colframe=popink,boxrule=2.8pt,arc=13pt,
      drop shadow=popink,left=18pt,right=18pt,top=12pt,bottom=13pt,after skip=10pt]
    \poppill{#2 \textbullet\ #3}{popink}{popcream}\hspace{5pt}\poppill{#4}{white}{popink}\par\vspace{8pt}
    {\popdisplay\fontsize{14}{14}\selectfont\color{popink}LEÇON #5}\par\vspace{4pt}
    {\raggedright\hyphenpenalty=10000\exhyphenpenalty=10000\popdisplay\fontsize{25}{27}\selectfont\color{popcream}\MakeUppercase{#1}\par}
    \vspace{8pt}
    {\popxbold\fontsize{9.5}{9.5}\selectfont\color{popink}\MakeUppercase{Durée conseillée : #6}}
  \end{tcolorbox}
  \begin{tcolorbox}[enhanced,colback=white,colframe=popink,boxrule=2.2pt,arc=10pt,
      drop shadow=popink,left=16pt,right=16pt,top=10pt,bottom=10pt,after skip=8pt]
    \poppill{Dans cette leçon}{poppurple}{white}\par\vspace{5pt}
    \popsemi\fontsize{9.3}{12.6}\selectfont\color{popink}#7
  \end{tcolorbox}
  \noindent\begin{minipage}[t]{0.487\linewidth}
    \begin{tcolorbox}[enhanced,colback=popgreen,colframe=popink,boxrule=2.2pt,arc=10pt,
        drop shadow=popink,left=11pt,right=11pt,top=10pt,bottom=10pt,height=3.1cm,valign=center]
      \tcbox[colback=white,colframe=popink,boxrule=1.3pt,arc=5pt,boxsep=0pt,nobeforeafter,
        left=3pt,right=3pt,top=3pt,bottom=3pt,valign=center]{\qrcode[height=1.9cm]{https://play.google.com/store/apps/details?id=com.luvvix.onbuch&hl=fr}}%
      \hspace{8pt}\begin{minipage}[c]{0.5\linewidth}
        {\popdisplay\fontsize{11}{12.5}\selectfont\color{white}\MakeUppercase{Télécharge OnBuch}}\par\vspace{4pt}
        {\raggedright\popsemi\fontsize{8.4}{11}\selectfont\color{white}Gratuit \textbullet\ Google Play\\Tuteur IA, annales, concours, résultats\par}
      \end{minipage}
    \end{tcolorbox}
  \end{minipage}\hfill
  \begin{minipage}[t]{0.487\linewidth}
    \begin{tcolorbox}[enhanced,colback=poppurple,colframe=popink,boxrule=2.2pt,arc=10pt,
        drop shadow=popink,left=11pt,right=11pt,top=10pt,bottom=10pt,height=3.1cm,valign=center]
      \tikz[baseline=-0.7ex]\node[circle,fill=white,draw=popink,line width=1.3pt,minimum size=1.6cm,inner sep=0]{\popdisplay\fontsize{24}{24}\selectfont\color{poppurple}+};%
      \hspace{8pt}\begin{minipage}[c]{0.6\linewidth}
        {\popdisplay\fontsize{11}{12.5}\selectfont\color{white}\MakeUppercase{Passe à OnBuch+}}\par\vspace{4pt}
        {\raggedright\popsemi\fontsize{8.4}{11}\selectfont\color{white}Tous les cours signés, Tuteur IA illimité, fiches PDF — onglet OnBuch+ de l'app\par}
      \end{minipage}
    \end{tcolorbox}
  \end{minipage}
  \vspace{8pt}
  \popcontenu
  \vfill
  \signatureonbuch
  \restoregeometry
  \newpage
}

% Rangée « Ce cours contient » (page de garde)
\newcommand{\poptile}[3]{% #1 couleur #2 gros texte #3 légende
  \begin{tcolorbox}[enhanced,colback=#1,colframe=popink,boxrule=2pt,arc=9pt,shadow={2.5pt}{-2.5pt}{0pt}{popink},
      left=6pt,right=6pt,top=9pt,bottom=9pt,halign=center,valign=center,height=1.75cm,width=\linewidth,nobeforeafter]
    {\popdisplay\fontsize{12.5}{13}\selectfont\color{white}\MakeUppercase{#2}}\par\vspace{3pt}
    {\centering\popsemi\fontsize{7.8}{9.5}\selectfont\color{white}#3\par}
  \end{tcolorbox}}
\newcommand{\popcontenu}{%
  \par\noindent{\popxboldls\fontsize{7.5}{7.5}\selectfont\color{popmuted}CE COURS SIGNÉ CONTIENT}\par\vspace{7pt}
  \noindent\begin{minipage}[t]{0.235\linewidth}\poptile{popblue}{Cours}{complet, illustré, pas à pas}\end{minipage}\hfill
  \begin{minipage}[t]{0.235\linewidth}\poptile{poporange}{Méthodes}{et exercices résolus}\end{minipage}\hfill
  \begin{minipage}[t]{0.235\linewidth}\poptile{poppink}{Exercices}{type examen + corrigés}\end{minipage}\hfill
  \begin{minipage}[t]{0.235\linewidth}\poptile{popgreen}{Fiche bilan}{l'essentiel à retenir}\end{minipage}\par}

% Sommaire
\newtcolorbox{sommaire}{enhanced,colback=white,colframe=popink,boxrule=2.2pt,arc=10pt,
  drop shadow=popink,left=18pt,right=18pt,top=13pt,bottom=13pt,before skip=6pt,after skip=20pt,
  before upper={\poppill{Sommaire}{poppurple}{white}\par\vspace{9pt}\popsemi\fontsize{10.6}{14.5}\selectfont\color{popink}}}

% Signature de fin de cours — poussée tout en bas de la dernière page
\newcommand{\findecours}{%
  \par\vfill\nopagebreak
  \begin{center}\popsquiggle{poporange}\end{center}\vspace{4pt}
  \signatureonbuch
}
\AtBeginDocument{\def\llbracket{\mathopen{[\mkern-3.5mu[}}\def\rrbracket{\mathclose{]\mkern-3.5mu]}}} % intervalles d'entiers (glyphe ⟦ absent de la police)
% Glyphes absents de Fira Math : redéfinis à partir de glyphes disponibles
\AtBeginDocument{%
  \def\setminus{\mathbin{\backslash}}\let\smallsetminus\setminus
  \def\vdots{\mathord{\vbox{\baselineskip3.2pt\lineskiplimit0pt\kern4pt\hbox{.}\hbox{.}\hbox{.}}}}%
  \def\ddots{\mathinner{\mkern1mu\raise6.5pt\hbox{.}\mkern2mu\raise3.6pt\hbox{.}\mkern2mu\raise0.7pt\hbox{.}\mkern1mu}}%
  \def\triangle{\mathord{\scalebox{0.9}{$\symup{\Delta}$}}}%
  \def\nabla{\mathord{\raisebox{\depth}{\scalebox{1}[-1]{$\symup{\Delta}$}}}}%
  \def\checkmark{\popcheck{popdark}}%
  \def\star{\mathbin{\ast}}\def\bigstar{\mathord{\ast}}%
  \def\diamond{\mathbin{\scalebox{0.6}{$\Diamond$}}}%
  \def\bigcup{\mathop{\mathchoice{\raisebox{-0.3em}{\scalebox{1.7}{$\cup$}}}{\scalebox{1.25}{$\cup$}}{\cup}{\cup}}\displaylimits}%
  \def\bigcap{\mathop{\mathchoice{\raisebox{-0.3em}{\scalebox{1.7}{$\cap$}}}{\scalebox{1.25}{$\cap$}}{\cap}{\cap}}\displaylimits}%
  \def\lesseqqgtr{\mathrel{\lessgtr}}\def\bigoplus{\mathop{\oplus}\displaylimits}%
}
\providecommand{\ssi}{si et seulement si} % abréviation fréquente
\AtBeginDocument{\providecommand{\wideparen}{\overparen}} % arc de cercle

\begin{document}

\pagegarde{La situation de la femme (en ville et au village)}{Allemand}{Tle A}{Chapitre 1 : Vie familiale et sociale}{AL01}{2 h}{Cette leçon propose l'étude d'un texte sur la situation de la femme en ville et au village, dans les pays germanophones et au Cameroun. Elle permet d'acquérir le vocabulaire thématique (égalité, éducation, foyer, travail), de consolider l'ordre des mots, les connecteurs logiques et la conjugaison au présent, et de produire un court commentaire écrit et oral.
\vspace{4pt}
\begin{multicols}{2}\small
\begin{itemize}
\item À la fin de cette leçon, je sais comprendre globalement puis en détail un texte allemand sur la situation de la femme en ville et au village.
\item À la fin de cette leçon, je sais employer correctement le vocabulaire thématique (die Frau, die Gleichberechtigung, die Schulbildung, der Haushalt, der Beruf, die Arbeit auf dem Feld) avec articles et pluriels.
\item À la fin de cette leçon, je sais appliquer la règle du verbe en 2e position dans la phrase principale et du verbe en fin de subordonnée.
\item À la fin de cette leçon, je sais utiliser correctement les connecteurs weil, denn, deshalb et obwohl pour exprimer la cause et la concession.
\item À la fin de cette leçon, je sais conjuguer au présent les verbes faibles, les verbes forts et les verbes de modalité.
\item À la fin de cette leçon, je sais répondre à des questions de compréhension par des phrases complètes en allemand.
\end{itemize}\end{multicols}}

\begin{sommaire}
\begin{multicols}{2}
\begin{enumerate}
\item Texte support et première compréhension
\item Lexique thématique : la femme, la famille, le travail
\item Grammaire 1 : l'ordre des mots
\item Grammaire 2 : les connecteurs weil, denn, deshalb, obwohl
\item Conjugaison : présent et verbes de modalité
\item Méthode du commentaire et production guidée
\item Activité d'intégration
\item Méthodes et exercices résolus
\item Exercices
\item Corrigés des exercices
\item Fiche bilan
\end{enumerate}
\end{multicols}
\end{sommaire}

\begin{objectifs}
\begin{itemize}
\item À la fin de cette leçon, je sais comprendre globalement puis en détail un texte allemand sur la situation de la femme en ville et au village.
\item À la fin de cette leçon, je sais employer correctement le vocabulaire thématique (die Frau, die Gleichberechtigung, die Schulbildung, der Haushalt, der Beruf, die Arbeit auf dem Feld) avec articles et pluriels.
\item À la fin de cette leçon, je sais appliquer la règle du verbe en 2e position dans la phrase principale et du verbe en fin de subordonnée.
\item À la fin de cette leçon, je sais utiliser correctement les connecteurs weil, denn, deshalb et obwohl pour exprimer la cause et la concession.
\item À la fin de cette leçon, je sais conjuguer au présent les verbes faibles, les verbes forts et les verbes de modalité.
\item À la fin de cette leçon, je sais répondre à des questions de compréhension par des phrases complètes en allemand.
\item À la fin de cette leçon, je sais résumer un texte et rédiger un court paragraphe argumenté sur la place de la femme.
\item À la fin de cette leçon, je sais présenter oralement en quelques phrases la situation de la femme dans mon environnement.
\end{itemize}
\end{objectifs}

\begin{prerequis}
\begin{itemize}
\item Conjugaison au présent des verbes réguliers et des auxiliaires haben/sein (classe de Première)
\item Ordre des mots de base dans la phrase déclarative (sujet – verbe – compléments)
\item Notions sur les verbes de modalité (können, müssen, wollen, dürfen, sollen) vues en Première
\item Vocabulaire général de la famille et de la vie quotidienne (die Familie, das Kind, das Haus)
\item Stratégies de compréhension écrite : repérage du thème, des mots-clés, des mots-outils
\end{itemize}
\end{prerequis}

\newpage

\coursec{Texte support et première compréhension}

\courssub{Situation d'accroche et problématique}

Au marché de Mokolo à Yaoundé, une mère vend ses légumes le matin, prépare le repas le midi et aide ses enfants à faire leurs devoirs le soir : elle cumule trois journées en une. En Allemagne aussi, on débat de la place de la femme : faut-il choisir entre \all{Beruf} (métier) et \all{Familie} ? Le mot allemand \all{Rabenmutter} (littéralement « mère corbeau »), qui désigne une mère jugée trop absente de son foyer, montre que le débat existe aussi dans les pays germanophones. Comment la situation de la femme diffère-t-elle entre la ville et le village, au Cameroun comme en Allemagne ? C'est ce que nous allons découvrir à travers un texte, son vocabulaire et les outils grammaticaux pour en parler.

\begin{aretenir}[Problématique de la leçon]
\all{Wie unterscheidet sich die Situation der Frau in der Stadt und auf dem Land?} — En quoi la situation de la femme diffère-t-elle entre la ville et la campagne ? Pour répondre à cette question, nous devons : comprendre un article de journal allemand, maîtriser le vocabulaire de la femme, de la famille et du travail, et savoir repérer les connecteurs logiques qui organisent les arguments.
\end{aretenir}

\courssub{Lecture du texte support}

\textbf{Consigne de classe.} Lecture silencieuse d'abord (3 minutes, crayon à la main), puis lecture à voix haute par l'enseignant. Pendant la lecture silencieuse, souligne les mots que tu ne comprends pas et entoure les mots que tu reconnais.

\begin{textebox}[Die Frau in Stadt und Land (article de journal)]
In deutschen Städten arbeiten viele Frauen heute in einem Beruf, weil sie eine gute Schulbildung haben. Sie sind Ärztinnen, Lehrerinnen, Ingenieurinnen oder Verkäuferinnen. Obwohl sie berufstätig sind, kümmern sie sich oft noch um den Haushalt und die Kinder. Viele Männer helfen inzwischen mit, aber die meiste Arbeit im Haus machen immer noch die Frauen. Deshalb sprechen Soziologen von einer „doppelten Belastung" : die Frau hat einen Beruf und gleichzeitig eine zweite Arbeit zu Hause.

Auf dem Land ist die Situation anders. Viele Frauen arbeiten dort auf dem Feld und haben wenig Zeit für eine Ausbildung. Manche Mädchen verlassen die Schule früh, weil die Familie ihre Hilfe braucht. Trotzdem gibt es auch positive Beispiele : in einigen Dörfern gründen Frauen kleine Geschäfte und verkaufen ihre Produkte auf dem Markt. So verdienen sie eigenes Geld und werden unabhängiger.

Deshalb fordern Politikerinnen mehr Gleichberechtigung zwischen Mann und Frau, in der Stadt und im Dorf. Sie wollen bessere Schulen für Mädchen, mehr Kindergärten und faire Löhne für alle. Denn nur wenn Frauen und Männer die gleichen Chancen haben, kann sich die Gesellschaft wirklich entwickeln. Die Gleichberechtigung ist also nicht nur ein Problem der Frauen, sondern eine Aufgabe für die ganze Gesellschaft.
\end{textebox}

\begin{textebox}[Glossaire du texte]
die Gleichberechtigung (-, f.) = l'égalité des droits ; die Schulbildung (-en, f.) = l'instruction scolaire ; der Haushalt (-e, m.) = le ménage ; der Beruf (-e, m.) = le métier ; das Feld (-er, n.) = le champ ; die Ausbildung (-en, f.) = la formation ; die Chance (-n, f.) = la chance, l'opportunité ; sich kümmern um (+ Akk.) = s'occuper de ; berufstätig = qui exerce un métier, actif ; die Belastung (-en, f.) = la charge ; der Lohn (Löhne, m.) = le salaire ; sich entwickeln = se développer ; die Aufgabe (-n, f.) = la tâche ; unabhängig = indépendant ; der Markt (Märkte, m.) = le marché.
\end{textebox}

\courssub{Repérage : thème, type de texte, source, public}

Avant de répondre aux questions, il faut identifier le document. C'est la première étape de tout commentaire de texte au Bac.

\begin{center}
\begin{tabularx}{\linewidth}{|X|X|}\hline
\rowcolor{popblueL} \textbf{Critère} & \textbf{Réponse} \\ \hline
Thème (\cle{das Thema}) & La situation de la femme en ville et au village (\all{die Situation der Frau in Stadt und Land}) \\ \hline
Type de texte (\cle{die Textsorte}) & Article de journal (\all{ein Zeitungsartikel}) : ton informatif, arguments, connecteurs logiques \\ \hline
Source probable & Un journal ou un magazine allemand (rubrique société) \\ \hline
Public visé (\cle{die Leser}) & Le grand public germanophone, des lecteurs intéressés par les questions de société \\ \hline
Intention de l'auteur & Informer et convaincre : il faut plus d'égalité entre hommes et femmes \\ \hline
\end{tabularx}
\end{center}

\begin{attention}[Ne confonds pas !]
Un article de journal n'est ni une lettre ni un récit personnel. Indices à repérer : la 3e personne (\all{viele Frauen}, \all{Politikerinnen}), le présent de vérité générale, les connecteurs d'argumentation (\all{weil}, \all{obwohl}, \all{deshalb}, \all{denn}) et l'absence de « je ». Au Bac, on te demande souvent : \all{Um welche Textsorte handelt es sich?} — De quel type de texte s'agit-il ?
\end{attention}

\courssub{Compréhension globale : les idées principales}

Le texte comporte trois paragraphes. Chacun développe une idée principale :

\begin{itemize}
\item \textbf{Paragraphe 1 — la ville :} les femmes instruites exercent un métier, mais elles assument encore l'essentiel du ménage et des enfants : c'est la « double charge » (\all{die doppelte Belastung}).
\item \textbf{Paragraphe 2 — le village :} les femmes travaillent aux champs, ont peu de formation, mais certaines créent de petits commerces et gagnent en indépendance.
\item \textbf{Paragraphe 3 — l'évolution :} des femmes politiques réclament l'égalité des droits, car le développement de la société en dépend.
\end{itemize}

\begin{popfigure}
\begin{center}
\begin{tikzpicture}[
  node distance=0.55cm and 1.1cm,
  box/.style={rectangle, rounded corners, draw=popblue, fill=popblueL, align=center, minimum width=3.6cm, minimum height=1cm, font=\small},
  arr/.style={-{Stealth[length=3mm]}, thick, popdark}]
\node[box, align=center] (a){1\textsuperscript{re} lecture\\ globale : thème,\\ type de texte};
\node[box, below=of a, align=center] (b){2\textsuperscript{e} lecture\\ détaillée : souligner\\ mots-clés et connecteurs};
\node[box, below=of b, align=center] (c){Questions :\\ répondre par des\\ phrases complètes};
\node[box, below=of c, align=center] (d){Résumé :\\ 3 idées principales\\ avec ses propres mots};
\draw[arr] (a) -- (b);
\draw[arr] (b) -- (c);
\draw[arr] (c) -- (d);
\end{tikzpicture}
\end{center}
\legende{Organigramme de la méthode de lecture en quatre étapes}
\end{popfigure}

\begin{methode}[Comprendre un texte au Bac : la démarche gagnante]
\begin{enumerate}
\item Lis d'abord les questions avant le texte : tu sauras ce que tu cherches.
\item Lis le texte une première fois pour le sens général, sans dictionnaire.
\item À la deuxième lecture, souligne les mots-clés et entoure les connecteurs (\all{weil}, \all{obwohl}, \all{deshalb}, \all{denn}) : ils signalent les arguments (cause, opposition, conséquence).
\item Réponds par une phrase complète, en reprenant les termes de la question, jamais par un simple « oui » ou « non ».
\item Pour les questions \all{richtig/falsch}, justifie toujours ta réponse par une citation courte du texte.
\end{enumerate}
\end{methode}

\courssub{Compréhension détaillée : questions de type Bac}

\begin{exoresolu}[Questions de compréhension sur le texte]
\textbf{Énoncé.} 1) Entscheiden Sie : richtig (r) oder falsch (f) ? a) In der Stadt haben viele Frauen keinen Beruf. b) Auf dem Land arbeiten viele Frauen auf dem Feld. 2) Warum sprechen Soziologen von einer „doppelten Belastung" ? 3) Was fordern die Politikerinnen ?
\tcblower
\textbf{Solution détaillée.} 1a) \textbf{falsch} — le texte dit : \all{„In deutschen Städten arbeiten viele Frauen heute in einem Beruf."} (beaucoup de femmes exercent un métier). 1b) \textbf{richtig} — \all{„Auf dem Land ... arbeiten viele Frauen auf dem Feld."} 2) \all{Soziologen sprechen von einer doppelten Belastung, weil die Frau einen Beruf hat und gleichzeitig eine zweite Arbeit zu Hause macht.} (Les sociologues parlent d'une double charge parce que la femme a un métier et en même temps un deuxième travail à la maison.) 3) \all{Die Politikerinnen fordern mehr Gleichberechtigung zwischen Mann und Frau, bessere Schulen für Mädchen, mehr Kindergärten und faire Löhne.} Remarque la méthode : chaque réponse reprend les termes de la question et cite le texte.
\end{exoresolu}

\begin{exercice}[À toi de jouer]
1) Richtig oder falsch ? Begründen Sie mit dem Text. a) Viele Mädchen auf dem Land verlassen die Schule früh. b) Die Gleichberechtigung ist nur ein Problem der Frauen. 2) Was machen manche Frauen in den Dörfern, um Geld zu verdienen ? 3) Warum ist die Gleichberechtigung wichtig für die Gesellschaft ?
\end{exercice}

\begin{corrige}[Exercice « À toi de jouer »]
1a) \textbf{richtig} — \all{„Manche Mädchen verlassen die Schule früh, weil die Familie ihre Hilfe braucht."} 1b) \textbf{falsch} — \all{„Die Gleichberechtigung ist ... eine Aufgabe für die ganze Gesellschaft."} 2) \all{Manche Frauen gründen kleine Geschäfte und verkaufen ihre Produkte auf dem Markt.} (Elles fondent de petits commerces et vendent leurs produits au marché.) 3) \all{Nur wenn Frauen und Männer die gleichen Chancen haben, kann sich die Gesellschaft wirklich entwickeln.} (Seulement si femmes et hommes ont les mêmes chances, la société peut vraiment se développer.)
\end{corrige}

\courssub{Le champ lexical de la femme et du travail dans le texte}

Cherchons maintenant dans le texte tous les mots qui appartiennent au champ lexical de la femme et du travail. C'est une compétence exigée au Bac : \all{Finden Sie Wörter aus dem Wortfeld ...}

\begin{center}
\begin{tabularx}{\linewidth}{|X|X|X|}\hline
\rowcolor{popblueL} \textbf{Deutsch} & \textbf{Français} & \textbf{Famille de mots} \\ \hline
die Frau (-en) & la femme & die Politikerin, die Ärztin, die Lehrerin \\ \hline
der Beruf (-e) & le métier & berufstätig (actif), die Berufstätigkeit \\ \hline
die Arbeit (-en) & le travail & arbeiten, der Arbeiter, die Arbeiterin \\ \hline
die Schulbildung (-en) & l'instruction & die Schule, bilden, die Ausbildung \\ \hline
der Haushalt (-e) & le ménage & das Haus, der Haushaltsartikel \\ \hline
das Feld (-er) & le champ & die Feldarbeit (travail des champs) \\ \hline
die Gleichberechtigung (-) & l'égalité des droits & gleich (égal), das Recht (le droit) \\ \hline
der Lohn (Löhne) & le salaire & verdienen (gagner), das Geld \\ \hline
\end{tabularx}
\end{center}

\begin{exemplebox}[Deux phrases clés du texte, traduites]
\all{Viele Frauen arbeiten in einem Beruf, weil sie eine gute Schulbildung haben.} = Beaucoup de femmes exercent un métier parce qu'elles ont une bonne instruction.

\all{Deshalb fordern Politikerinnen mehr Gleichberechtigung.} = C'est pourquoi des femmes politiques réclament plus d'égalité.
\end{exemplebox}

\begin{attention}[Faux amis et pièges de genre]
\begin{itemize}
\item \all{die Ausbildung} = la formation professionnelle, pas « l'éducation » en général (qui se dit plutôt \all{die Erziehung}).
\item \all{der Beruf} est masculin : \all{ein guter Beruf}, jamais « eine Beruf ».
\item Les métiers féminins se forment souvent avec \all{-in} : \all{der Lehrer} $\rightarrow$ \all{die Lehrerin} ; \all{der Arzt} $\rightarrow$ \all{die Ärztin} (avec Umlaut !).
\item \all{das Feld} a le pluriel \all{die Felder} : \all{auf dem Feld arbeiten} = travailler aux champs.
\end{itemize}
\end{attention}

\courssub{Synthèse : tableau Stadt / Land à compléter}

Complète ce tableau avec les informations du texte (activité en binôme, puis correction collective) :

\begin{popfigure}
\begin{center}
\begin{tikzpicture}[
  cell/.style={rectangle, draw=popdark, align=left, text width=5.6cm, minimum height=2.6cm, inner sep=6pt, font=\small},
  head/.style={rectangle, draw=popdark, fill=poporangeL, align=center, text width=5.6cm, minimum height=0.8cm, font=\small\bfseries}]
\node[head] (h1) {Die Stadt (la ville)};
\node[head, right=0cm of h1] (h2) {Das Land (le village)};
\node[cell, below=0cm of h1, align=center] (c1){Frauen haben eine gute Schulbildung und einen Beruf (Ärztin, Lehrerin ...)\\[2pt] Aber : doppelte Belastung = Beruf + Haushalt + Kinder};
\node[cell, below=0cm of h2, align=center] (c2){Frauen arbeiten auf dem Feld, wenig Zeit für eine Ausbildung\\[2pt] Aber : manche gründen kleine Geschäfte und werden unabhängiger};
\end{tikzpicture}
\end{center}
\legende{Tableau comparatif Stadt / Land complété d'après le texte}
\end{popfigure}

\begin{savaistu}[L'égalité dans la Loi fondamentale allemande]
En Allemagne, l'égalité entre hommes et femmes est inscrite dans la Loi fondamentale (\all{das Grundgesetz}, article 3) depuis 1949 : \all{„Männer und Frauen sind gleichberechtigt."} (Les hommes et les femmes ont les mêmes droits.) Pourtant, le débat sur le partage des tâches domestiques reste d'actualité, exactement comme au Cameroun, où les femmes assurent souvent à la fois le commerce, les champs et le foyer. Le mot \all{Rabenmutter} montre que les mentalités évoluent lentement, même quand les lois changent.
\end{savaistu}

\begin{experience}[Production orale guidée : mini-débat]
Par groupes de quatre. Deux élèves défendent la thèse : \all{„Die Frau soll nur zu Hause bleiben."} ; deux autres défendent : \all{„Die Frau soll einen Beruf haben."} Utilisez obligatoirement trois connecteurs du texte (\all{weil}, \all{obwohl}, \all{deshalb}) et cinq mots du glossaire. Exemple de départ : \all{Ich meine, dass Frauen arbeiten sollen, weil ...} (Je pense que les femmes devraient travailler parce que ...)
\end{experience}

\coursec{Lexique thématique : la femme, la famille, le travail}

Après la lecture globale du texte support, nous allons maintenant construire méthodiquement le vocabulaire indispensable pour parler de la situation de la femme, en ville comme au village. Un bon lexique ne s'apprend pas mot par mot : il s'organise en champs thématiques, en familles de mots et en réseaux. C'est exactement la démarche que nous allons suivre ici, en cinq champs : l'égalité, l'éducation, le foyer, le travail et la vie rurale.

\courssub{Observer d'abord : un mini-texte pour découvrir le lexique}

\begin{textebox}[„Aminas Alltag in Yaoundé“]
Amina ist 35 Jahre alt und lebt mit ihrer Familie in Yaoundé. Sie ist Lehrerin und arbeitet jeden Tag in einer Schule. Ihr Mann arbeitet auch, aber Amina macht noch den größten Teil der Hausarbeit: Sie kocht, putzt und kümmert sich um die drei Kinder. Am Abend erzieht sie ihre Kinder und hilft ihnen bei den Hausaufgaben. Ihre Mutter lebt auf dem Land, in einem kleinen Dorf. Dort arbeiten viele Frauen auf dem Feld und bringen die Ernte auf den Markt. Amina sagt: „Die Gleichberechtigung ist wichtig. Mädchen müssen eine gute Schulbildung bekommen, dann finden sie später einen guten Beruf.“
\end{textebox}

\textbf{Glossaire :} \all{der Alltag} (le quotidien) ; \all{die Hausarbeit} (les tâches ménagères) ; \all{sich kümmern um} (s'occuper de) ; \all{die Ernte} (la récolte) ; \all{der Markt, die Märkte} (le marché) ; \all{die Gleichberechtigung} (l'égalité des droits) ; \all{die Schulbildung} (l'instruction scolaire).

Ce court texte contient déjà presque tout le lexique de la leçon. Observons-le, classons-le, puis apprenons-le.

\courssub{Champ 1 — La femme et l'égalité}

\begin{definition}[die Gleichberechtigung]
\all{Die Gleichberechtigung} (féminin, sans pluriel usuel) désigne l'égalité des droits entre les hommes et les femmes. C'est un mot composé : \all{gleich} (égal) + \all{die Berechtigung} (le droit, la légitimation). On dit : \all{die Gleichberechtigung von Mann und Frau} (l'égalité entre l'homme et la femme).
\end{definition}

\begin{center}
\begin{tabularx}{\linewidth}{|X|X|X|}\hline
\rowcolor{popblueL} Allemand & Français & Remarque \\ \hline
die Frau, -en & la femme & aussi « Madame » : \all{Frau Mbarga} \\ \hline
die Gleichberechtigung & l'égalité des droits & mot composé, sans pluriel \\ \hline
die Rechte (Pl.) & les droits & le singulier : \all{das Recht, -e} \\ \hline
die Rolle, -n & le rôle & \all{die Rolle der Frau} : le rôle de la femme \\ \hline
die Emanzipation & l'émancipation & faux ami partiel : sens politique et social \\ \hline
der Mann, -¨er & l'homme, le mari & attention au pluriel irrégulier \all{Männer} \\ \hline
\end{tabularx}
\end{center}

\begin{exemplebox}[Exemples traduits]
\all{Die Gleichberechtigung von Mann und Frau ist ein wichtiges Thema.} « L'égalité des droits entre l'homme et la femme est un sujet important. »

\all{Die Rolle der Frau hat sich verändert.} « Le rôle de la femme a changé. »

\all{Frauen kämpfen für ihre Rechte.} « Les femmes luttent pour leurs droits. »
\end{exemplebox}

\courssub{Champ 2 — L'éducation}

\begin{center}
\begin{tabularx}{\linewidth}{|X|X|X|}\hline
\rowcolor{popblueL} Allemand & Français & Remarque \\ \hline
die Schulbildung & l'instruction scolaire & \all{Schule} + \all{Bildung} \\ \hline
die Ausbildung, -en & la formation professionnelle & \all{eine Ausbildung machen} : suivre une formation \\ \hline
die Schule, -n & l'école & \all{in die Schule gehen} : aller à l'école \\ \hline
lernen & apprendre & verbe faible : \all{ich lerne, sie lernt} \\ \hline
studieren & étudier (à l'université) & \all{an der Universität studieren} \\ \hline
das Mädchen, - & la fille, la jeune fille & toujours neutre ! \\ \hline
\end{tabularx}
\end{center}

\begin{attention}[lernen ou studieren ?]
En français, « étudier » couvre tout. En allemand, on distingue : \all{lernen} = apprendre (à l'école, une langue, une leçon) ; \all{studieren} = faire des études universitaires. On dit donc \all{Das Mädchen lernt in der Schule} (la fille apprend à l'école) mais \all{Sie studiert Medizin} (elle fait des études de médecine). Ne dites jamais \all{sie studiert in der Schule}.
\end{attention}

\courssub{Champ 3 — Le foyer}

\begin{center}
\begin{tabularx}{\linewidth}{|X|X|X|}\hline
\rowcolor{popblueL} Allemand & Français & Remarque \\ \hline
der Haushalt, -e & le ménage, le foyer & \all{den Haushalt machen} : faire le ménage \\ \hline
die Hausarbeit & les tâches ménagères & \all{Haus} + \all{Arbeit} \\ \hline
kochen & cuisiner, faire la cuisine & verbe faible \\ \hline
putzen & nettoyer, faire le ménage & verbe faible \\ \hline
sich kümmern um (+ Akk.) & s'occuper de & verbe réfléchi : \all{ich kümmere mich um...} \\ \hline
die Kinder erziehen & élever les enfants & verbe fort : \all{erzieht, erzog, hat erzogen} \\ \hline
das Kind, -er & l'enfant & pluriel irrégulier \all{Kinder} \\ \hline
\end{tabularx}
\end{center}

\begin{exemplebox}[Exemples traduits]
\all{Sie erzieht ihre Kinder und macht den Haushalt.} « Elle élève ses enfants et fait le ménage. »

\all{Am Morgen kocht sie, dann putzt sie das Haus.} « Le matin, elle cuisine, puis elle nettoie la maison. »

\all{Sie kümmert sich um ihre kranke Mutter.} « Elle s'occupe de sa mère malade. »
\end{exemplebox}

\begin{attention}[Verbe réfléchi obligatoire]
\all{Sich kümmern um} est toujours réfléchi : le pronom réfléchi (\all{mich, dich, sich, uns, euch, sich}) ne se supprime jamais, contrairement au français où « s'occuper » peut se simplifier mentalement. Erreur typique : \all{Sie kümmert um die Kinder} — faux ! Il faut \all{Sie kümmert sich um die Kinder}.
\end{attention}

\courssub{Champ 4 — Le travail}

\begin{center}
\begin{tabularx}{\linewidth}{|X|X|X|}\hline
\rowcolor{popblueL} Allemand & Français & Remarque \\ \hline
der Beruf, -e & le métier, la profession & \all{Was sind Sie von Beruf?} : quel est votre métier ? \\ \hline
die Arbeit, -en & le travail (en général) & \all{zur Arbeit gehen} : aller au travail \\ \hline
arbeiten & travailler & verbe faible : \all{er arbeitet} (e intercalaire) \\ \hline
berufstätig sein & exercer une profession & \all{Sie ist berufstätig.} : elle travaille (professionnellement) \\ \hline
verdienen & gagner (de l'argent) & \all{Geld verdienen} : gagner de l'argent \\ \hline
der Lohn, -¨e & le salaire, la paie & pluriel \all{Löhne} \\ \hline
\end{tabularx}
\end{center}

\begin{attention}[die Arbeit ou der Beruf ?]
\all{Die Arbeit} désigne le travail en général, l'activité, la tâche : \all{Die Arbeit auf dem Feld ist schwer} (le travail aux champs est dur). \all{Der Beruf} désigne le métier exercé, la profession qualifiée : \all{Ihr Beruf ist Lehrerin} (son métier est institutrice). Un francophone traduit souvent « travail » automatiquement par \all{Arbeit} : vérifiez toujours si le sens est « activité » (\all{Arbeit}) ou « métier » (\all{Beruf}).
\end{attention}

\begin{propriete}[Le féminin des noms de métiers : le suffixe -in]
Pour former le féminin d'un nom de métier ou de fonction, l'allemand ajoute le suffixe \cle{-in} au nom masculin, et l'article \all{der} devient \all{die} :

\all{der Lehrer} → \all{die Lehrerin} (l'instituteur / l'institutrice) ; \all{der Arbeiter} → \all{die Arbeiterin} (l'ouvrier / l'ouvrière) ; \all{der Politiker} → \all{die Politikerin} (l'homme politique / la femme politique) ; \all{der Bauer} → \all{die Bäuerin} (le paysan / la paysanne, avec Umlaut).

Au pluriel, le féminin prend \cle{-innen} : \all{die Lehrerinnen}, \all{die Arbeiterinnen}, \all{die Politikerinnen}. Le masculin pluriel, lui, reste souvent identique au singulier : \all{die Lehrer}, \all{die Arbeiter}.

Attention : les noms de métiers féminins en \all{-in} portent toujours l'article \all{die}, même si le mot se termine par une consonne.
\end{propriete}

\begin{exemplebox}[Exemples traduits]
\all{Die Frauen arbeiten auf dem Feld.} « Les femmes travaillent aux champs. »

\all{Meine Schwester ist Lehrerin und verdient gut.} « Ma sœur est institutrice et gagne bien sa vie. »

\all{Viele Frauen sind heute berufstätig.} « Beaucoup de femmes exercent aujourd'hui une profession. »

\all{Der Lohn für diese Arbeit ist klein.} « Le salaire pour ce travail est faible. »
\end{exemplebox}

\courssub{Champ 5 — La vie rurale}

\begin{center}
\begin{tabularx}{\linewidth}{|X|X|X|}\hline
\rowcolor{popblueL} Allemand & Français & Remarque \\ \hline
das Dorf, -¨er & le village & pluriel \all{Dörfer} \\ \hline
das Feld, -er & le champ & \all{auf dem Feld arbeiten} : travailler aux champs \\ \hline
die Ernte, -n & la récolte & \all{die Ernte einbringen} : rentrer la récolte \\ \hline
das Land, -¨er & le pays ; la campagne & \all{auf dem Land} : à la campagne \\ \hline
die Stadt, -¨e & la ville & pluriel \all{Städte} \\ \hline
der Markt, -¨e & le marché & \all{auf den Markt gehen} : aller au marché \\ \hline
\end{tabularx}
\end{center}

\begin{propriete}[Stadt ou Dorf ? Les expressions à connaître]
\all{Die Stadt} (pluriel \all{Städte}) s'oppose à \all{das Dorf} (pluriel \all{Dörfer}). Les expressions figées sont :

\all{in der Stadt leben} = vivre en ville ; \all{auf dem Land leben} = vivre à la campagne ; \all{in einem Dorf leben} = vivre dans un village ; \all{in die Stadt ziehen} = s'installer en ville.

Notez la différence de préposition : \all{in der Stadt} mais \all{auf dem Land}. C'est un usage figé, il faut l'apprendre tel quel.
\end{propriete}

\begin{attention}[« auf dem Land » ne veut pas dire « sur la terre » !]
\all{Auf dem Land} signifie « à la campagne » : \all{Auf dem Land leben viele Familien.} « À la campagne vivent beaucoup de familles. » Ne confondez pas :
\begin{itemize}
\item \all{das Land} = le pays (\all{Kamerun ist ein Land in Afrika}) ou la campagne (\all{auf dem Land}) ;
\item \all{die Erde} = la terre (la planète ou le sol) : \all{Die Erde ist rund} (la Terre est ronde).
\end{itemize}
Autre piège : \all{die Stadt} est féminin (\all{die}), alors que le français « ville » peut faire hésiter ; et \all{das Dorf} est neutre, jamais \all{der Dorf}.
\end{attention}

\courssub{Les familles de mots : apprendre en réseau}

L'allemand construit des familles de mots très régulières. Les connaître multiplie votre vocabulaire sans effort.

\begin{propriete}[Famille de arbeiten et famille de bilden]
Autour du verbe \all{arbeiten} (travailler) : \all{die Arbeit} (le travail), \all{der Arbeiter} (l'ouvrier), \all{die Arbeiterin} (l'ouvrière), \all{berufstätig} (qui exerce une profession), \all{der Beruf} (le métier).

Autour du verbe \all{bilden} (former, éduquer) : \all{die Bildung} (l'éducation, la formation), \all{die Schulbildung} (l'instruction scolaire), \all{die Ausbildung} (la formation professionnelle).

Méthode : repérez la racine commune (\all{arbeit-}, \all{bild-}) et les suffixes (\all{-ung} forme des noms féminins d'action, \all{-er} forme des noms d'agents masculins).
\end{propriete}

\begin{popfigure}
\begin{tikzpicture}[mindmap, grow cyclic, every node/.style={concept, align=center, font=\small},
root concept/.append style={concept color=popblue, font=\bfseries},
level 1/.append style={level distance=3.2cm, sibling angle=72},
level 1 concept/.append style={concept color=poporange}]
\node[root concept]{die Frau}
child { node[align=center]{Gleichberechtigung\\die Rechte\\die Rolle} }
child { node[align=center]{Schulbildung\\die Schule\\lernen, studieren} }
child { node[align=center]{Haushalt\\kochen, putzen\\Kinder erziehen} }
child { node[align=center]{Beruf\\arbeiten\\verdienen} }
child { node[align=center]{Arbeit auf dem Feld\\das Dorf\\die Ernte} };
\end{tikzpicture}
\legende{Carte mentale du lexique : cinq champs autour de \all{die Frau}}
\end{popfigure}

\begin{popfigure}
\begin{tikzpicture}[
grow=down,
level distance=1.4cm,
level 1/.style={sibling distance=3.4cm},
level 2/.style={sibling distance=2.6cm},
every node/.style={draw=popblue, rounded corners=2pt, fill=popblueL, align=center, font=\small, inner sep=3pt},
edge from parent/.style={draw=popdark, -{Stealth}}]
\node[fill=poporangeL, draw=poporange, font=\bfseries\small]{arbeiten (travailler)}
child { node[align=center]{die Arbeit\\le travail} }
child { node[align=center]{der Arbeiter\\l'ouvrier}
child { node[align=center]{die Arbeiterin\\l'ouvrière} } }
child { node[align=center]{berufstätig\\qui travaille} }
child { node[align=center]{der Beruf\\le métier} };
\end{tikzpicture}
\legende{Arbre de la famille de mots : \all{arbeiten}}
\end{popfigure}

\begin{savaistu}[Les mots composés, la clé du Bac]
L'allemand compose facilement les mots : \all{Haus} + \all{Arbeit} = \all{die Hausarbeit} ; \all{Schule} + \all{Bildung} = \all{die Schulbildung} ; \all{Feld} + \all{Arbeit} = \all{die Feldarbeit}. Au Bac, vous rencontrerez des mots composés inconnus : ne paniquez pas, découpez-les ! Le dernier élément donne le sens général et le genre grammatical (\all{die Arbeit} → \all{die Hausarbeit}, féminin), les éléments précédents précisent le sens. C'est une stratégie de compréhension redoutablement efficace.
\end{savaistu}

\courssub{Texte d'entraînement au lexique}

\begin{textebox}[„Frauen in Stadt und Land“]
In Kamerun ist die Situation der Frau in der Stadt anders als auf dem Land. In der Stadt, zum Beispiel in Douala oder Yaoundé, haben viele Frauen einen Beruf: Sie sind Lehrerinnen, Ärztinnen oder Verkäuferinnen. Sie sind berufstätig und verdienen einen Lohn. Trotzdem machen sie oft noch die ganze Hausarbeit: Sie kochen, putzen und kümmern sich um die Kinder. Die Gleichberechtigung ist hier ein großes Thema, denn viele Männer helfen im Haushalt noch wenig.

Auf dem Land ist das Leben der Frauen oft schwerer. In den Dörfern arbeiten die Frauen auf dem Feld. Sie pflanzen Mais und Maniok, und nach der Ernte bringen sie die Produkte auf den Markt. Diese Arbeit ist hart, aber der Lohn ist klein. Viele Mädchen im Dorf können nicht lange in die Schule gehen, weil die Familien arm sind. Ohne Schulbildung finden sie später keinen guten Beruf.

Heute kämpfen viele Organisationen für die Rechte der Frauen. Sie sagen: Jedes Mädchen hat ein Recht auf Bildung. Wenn Mädchen lernen und später studieren können, verändert sich die Rolle der Frau in der Familie und in der Gesellschaft. Dann können Frauen auch in der Politik arbeiten und für mehr Gleichberechtigung kämpfen — in der Stadt und auf dem Land.
\end{textebox}

\textbf{Glossaire :} \all{die Ärztin, -nen} (la femme médecin) ; \all{die Verkäuferin, -nen} (la vendeuse) ; \all{trotzdem} (pourtant, néanmoins) ; \all{pflanzen} (planter) ; \all{der Mais} (le maïs) ; \all{der Maniok} (le manioc) ; \all{hart} (dur) ; \all{arm} (pauvre) ; \all{die Gesellschaft, -en} (la société) ; \all{die Politik} (la politique).

\textbf{Questions de compréhension :}
\begin{enumerate}
\item \all{Welche Berufe haben Frauen in der Stadt?} (Quels métiers les femmes exercent-elles en ville ?)
\item \all{Was machen die Frauen auf dem Land nach der Ernte?} (Que font les femmes à la campagne après la récolte ?)
\item \all{Warum können viele Mädchen im Dorf nicht lange zur Schule gehen?} (Pourquoi beaucoup de filles au village ne peuvent-elles pas aller longtemps à l'école ?)
\item \all{Was verändert sich, wenn Mädchen studieren können?} (Qu'est-ce qui change quand les filles peuvent faire des études ?)
\end{enumerate}

\courssub{Exercice résolu et entraînement}

\begin{exoresolu}[Lexique : complétez avec le mot qui convient]
Complétez avec : \all{die Gleichberechtigung, der Haushalt, das Feld, die Schulbildung, der Lohn, berufstätig}.

1. Viele Frauen in Yaoundé sind heute \dots\ und arbeiten als Lehrerinnen.
2. Die Mädchen brauchen eine gute \dots , um später einen Beruf zu finden.
3. Die Frauen arbeiten auf dem \dots\ und bringen die Ernte auf den Markt.
4. Meine Mutter macht den ganzen \dots : Sie kocht und putzt.
5. Für diese harte Arbeit bekommen die Frauen nur einen kleinen \dots .
6. \dots\ bedeutet: Männer und Frauen haben die gleichen Rechte.
\tcblower
\textbf{Solution détaillée :}

1. \all{berufstätig} — l'indice est « arbeiten als Lehrerinnen » : elles exercent une profession.
2. \all{Schulbildung} — l'indice est « Mädchen » et « um später einen Beruf zu finden » : il s'agit de l'instruction scolaire (et non \all{Ausbildung}, qui est la formation professionnelle).
3. \all{Feld} — l'expression figée est \all{auf dem Feld arbeiten} (travailler aux champs) ; datif après \all{auf} (position).
4. \all{Haushalt} — les indices « kocht und putzt » renvoient aux tâches ménagères ; expression : \all{den Haushalt machen}.
5. \all{Lohn} — « bekommen » + « klein » évoquent le salaire ; \all{der Lohn} est masculin : \all{einen kleinen Lohn} (accusatif).
6. \all{Gleichberechtigung} — la définition « Männer und Frauen haben die gleichen Rechte » correspond exactement à l'égalité des droits.
\end{exoresolu}

\begin{exercice}[À vous de jouer]
1. Donnez le féminin des métiers suivants : \all{der Lehrer, der Arbeiter, der Politiker, der Verkäufer, der Bauer}.
2. Traduisez en allemand : a) Les femmes travaillent aux champs. b) À la campagne vivent beaucoup de familles. c) Elle s'occupe de ses enfants et fait le ménage. d) Ma sœur est institutrice et gagne un bon salaire.
3. Découpez les mots composés et donnez leur sens : \all{die Hausarbeit, die Schulbildung, die Feldarbeit, die Kindererziehung}.
\end{exercice}

\begin{corrige}[Exercice « À vous de jouer »]
1. \all{die Lehrerin, die Arbeiterin, die Politikerin, die Verkäuferin, die Bäuerin} (attention à l'Umlaut de \all{Bäuerin}).
2. a) \all{Die Frauen arbeiten auf dem Feld.} b) \all{Auf dem Land leben viele Familien.} (verbe en 2e position !) c) \all{Sie kümmert sich um ihre Kinder und macht den Haushalt.} d) \all{Meine Schwester ist Lehrerin und verdient einen guten Lohn.}
3. \all{Haus + Arbeit} : les tâches ménagères ; \all{Schule + Bildung} : l'instruction scolaire ; \all{Feld + Arbeit} : le travail des champs ; \all{Kinder + Erziehung} : l'éducation des enfants.
\end{corrige}

\begin{aretenir}[L'essentiel du lexique]
\begin{itemize}
\item Cinq champs : égalité (\all{die Gleichberechtigung, die Rechte}), éducation (\all{die Schulbildung, lernen, studieren}), foyer (\all{der Haushalt, sich kümmern um, erziehen}), travail (\all{der Beruf, die Arbeit, verdienen, der Lohn}), vie rurale (\all{das Dorf, das Feld, die Ernte}).
\item Féminin des métiers : suffixe \all{-in}, pluriel \all{-innen} (\all{die Lehrerin → die Lehrerinnen}).
\item Expressions figées : \all{in der Stadt leben}, \all{auf dem Land leben}, \all{auf dem Feld arbeiten}, \all{den Haushalt machen}.
\item Pièges : \all{auf dem Land} = à la campagne ; \all{die Arbeit} (activité) ≠ \all{der Beruf} (métier) ; \all{lernen} ≠ \all{studieren}.
\item Stratégie : découper les mots composés pour deviner leur sens.
\end{itemize}
\end{aretenir}

Ce lexique est maintenant votre outil de travail. Dans la prochaine section, nous verrons comment ordonner ces mots dans la phrase allemande, car un vocabulaire riche ne suffit pas : il faut encore respecter la place du verbe.

\coursec{Grammaire 1 : l'ordre des mots}

Dans la section précédente, nous avons construit le vocabulaire du thème : \all{die Frau}, \all{die Gleichberechtigung}, \all{die Schulbildung}, \all{der Haushalt}, \all{der Beruf}, \all{die Arbeit auf dem Feld}. Nous savons maintenant \emph{de quoi} parler. Il reste à apprendre \emph{comment} assembler ces mots en phrases allemandes correctes. Or c'est précisément là que le francophone trébuche le plus souvent : en allemand, ce n'est pas le sujet qui commande la phrase, c'est \cle{le verbe}. Cette section est donc la charpente de toute votre expression écrite et orale en Terminale.

\courssub{Observation : deux phrases, un même sens, un même verbe}

Observons d'abord deux phrases tirées de notre thème.

\begin{exemplebox}[Observer avant de comprendre]
\all{Viele Frauen in Douala arbeiten heute in einem Beruf.}\\
« Beaucoup de femmes à Douala exercent aujourd'hui un métier. »

\all{Heute arbeiten viele Frauen in Douala in einem Beruf.}\\
« Aujourd'hui, beaucoup de femmes à Douala exercent un métier. »
\end{exemplebox}

Les deux phrases disent exactement la même chose. Pourtant, l'ordre des mots a changé : dans la première, le sujet \all{viele Frauen} est en tête ; dans la seconde, c'est le complément de temps \all{heute} qui ouvre la phrase, et le sujet a reculé \emph{après} le verbe. Une seule chose n'a pas bougé : le verbe conjugué \all{arbeiten} occupe dans les deux cas la \cle{deuxième position}.

Retenez bien ce que l'allemand appelle une « position » : ce n'est pas un mot, c'est un \emph{groupe de sens}. \all{Viele Frauen in Douala} (quatre mots) ne compte que pour \emph{une} position, la première. Le verbe vient ensuite, quoi qu'il arrive.

\courssub{La règle fondamentale : le verbe en 2e position}

\begin{propriete}[Le verbe conjugué en 2e position (phrase principale)]
Dans une \textbf{phrase principale affirmative} allemande, le verbe conjugué occupe \textbf{toujours la deuxième position}, quel que soit l'élément placé en tête (sujet, complément de temps, de lieu, de manière).
\begin{itemize}
\item Si le sujet est en tête : Sujet + \textbf{VERBE} + reste. \all{Die Frau arbeitet auf dem Feld.}
\item Si un complément est en tête : Complément + \textbf{VERBE} + sujet + reste. C'est l'\cle{inversion} sujet-verbe : \all{Auf dem Feld arbeitet die Frau.}
\end{itemize}
La première position ne peut contenir qu'\textbf{un seul} groupe de sens.
\end{propriete}

\begin{popfigure}
\begin{tikzpicture}[
  box/.style={draw=popblue, thick, rounded corners=3pt, fill=popblueL, align=center, minimum height=1.1cm, font=\small},
  vbox/.style={draw=poporange, very thick, rounded corners=3pt, fill=poporangeL, align=center, minimum height=1.1cm, font=\small\bfseries},
  lbl/.style={font=\footnotesize\itshape, text=popmuted}
]
\node[box, align=center] (p1) at (0,0){Position 1\\Sujet \emph{ou} complément};
\node[vbox, align=center] (p2) at (4.6,0){Position 2\\VERBE conjugué};
\node[box, align=center] (p3) at (9.2,0){Reste de la phrase\\(sujet, compléments)};
\node[lbl] at (0,-1.15) {\all{Viele Frauen} / \all{Heute}};
\node[lbl] at (4.6,-1.15) {\all{arbeiten} — ne bouge jamais};
\node[lbl] at (9.2,-1.15) {\all{heute in einem Beruf}};
\draw[-{Stealth}, thick, popdark] (p1) -- (p2);
\draw[-{Stealth}, thick, popdark] (p2) -- (p3);
\end{tikzpicture}
\legende{La phrase principale allemande : le verbe conjugué est le pivot fixe de la 2e position.}
\end{popfigure}

\begin{exemplebox}[L'inversion en action]
\all{In der Stadt haben die Frauen mehr Möglichkeiten.}\\
« En ville, les femmes ont plus de possibilités. »

\all{Auf dem Dorf kümmern sich viele Frauen um den Haushalt und die Arbeit auf dem Feld.}\\
« Au village, beaucoup de femmes s'occupent du ménage et du travail aux champs. »

\all{Seit einigen Jahren kämpfen viele Organisationen für die Gleichberechtigung.}\\
« Depuis quelques années, de nombreuses organisations luttent pour l'égalité des droits. »

\all{In Yaoundé besuchen fast alle Mädchen die Schule.}\\
« À Yaoundé, presque toutes les filles vont à l'école. »
\end{exemplebox}

Remarquez que le français, lui, ne bouge jamais : « En ville, \emph{les femmes ont}... » — le sujet précède toujours le verbe. L'allemand fait le contraire : dès qu'un complément ouvre la phrase, le verbe se précipite en 2e position et \emph{pousse} le sujet derrière lui.

\begin{attention}[L'erreur n° 1 du francophone]
Après un complément placé en tête, ne conservez \textbf{jamais} l'ordre sujet + verbe du français.

On dit : \all{Heute arbeiten viele Frauen in einem Beruf.}\\
On ne dit \textbf{pas} : \emph{Heute viele Frauen arbeiten...} (calque du français « Aujourd'hui beaucoup de femmes travaillent »).

Autre piège : ne placez pas deux compléments avant le verbe. \emph{Heute in der Stadt arbeiten...} est fautif : la première position n'accepte qu'un seul groupe de sens. Dites : \all{Heute arbeiten viele Frauen in der Stadt.} ou \all{In der Stadt arbeiten heute viele Frauen.}
\end{attention}

\courssub{La subordonnée : le verbe conjugué passe en dernière position}

Observons maintenant ce qui se passe quand la phrase est introduite par un connecteur subordonnant comme \all{weil} (parce que).

\begin{exemplebox}[Observer le déplacement du verbe]
Phrase principale : \all{Sie haben eine gute Schulbildung.}\\
« Elles ont une bonne instruction. »

Subordonnée : \all{..., weil sie eine gute Schulbildung haben.}\\
« ..., parce qu'elles ont une bonne instruction. »
\end{exemplebox}

Le verbe \all{haben}, qui était en 2e position, a été \emph{expulsé} à la fin de la phrase, derrière tous les compléments. C'est la deuxième grande loi de la syntaxe allemande.

\begin{propriete}[Le verbe en position finale (phrase subordonnée)]
Dans une \textbf{subordonnée} introduite par un connecteur subordonnant — \all{weil} (parce que), \all{obwohl} (bien que), \all{dass} (que), \all{wenn} (quand, si) — le verbe conjugué se place \textbf{en dernière position} de la subordonnée.

Structure : Connecteur + sujet + compléments + \textbf{VERBE}.

La subordonnée est \textbf{toujours séparée de la principale par une virgule}, sans exception.
\end{propriete}

\begin{popfigure}
\begin{tikzpicture}[
  box/.style={draw=popgreen, thick, rounded corners=3pt, fill=popgreenL, align=center, minimum height=1cm, font=\small},
  vbox/.style={draw=poporange, very thick, rounded corners=3pt, fill=poporangeL, align=center, minimum height=1cm, font=\small\bfseries},
  lbl/.style={font=\footnotesize\itshape, text=popmuted}
]
\node[box, align=center] (c) at (0,0){Connecteur\\\all{weil}};
\node[box, align=center] (s) at (2.6,0){Sujet\\\all{sie}};
\node[box, align=center] (k) at (5.6,0){Compléments\\\all{eine gute Schulbildung}};
\node[vbox, align=center] (v) at (9.6,0){VERBE\\\all{haben}};
\draw[-{Stealth}, thick, popdark] (c) -- (s);
\draw[-{Stealth}, thick, popdark] (s) -- (k);
\draw[-{Stealth}, thick, popdark] (k) -- (v);
\node[lbl] (src) at (4.6,1.5) {Position 2 de la principale : \all{Sie \textbf{haben} eine gute Schulbildung.}};
\draw[-{Stealth}, very thick, poppink, bend right=25] (src.east) to node[midway, above, font=\footnotesize\itshape, text=poppink]{le verbe migre en fin} (v.north);
\end{tikzpicture}
\legende{Dans la subordonnée, le connecteur « aspire » le verbe conjugué vers la dernière position.}
\end{popfigure}

\begin{exemplebox}[Subordonnées avec weil, obwohl, dass, wenn]
\all{Viele Mädchen verlassen die Schule früh, weil ihre Familien arm sind.}\\
« Beaucoup de filles quittent l'école tôt parce que leurs familles sont pauvres. »

\all{Obwohl sie berufstätig sind, kümmern sie sich um den Haushalt.}\\
« Bien qu'elles travaillent, elles s'occupent du ménage. »

\all{Ich denke, dass die Gleichberechtigung sehr wichtig ist.}\\
« Je pense que l'égalité des droits est très importante. »

\all{Wenn eine Frau einen Beruf hat, ist sie unabhängiger.}\\
« Quand une femme a un métier, elle est plus indépendante. »
\end{exemplebox}

Notez le dernier exemple : quand la subordonnée est \emph{en tête} (\all{Obwohl sie berufstätig sind, ...} ; \all{Wenn eine Frau einen Beruf hat, ...}), elle occupe à elle seule la \textbf{première position} de la phrase. Le verbe de la principale doit donc venir \emph{immédiatement après la virgule} : \all{..., kümmern sie sich...} et non \emph{..., sie kümmern sich...}. C'est l'inversion qui s'applique à nouveau — la règle de la 2e position ne connaît pas d'exception.

\begin{aretenir}[Les deux lois de la phrase allemande]
\begin{enumerate}
\item \textbf{Principale} : le verbe conjugué est en \textbf{2e position}, toujours.
\item \textbf{Subordonnée} (après \all{weil}, \all{obwohl}, \all{dass}, \all{wenn}...) : le verbe conjugué est en \textbf{dernière position}, toujours.
\item La virgule sépare obligatoirement les deux phrases.
\item Subordonnée en tête = position 1 : le verbe de la principale suit la virgule directement.
\end{enumerate}
\end{aretenir}

\courssub{Comparaison français / allemand}

\begin{center}
\begin{tabularx}{\linewidth}{|X|X|X|}
\hline
\rowcolor{popblueL} Français & Deutsch & Remarque \\
\hline
Beaucoup de femmes travaillent aujourd'hui. & \all{Viele Frauen arbeiten heute.} & Ordre identique : sujet en tête. \\
\hline
Aujourd'hui, beaucoup de femmes travaillent. & \all{Heute arbeiten viele Frauen.} & L'allemand inverse ; le français non. \\
\hline
Elles restent au village parce qu'elles aiment la vie rurale. & \all{Sie bleiben auf dem Dorf, weil sie das Landleben lieben.} & En allemand, le verbe de la subordonnée va à la fin ; en français, l'ordre ne change pas. \\
\hline
Bien qu'elles travaillent, elles s'occupent du ménage. & \all{Obwohl sie berufstätig sind, kümmern sie sich um den Haushalt.} & Après la virgule, verbe de la principale d'abord. \\
\hline
\end{tabularx}
\end{center}

La différence profonde est la suivante : en français, c'est le \emph{sujet} qui est le repère fixe de la phrase ; en allemand, c'est le \emph{verbe}. Le français demande : « qui fait l'action ? » ; l'allemand demande d'abord : « où est le verbe ? ». Tant que vous chercherez le sujet avant le verbe, vous construirez des phrases françaises avec des mots allemands.

\courssub{Mise en pratique sur un texte}

\begin{textebox}[Die Frauen in meiner Familie]
\all{Meine Großmutter lebt in einem Dorf in der Westregion. Jeden Tag steht sie um fünf Uhr auf, weil sie auf dem Feld arbeiten muss. Sie hat nie eine Schule besucht, denn ihre Eltern waren sehr arm. Heute ist die Situation anders: Meine Schwester studiert in Yaoundé, obwohl das Studium teuer ist. Meine Eltern sagen immer, dass die Schulbildung für Mädchen genauso wichtig ist wie für Jungen. Wenn meine Schwester ihr Studium beendet hat, will sie Ärztin werden. Ich bin stolz auf sie, weil sie für ihre Rechte kämpft. In der Stadt haben die Frauen mehr Chancen, aber auf dem Dorf ist das Leben für viele Frauen noch schwer. Deshalb müssen wir alle für die Gleichberechtigung arbeiten.}

\emph{Glossaire :} \all{das Dorf, die Dörfer} (village) — \all{aufstehen} (se lever) — \all{besuchen} (fréquenter, visiter) — \all{das Studium} (les études universitaires) — \all{die Ärztin, -nen} (la femme médecin) — \all{stolz} (fier) — \all{die Chance, -n} (la chance, l'opportunité) — \all{schwer} (dur, difficile).
\end{textebox}

\textbf{Questions de grammaire sur le texte :}
\begin{enumerate}
\item Relevez trois phrases principales avec un complément en tête et soulignez le verbe en 2e position.
\item Relevez toutes les subordonnées et justifiez la place du verbe conjugué.
\item Pourquoi y a-t-il une virgule devant \all{weil}, \all{obwohl}, \all{dass}, \all{wenn} ?
\end{enumerate}

\courssub{Méthode : vérifier une phrase allemande en trois étapes}

\begin{methode}[Le contrôle en trois points]
Avant de rendre toute copie (thème, résumé, production écrite), vérifiez chaque phrase ainsi :
\begin{enumerate}
\item \textbf{Trouvez le verbe conjugué} de chaque proposition (attention : il peut y en avoir plusieurs si la phrase contient une subordonnée).
\item \textbf{Vérifiez sa position} : 2e position dans la principale ; dernière position dans la subordonnée. Si un complément ouvre la principale, le sujet doit être \emph{après} le verbe.
\item \textbf{Vérifiez la virgule} : il y en a toujours une devant \all{weil}, \all{obwohl}, \all{dass}, \all{wenn}, et entre une subordonnée en tête et la principale qui suit.
\end{enumerate}
\end{methode}

\begin{exoresolu}[Transformer une principale en subordonnée]
Énoncé : transformez la phrase \all{Viele Frauen in der Stadt haben einen Beruf.} en subordonnée introduite par \all{weil}, puis complétez : \all{Sie sind unabhängig, ...}
\tcblower
\textbf{Étape 1} — J'identifie le verbe conjugué : \all{haben} (2e position dans la principale).\\
\textbf{Étape 2} — J'introduis \all{weil} : le verbe doit passer en dernière position. Le sujet \all{viele Frauen} suit le connecteur, les compléments \all{in der Stadt} et \all{einen Beruf} se placent entre le sujet et le verbe.\\
\textbf{Étape 3} — Je n'oublie pas la virgule devant \all{weil}.

Solution : \all{Sie sind unabhängig, weil viele Frauen in der Stadt einen Beruf haben.}\\
« Elles sont indépendantes parce que beaucoup de femmes en ville ont un métier. »

Transformation inverse : \all{Obwohl das Leben auf dem Dorf schwer ist, ...} redonne en principale : \all{Das Leben auf dem Dorf ist schwer.} — le verbe \all{ist} revient en 2e position.
\end{exoresolu}

\begin{exercice}[À vous de jouer]
\begin{enumerate}
\item Remettez dans l'ordre : \all{heute / viele Frauen / arbeiten / in einem Beruf}.
\item Transformez en subordonnée avec \all{weil} : \all{Die Mädchen besuchen die Schule.} (\all{Ich bin froh, ...})
\item Transformez en subordonnée avec \all{obwohl} et placez-la en tête : \all{Sie arbeitet auf dem Feld. Sie verdient wenig Geld.}
\item Traduisez : « Au village, les femmes travaillent beaucoup parce qu'elles doivent nourrir la famille. »
\end{enumerate}
\end{exercice}

\begin{corrige}[Exercice : À vous de jouer]
\begin{enumerate}
\item \all{Heute arbeiten viele Frauen in einem Beruf.} (complément en tête : inversion ; le verbe reste en 2e position).
\item \all{Ich bin froh, weil die Mädchen die Schule besuchen.} (verbe \all{besuchen} à la fin, virgule devant \all{weil}).
\item \all{Obwohl sie auf dem Feld arbeitet, verdient sie wenig Geld.} (subordonnée en tête = position 1 ; le verbe \all{verdient} suit la virgule immédiatement).
\item \all{Auf dem Dorf arbeiten die Frauen viel, weil sie die Familie ernähren müssen.} (inversion après \all{auf dem Dorf} ; \all{müssen} à la fin de la subordonnée, l'infinitif \all{ernähren} juste avant lui).
\end{enumerate}
\end{corrige}

\begin{attention}[Récapitulatif des pièges]
\begin{itemize}
\item \emph{Heute viele Frauen arbeiten...} : faux. Complément en tête $\Rightarrow$ inversion : \all{Heute arbeiten viele Frauen...}
\item \emph{..., weil sie haben eine gute Schulbildung.} : faux. Après \all{weil}, le verbe va à la fin : \all{..., weil sie eine gute Schulbildung haben.}
\item Oublier la virgule devant la subordonnée : faute de ponctuation systématiquement sanctionnée au baccalauréat.
\item \emph{Obwohl sie arbeiten, sie kümmern sich um den Haushalt.} : faux. Après une subordonnée en tête, le verbe de la principale vient d'abord : \all{Obwohl sie arbeiten, kümmern sie sich um den Haushalt.}
\end{itemize}
\end{attention}

Vous maîtrisez désormais les deux lois qui gouvernent toute phrase allemande. La section suivante approfondira les connecteurs \all{weil}, \all{denn}, \all{deshalb} et \all{obwohl} — et vous verrez que tout ce que nous venons d'apprendre sur la place du verbe y trouve immédiatement son application.

\coursec{Grammaire 2 : les connecteurs weil, denn, deshalb, obwohl}

Dans la section précédente, nous avons établi la règle d'or de la phrase allemande : le verbe conjugué occupe toujours la deuxième position dans la phrase principale, et se rejette à la fin dans la subordonnée. Nous allons maintenant mettre cette règle à l'épreuve avec quatre petits mots qui relient les idées : \all{weil}, \all{denn}, \all{deshalb} et \all{obwohl}. Ils servent à exprimer la cause, la conséquence et la concession — exactement ce dont vous avez besoin pour commenter le texte sur la situation de la femme. Mais attention : chacun impose une construction différente. C'est l'un des points les plus lucratifs — et les plus piégeux — du Bac.

\courssub{Observation : quatre phrases, quatre constructions}

\begin{textebox}[Mini-texte d'observation : Amina zwischen Feld und Beruf]
Amina arbeitet jeden Tag auf dem Feld, \textbf{weil} ihre Familie das Geld braucht. Sie verkauft auch Gemüse auf dem Markt, \textbf{denn} sie will ihren Kindern die Schule bezahlen. Sie hat wenig Zeit, \textbf{deshalb} steht sie schon um fünf Uhr morgens auf. \textbf{Obwohl} die Arbeit sehr hart ist, ist Amina stolz auf ihr Leben. Viele Frauen im Dorf arbeiten wie sie, \textbf{weil} die Männer oft in der Stadt Arbeit suchen. Die Frauen sind müde, \textbf{deshalb} fordern sie jetzt mehr Hilfe und mehr Rechte.
\end{textebox}

\begin{definition}[Glossaire]
\begin{itemize}
\item \all{das Feld, die Felder} : le champ ; \all{auf dem Feld arbeiten} : travailler aux champs
\item \all{das Gemüse (kein Plural / nur Singular)} : les légumes ; \all{der Markt, die Märkte} : le marché
\item \all{aufstehen (steht auf, stand auf, ist aufgestanden)} : se lever
\item \all{stolz (auf + Akk.)} : fier (de)
\item \all{das Dorf, die Dörfer} : le village ; \all{die Stadt, die Städte} : la ville
\item \all{fordern} : exiger, réclamer ; \all{das Recht, die Rechte} : le droit
\end{itemize}
\end{definition}

Observons la place du verbe conjugué dans chaque cas :

\begin{center}
\begin{tabularx}{\linewidth}{|l|X|l|}
\hline
\rowcolor{popblueL} Connecteur & Phrase & Place du verbe \\
\hline
\all{weil} & \all{..., weil ihre Familie das Geld \textbf{braucht}.} & à la fin \\
\hline
\all{denn} & \all{..., denn sie \textbf{will} ihren Kindern die Schule bezahlen.} & 2e position \\
\hline
\all{deshalb} & \all{..., deshalb \textbf{steht} sie ... auf.} & juste après (inversion) \\
\hline
\all{obwohl} & \all{Obwohl die Arbeit sehr hart \textbf{ist}, ...} & à la fin \\
\hline
\end{tabularx}
\end{center}

Quatre connecteurs, trois constructions différentes. Voyons chacun en détail.

\courssub{weil (parce que) : le subordonnant de cause}

\begin{propriete}[Règle : weil]
\all{weil} (parce que) introduit une \cle{subordonnée de cause}. Le verbe conjugué se place \textbf{à la fin} de la subordonnée. Si la subordonnée précède la principale, celle-ci commence par son verbe (inversion).
\begin{itemize}
\item \all{Sie arbeitet, weil sie Geld verdienen will.} — principale + subordonnée.
\item \all{Weil sie Geld verdienen will, arbeitet sie.} — subordonnée en tête : la principale inverse (verbe puis sujet).
\end{itemize}
Attention : s'il y a un infinitif ou un participe, le verbe conjugué passe \emph{après} eux : \all{..., weil sie Geld \emph{verdienen} \textbf{will}} (et non \all{will verdienen}).
\end{propriete}

\begin{exemplebox}[Exemples traduits]
\begin{itemize}
\item \all{Sie arbeitet, weil sie Geld verdienen will.} « Elle travaille parce qu'elle veut gagner de l'argent. »
\item \all{Viele Mädchen gehen nicht zur Schule, weil die Familien arm sind.} « Beaucoup de filles ne vont pas à l'école parce que les familles sont pauvres. »
\item \all{Weil die Frauen im Dorf viel arbeiten, haben sie wenig Freizeit.} « Parce que les femmes du village travaillent beaucoup, elles ont peu de temps libre. »
\item \all{Er bleibt zu Hause, weil er krank geworden ist.} « Il reste à la maison parce qu'il est tombé malade. » (au Perfekt : le participe \all{geworden} précède l'auxiliaire \all{ist})
\end{itemize}
\end{exemplebox}

\courssub{denn (car) : le coordonnant de cause}

\begin{propriete}[Règle : denn]
\all{denn} (car) est une \cle{conjonction de coordination} : il relie deux phrases principales. Il occupe la « position zéro » et ne compte pas dans la phrase : l'ordre des mots reste \textbf{normal}, le verbe conjugué garde sa deuxième position. Il n'y a \textbf{pas d'inversion} après \all{denn}.
\par
Contrairement à \all{weil}, \all{denn} ne peut jamais commencer une phrase : la proposition avec \all{denn} suit toujours la première.
\end{propriete}

\begin{exemplebox}[Exemples traduits]
\begin{itemize}
\item \all{Sie arbeitet, denn sie will Geld verdienen.} « Elle travaille, car elle veut gagner de l'argent. »
\item \all{Amina steht früh auf, denn der Markt öffnet um sechs Uhr.} « Amina se lève tôt, car le marché ouvre à six heures. »
\item \all{Die Frauen fordern mehr Rechte, denn sie tragen die Hauptlast der Arbeit.} « Les femmes réclament plus de droits, car elles portent la charge principale du travail. »
\end{itemize}
\end{exemplebox}

\begin{attention}[Piège classique du Bac : « car » n'est pas « weil » + ordre normal]
Les deux constructions suivantes sont fausses :
\begin{itemize}
\item \all{Sie arbeitet, weil sie will Geld verdienen.} — FAUX : avec \all{weil}, le verbe va à la fin.
\item \all{Denn sie will Geld verdienen, arbeitet sie.} — FAUX : \all{denn} ne peut pas ouvrir la phrase.
\end{itemize}
Réflexe à acquérir : si vous voulez garder l'ordre normal (verbe en 2e position), utilisez \all{denn} ; si vous utilisez \all{weil}, envoyez le verbe à la fin. En version, « car » se traduit presque toujours par \all{denn}, et « parce que » par \all{weil}.
\end{attention}

\courssub{deshalb (c'est pourquoi, donc) : l'adverbe de conséquence}

\begin{propriete}[Règle : deshalb]
\all{deshalb} (aussi \all{deswegen}, \all{darum}) n'est \textbf{pas} une conjonction mais un \cle{adverbe}. Placé en tête de phrase, il occupe la \textbf{première position} et provoque donc l'\textbf{inversion} : le verbe conjugué vient immédiatement après, puis le sujet.
\par
\all{Sie will Geld verdienen, deshalb arbeitet sie.} — jamais \all{deshalb sie arbeitet}.
\par
On peut aussi placer \all{deshalb} après le verbe : \all{Sie arbeitet deshalb sehr viel.} « C'est pourquoi elle travaille beaucoup. »
\end{propriete}

\begin{exemplebox}[Exemples traduits]
\begin{itemize}
\item \all{Sie will Geld verdienen, deshalb arbeitet sie.} « Elle veut gagner de l'argent, c'est pourquoi elle travaille. »
\item \all{Die Familie ist arm, deshalb kann das Mädchen nicht zur Schule gehen.} « La famille est pauvre, c'est pourquoi la fille ne peut pas aller à l'école. »
\item \all{Die Frauen sind müde, deshalb fordern sie mehr Hilfe.} « Les femmes sont fatiguées, donc elles réclament plus d'aide. »
\end{itemize}
\end{exemplebox}

\begin{attention}[deshalb n'est pas un subordonnant]
Erreur très fréquente des francophones, qui calquent « donc elle travaille » : \all{deshalb sie arbeitet} est impossible. Après \all{deshalb}, le verbe vient \emph{tout de suite} : \all{deshalb arbeitet sie}. Retenez le schéma : \textbf{deshalb + verbe + sujet}. De même, ne confondez pas la logique : \all{weil} et \all{denn} introduisent la \emph{cause}, \all{deshalb} introduit la \emph{conséquence}. « Elle est fatiguée, donc elle dort » = \all{Sie ist müde, deshalb schläft sie} — et non \all{weil sie müde ist} !
\end{attention}

\courssub{obwohl (bien que) : le subordonnant de concession}

\begin{propriete}[Règle : obwohl]
\all{obwohl} (bien que, quoique) introduit une \cle{subordonnée de concession} : il exprime une opposition, un obstacle qui n'empêche pas l'action. Comme \all{weil}, c'est un subordonnant : le verbe conjugué va \textbf{à la fin}. Si la subordonnée est en tête, la principale inverse.
\par
Remarque culturelle : le français « bien que » exige le subjonctif ; l'allemand \all{obwohl} se construit avec l'\textbf{indicatif} — aucun mode spécial n'est nécessaire.
\end{propriete}

\begin{exemplebox}[Exemples traduits]
\begin{itemize}
\item \all{Obwohl sie müde ist, arbeitet sie.} « Bien qu'elle soit fatiguée, elle travaille. »
\item \all{Sie lächelt, obwohl die Arbeit hart ist.} « Elle sourit, bien que le travail soit dur. »
\item \all{Obwohl viele Frauen keinen Beruf gelernt haben, ernähren sie ihre Familien.} « Bien que beaucoup de femmes n'aient pas appris de métier, elles nourrissent leurs familles. »
\item \all{Obwohl er krank ist, geht er zur Arbeit.} « Bien qu'il soit malade, il va au travail. »
\end{itemize}
\end{exemplebox}

\courssub{Tableau récapitulatif et comparaison avec le français}

\begin{center}
\begin{tabularx}{\linewidth}{|l|l|l|X|}
\hline
\rowcolor{popblueL} Connecteur & Sens & Type & Place du verbe \\
\hline
\all{weil} & parce que (cause) & subordonnant & à la fin \\
\hline
\all{denn} & car (cause) & coordonnant & 2e position (ordre normal) \\
\hline
\all{deshalb} & c'est pourquoi, donc (conséquence) & adverbe & juste après (inversion) \\
\hline
\all{obwohl} & bien que (concession) & subordonnant & à la fin \\
\hline
\end{tabularx}
\end{center}

\begin{center}
\begin{tabularx}{\linewidth}{|X|X|X|}
\hline
\rowcolor{popblueL} Français & Deutsch & Remarque \\
\hline
Elle travaille \textbf{parce qu'}elle veut gagner de l'argent. & \all{Sie arbeitet, \textbf{weil} sie Geld verdienen \textbf{will}.} & verbe à la fin \\
\hline
Elle travaille, \textbf{car} elle veut gagner de l'argent. & \all{Sie arbeitet, \textbf{denn} sie \textbf{will} Geld verdienen.} & ordre normal \\
\hline
Elle veut gagner de l'argent, \textbf{c'est pourquoi} elle travaille. & \all{Sie will Geld verdienen, \textbf{deshalb arbeitet} sie.} & inversion \\
\hline
\textbf{Bien qu'}elle soit fatiguée, elle travaille. & \all{\textbf{Obwohl} sie müde \textbf{ist}, \textbf{arbeitet} sie.} & fin + inversion dans la principale \\
\hline
\end{tabularx}
\end{center}

\begin{popfigure}
\begin{tikzpicture}[
  conn/.style={rectangle, rounded corners=3pt, draw=popblue, fill=popblueL, align=center, font=\small, minimum height=1.1cm},
  lab/.style={font=\footnotesize\itshape, text=popmuted}
]
\node[conn, align=center] (weil) at (0,2.4){\all{weil}\\parce que};
\node[conn, align=center] (denn) at (3.4,2.4){\all{denn}\\car};
\node[conn, draw=poporange, fill=poporangeL, align=center] (deshalb) at (6.8,2.4){\all{deshalb}\\c'est pourquoi};
\node[conn, draw=poppurple, fill=poppurpleL, align=center] (obwohl) at (10.2,2.4){\all{obwohl}\\bien que};
\node[lab] at (0,1.5) {subordonnant};
\node[lab] at (3.4,1.5) {coordonnant};
\node[lab] at (6.8,1.5) {adverbe};
\node[lab] at (10.2,1.5) {subordonnant};
\node[lab, text=popdark] at (0,0.9) {verbe à la fin};
\node[lab, text=popdark] at (3.4,0.9) {verbe en 2e position};
\node[lab, text=popdark] at (6.8,0.9) {inversion sujet-verbe};
\node[lab, text=popdark] at (10.2,0.9) {verbe à la fin};
\end{tikzpicture}
\legende{Les quatre connecteurs : sens, nature grammaticale et construction exigée.}
\end{popfigure}

\begin{popfigure}
\begin{tikzpicture}[
  box/.style={rectangle, rounded corners=4pt, align=center, font=\small, minimum height=1cm},
  fleche/.style={-{Stealth[length=3mm]}, thick}
]
\node[box, draw=popgreen, fill=popgreenL, align=center] (cause) at (0,0){\textbf{CAUSE}\\\all{weil} / \all{denn}\\Elle veut gagner\\de l'argent};
\node[box, draw=poporange, fill=poporangeL, align=center] (cons) at (6.2,0){\textbf{CONSÉQUENCE}\\\all{deshalb}\\donc elle travaille};
\node[box, draw=poppurple, fill=poppurpleL, align=center] (obst) at (0,-2.6){\textbf{OBSTACLE}\\\all{obwohl}\\bien qu'elle soit fatiguée};
\node[box, draw=popblue, fill=popblueL, align=center] (res) at (6.2,-2.6){\textbf{RÉSULTAT}\\elle travaille\\quand même};
\draw[fleche, popgreen] (cause) -- (cons);
\draw[fleche, poppurple] (obst) -- (res);
\end{tikzpicture}
\legende{Relations logiques : \all{weil}/\all{denn} expriment la cause, \all{deshalb} la conséquence, \all{obwohl} l'opposition.}
\end{popfigure}

\courssub{Méthode : choisir le bon connecteur}

\begin{methode}[Trois étapes pour ne plus se tromper]
\begin{enumerate}
\item \textbf{Identifier la relation logique.} Est-ce une cause (pourquoi ?), une conséquence (résultat ?) ou une concession (malgré quoi ?) ? Cause $\rightarrow$ \all{weil} ou \all{denn} ; conséquence $\rightarrow$ \all{deshalb} ; concession $\rightarrow$ \all{obwohl}.
\item \textbf{Vérifier la construction exigée.} Subordonnant (\all{weil}, \all{obwohl}) : verbe à la fin. Coordonnant (\all{denn}) : ordre normal. Adverbe (\all{deshalb}) : inversion immédiate.
\item \textbf{Relire la phrase complète.} Contrôlez la place de chaque verbe, la virgule avant \all{weil}/\all{denn}/\all{obwohl}, et la majuscule aux noms.
\end{enumerate}
\end{methode}

\begin{exoresolu}[Relier deux phrases]
Énoncé : reliez les deux phrases avec le connecteur indiqué. Phrases : \all{Die Frau arbeitet auf dem Feld. Sie will ihre Kinder ernähren.}
\begin{enumerate}
\item avec \all{weil} ; 2. avec \all{denn} ; 3. avec \all{deshalb} ; 4. avec \all{obwohl} (remplacez la deuxième phrase par : \all{Sie ist müde}).
\end{enumerate}
\tcblower
Solution détaillée :
\begin{enumerate}
\item Cause avec subordonnant, verbe à la fin : \all{Die Frau arbeitet auf dem Feld, weil sie ihre Kinder ernähren will.} « La femme travaille aux champs parce qu'elle veut nourrir ses enfants. » (L'infinitif \all{ernähren} précède le verbe modal \all{will}.)
\item Cause avec coordonnant, ordre normal : \all{Die Frau arbeitet auf dem Feld, denn sie will ihre Kinder ernähren.} « …, car elle veut nourrir ses enfants. »
\item Conséquence avec inversion : \all{Die Frau will ihre Kinder ernähren, deshalb arbeitet sie auf dem Feld.} « La femme veut nourrir ses enfants, c'est pourquoi elle travaille aux champs. » (On inverse l'ordre logique : la cause d'abord, la conséquence ensuite.)
\item Concession, verbe à la fin, puis inversion dans la principale : \all{Obwohl sie müde ist, arbeitet die Frau auf dem Feld.} « Bien qu'elle soit fatiguée, la femme travaille aux champs. »
\end{enumerate}
\end{exoresolu}

\begin{exercice}[À vous de jouer]
Reliez chaque paire de phrases avec le connecteur entre parenthèses, puis traduisez en français.
\begin{enumerate}
\item \all{Viele Mädchen bleiben zu Hause. Sie müssen im Haushalt helfen.} (\all{weil})
\item \all{Amina verkauft Gemüse auf dem Markt. Sie braucht Geld für die Schule ihrer Kinder.} (\all{denn})
\item \all{Die Frauen im Dorf haben wenig Freizeit. Sie sind oft sehr müde.} (\all{deshalb})
\item \all{Die Arbeit ist hart. Die Frauen sind stolz auf ihre Familie.} (\all{obwohl})
\item Traduisez en allemand : « Elle va en ville, car elle cherche du travail. » puis « Elle cherche du travail, c'est pourquoi elle va en ville. »
\end{enumerate}
\end{exercice}

\begin{corrige}[Exercice « À vous de jouer »]
\begin{enumerate}
\item \all{Viele Mädchen bleiben zu Hause, weil sie im Haushalt helfen müssen.} « Beaucoup de filles restent à la maison parce qu'elles doivent aider au ménage. »
\item \all{Amina verkauft Gemüse auf dem Markt, denn sie braucht Geld für die Schule ihrer Kinder.} « Amina vend des légumes au marché, car elle a besoin d'argent pour l'école de ses enfants. »
\item \all{Die Frauen im Dorf haben wenig Freizeit, deshalb sind sie oft sehr müde.} « Les femmes du village ont peu de temps libre, c'est pourquoi elles sont souvent très fatiguées. »
\item \all{Obwohl die Arbeit hart ist, sind die Frauen stolz auf ihre Familie.} « Bien que le travail soit dur, les femmes sont fières de leur famille. »
\item \all{Sie geht in die Stadt, denn sie sucht Arbeit.} / \all{Sie sucht Arbeit, deshalb geht sie in die Stadt.}
\end{enumerate}
\end{corrige}

\begin{aretenir}[L'essentiel]
\begin{itemize}
\item \all{weil} et \all{obwohl} : subordonnants $\rightarrow$ verbe conjugué \textbf{à la fin}.
\item \all{denn} : coordonnant $\rightarrow$ ordre \textbf{normal}, verbe en 2e position ; jamais en tête de phrase.
\item \all{deshalb} : adverbe $\rightarrow$ \textbf{inversion} : \all{deshalb + verbe + sujet}.
\item Logique : \all{weil}/\all{denn} = cause ; \all{deshalb} = conséquence ; \all{obwohl} = concession (avec indicatif, sans subjonctif !).
\end{itemize}
\end{aretenir}

\coursec{Conjugaison : présent et verbes de modalité}

Dans la section précédente, nous avons appris à relier nos idées avec \all{weil}, \all{denn}, \all{deshalb} et \all{obwohl}. Mais pour que ces phrases soient correctes, il faut des verbes bien conjugués. Revenons au texte sur la situation de la femme : on y lit \all{Die Frau arbeitet viel}, \all{sie liest Zeitung}, \all{sie kann nicht studieren}, \all{weil sie auf dem Feld arbeiten muss}. Pourquoi \all{arbeitet} mais \all{liest} ? Pourquoi \all{muss} se trouve-t-il tantôt en deuxième position, tantôt à la fin ? C'est l'objet de cette section : maîtriser le présent des verbes faibles et forts, puis les verbes de modalité, outils indispensables pour parler des droits, des devoirs et des possibilités des femmes.

\courssub{Le présent des verbes faibles}

\begin{definition}[Verbe faible]
Un \cle{verbe faible} (\all{schwaches Verb}) est un verbe \emph{régulier} : son radical ne change jamais. On ajoute au radical les terminaisons du présent : \textbf{-e, -st, -t, -en, -t, -en}. Exemple : \all{machen} → \all{ich mache, du machst, er macht, wir machen, ihr macht, sie machen}.
\end{definition}

\begin{propriete}[Terminaisons du présent]
\begin{center}
\begin{tabular}{|l|l|l|}
\hline
\rowcolor{popblueL} Personne & Terminaison & Exemple : machen \\
\hline
ich & -e & ich mache \\
\hline
du & -st & du machst \\
\hline
er / sie / es & -t & sie macht \\
\hline
wir & -en & wir machen \\
\hline
ihr & -t & ihr macht \\
\hline
sie / Sie & -en & sie machen \\
\hline
\end{tabular}
\end{center}
\all{Sie macht den Haushalt.} « Elle fait le ménage. » — \all{Wir lernen Deutsch.} « Nous apprenons l'allemand. »
\end{propriete}

\begin{attention}[Les verbes dont le radical finit en -t ou -d]
Un verbe comme \all{arbeiten} (radical \all{arbeit-}) ne peut pas recevoir directement \all{-st} ou \all{-t} : on intercale un \textbf{e} de liaison. On obtient donc \all{du arbeit\textbf{est}}, \all{er arbeit\textbf{et}}, \all{ihr arbeit\textbf{et}}. Même phénomène avec \all{bilden} → \all{du bild\textbf{est}}, \all{er bild\textbf{et}} ; \all{warten} → \all{du wartest}. Erreur fréquente des francophones : écrire \all{du arbeitst} — c'est faux !
\end{attention}

\begin{attention}[Cas particuliers : verbes en -eln et -ern ; haben, sein, werden]
\begin{itemize}
\item Verbes en \textbf{-eln} : à la 1\textsuperscript{re} personne du singulier, le \textbf{e} du radical \emph{peut} s'élider (tomber) : \all{sammeln} → \all{ich sammel} (forme courante) ou \all{ich sammele} (forme complète) ; \all{du sammelst}, \all{wir sammeln}.
\item Verbes en \textbf{-ern} : réguliers : \all{erwidern} → \all{ich erwidere}.
\item \all{haben} : \all{ich habe, du hast, er hat, wir haben, ihr habt, sie haben}.
\item \all{sein} (irrégulier, à connaître par cœur) : \all{ich bin, du bist, er ist, wir sind, ihr seid, sie sind}.
\item \all{werden} : \all{ich werde, du wirst, er wird, wir werden, ihr werdet, sie werden}. \all{Sie will Ärztin werden.} « Elle veut devenir médecin. »
\end{itemize}
\end{attention}

\courssub{Le présent des verbes forts}

Observons d'abord trois phrases tirées de notre thème :

\all{Die Frau liest ein Buch.} « La femme lit un livre. »

\all{Sie fährt mit dem Bus in die Stadt.} « Elle va en ville en bus. »

\all{Sie spricht mit den Nachbarn.} « Elle parle avec les voisines. »

Le radical a changé : \all{lesen} → \all{liest}, \all{fahren} → \all{fährt}, \all{sprechen} → \all{spricht}. Ce changement n'a lieu qu'à la \textbf{2\textsuperscript{e} et à la 3\textsuperscript{e} personne du singulier} (\all{du} et \all{er/sie/es}).

\begin{propriete}[Le changement de voyelle des verbes forts]
Les \cle{verbes forts} (\all{starke Verben}) modifient leur voyelle radicale à la 2\textsuperscript{e} et 3\textsuperscript{e} personne du singulier, selon trois schémas principaux :
\begin{itemize}
\item \textbf{e → i} : \all{sprechen} → \all{du sprichst, er spricht} ; \all{helfen} → \all{er hilft} ; \all{essen} → \all{er isst} ; \all{geben} → \all{er gibt}.
\item \textbf{e → ie} : \all{lesen} → \all{du liest, sie liest} ; \all{sehen} → \all{er sieht} ; \all{befehlen} → \all{er befiehlt}.
\item \textbf{a → ä} : \all{fahren} → \all{du fährst, sie fährt} ; \all{schlafen} → \all{er schläft} ; \all{tragen} → \all{sie trägt} ; \all{waschen} → \all{er wäscht}.
\end{itemize}
Aux autres personnes (\all{ich, wir, ihr, sie}), le radical reste identique à l'infinitif : \all{ich lese, wir lesen, ihr lest}. Les verbes faibles, eux, ne changent \emph{jamais} leur voyelle.
\end{propriete}

\begin{popfigure}
\begin{center}
\begin{tikzpicture}[font=\small]
\foreach \x/\verb/\voy in {0/arbeiten/-, 4.6/lesen/e$\rightarrow$ie, 9.2/fahren/a$\rightarrow$\"a}{
  \node[fill=popblue, text=white, rounded corners=2pt, minimum width=3.6cm, font=\bfseries] at (\x+1.8,5.4) {\verb};
}
\foreach \y/\p in {4.7/ich, 4.0/du, 3.3/er\,sie\,es, 2.6/wir, 1.9/ihr, 1.2/sie}{
  \node[fill=popcream, draw=popline, rounded corners=1pt, minimum width=1.2cm] at (0.6,\y) {\p};
  \node[fill=popcream, draw=popline, rounded corners=1pt, minimum width=1.2cm] at (5.2,\y) {\p};
  \node[fill=popcream, draw=popline, rounded corners=1pt, minimum width=1.2cm] at (9.8,\y) {\p};
}
\node at (2.6,4.7) {arbeite};
\node at (2.6,4.0) {arbeit\textcolor{poporange}{\bfseries est}};
\node at (2.6,3.3) {arbeit\textcolor{poporange}{\bfseries et}};
\node at (2.6,2.6) {arbeiten};
\node at (2.6,1.9) {arbeit\textcolor{poporange}{\bfseries et}};
\node at (2.6,1.2) {arbeiten};
\node at (7.2,4.7) {lese};
\node at (7.2,4.0) {l\textcolor{poporange}{\bfseries ie}st};
\node at (7.2,3.3) {l\textcolor{poporange}{\bfseries ie}st};
\node at (7.2,2.6) {lesen};
\node at (7.2,1.9) {lest};
\node at (7.2,1.2) {lesen};
\node at (11.8,4.7) {fahre};
\node at (11.8,4.0) {f\textcolor{poporange}{\bfseries \"a}hrst};
\node at (11.8,3.3) {f\textcolor{poporange}{\bfseries \"a}hrt};
\node at (11.8,2.6) {fahren};
\node at (11.8,1.9) {fahrt};
\node at (11.8,1.2) {fahren};
\end{tikzpicture}
\end{center}
\legende{Conjugaison au présent : \all{arbeiten} (faible, avec e de liaison), \all{lesen} et \all{fahren} (forts, voyelle modifiée en orange).}
\end{popfigure}

\begin{exemplebox}[Verbes forts dans notre thème]
\all{Das Mädchen liest jeden Tag.} « La fillette lit chaque jour. »

\all{Die B\"auerin tr\"agt Wasser vom Fluss.} « La paysanne porte l'eau de la rivière. »

\all{Er hilft seiner Frau im Haushalt.} « Il aide sa femme dans les tâches ménagères. »

\all{Sie spricht drei Sprachen.} « Elle parle trois langues. »
\end{exemplebox}

\courssub{Les verbes de modalité au présent}

Pour parler de la situation de la femme — ce qu'elle \emph{peut} faire, ce qu'elle \emph{doit} faire, ce qu'elle a le \emph{droit} de faire, ce qu'elle \emph{veut} —, l'allemand utilise six verbes de modalité (plus \all{möchte}). Ce sont des verbes irréguliers : le singulier (\all{ich, du, er}) a un radical modifié, et la 1\textsuperscript{re} et la 3\textsuperscript{e} personne du singulier sont identiques (\all{ich kann, er kann}).


\begin{center}
\begin{tabularx}{\linewidth}{|l|X|X|X|}
\hline
\rowcolor{popblueL} Infinitif & ich / er & Sens & Exemple \\
\hline
k\"onnen & kann & pouvoir, savoir (capacité) & Sie kann lesen. \\
\hline
m\"ussen & muss & devoir (nécessité) & Sie muss arbeiten. \\
\hline
d\"urfen & darf & avoir le droit (permission) & Sie darf studieren. \\
\hline
wollen & will & vouloir (volonté) & Sie will \"Arztin werden. \\
\hline
sollen & soll & devoir (ordre, conseil) & Sie soll mehr lernen. \\
\hline
m\"ogen & mag & aimer & Sie mag ihren Beruf. \\
\hline
m\"ochte & m\"ochte & voudrais (politesse) & Ich m\"ochte Lehrerin werden. \\
\hline
\end{tabularx}
\end{center}

\begin{propriete}[La construction avec verbe de modalité (phrase principale)]
Avec un verbe de modalité, la phrase allemande forme une \emph{embrassade} (\all{Satzklammer}) : le \cle{modal conjugué} occupe la \textbf{2\textsuperscript{e} position} et le \cle{verbe lexical} reste à l'\textbf{infinitif}, rejeté à la \textbf{toute fin} de la phrase.

\all{Die Frau \textbf{muss} viel \textbf{arbeiten}.} « La femme doit beaucoup travailler. »

\all{Die M\"adchen \textbf{k\"onnen} heute zur Schule \textbf{gehen}.} « Les filles peuvent aller à l'école aujourd'hui. »

Attention : en français, l'infinitif suit immédiatement le modal (« elle doit travailler ») ; en allemand, tout le reste de la phrase passe \emph{avant} l'infinitif.
\end{propriete}

\begin{popfigure}
\begin{center}
\begin{tikzpicture}[font=\small, node distance=0.4cm]
\node[fill=popblue, text=white, rounded corners=2pt, minimum width=2.2cm, align=center] (s) {SUJET\\Die Frau};
\node[fill=poporange, text=white, rounded corners=2pt, minimum width=2.6cm, align=center, right=of s] (m) {MODAL (2\textsuperscript{e} pos.)\\muss};
\node[fill=popcream, draw=popline, rounded corners=2pt, minimum width=3.2cm, align=center, right=of m] (c) {COMPLÉMENTS\\auf dem Feld viel};
\node[fill=popgreen, text=white, rounded corners=2pt, minimum width=2.8cm, align=center, right=of c] (i) {INFINITIF (fin)\\arbeiten};
\draw[-{Stealth}, thick, popdark] (s) -- (m);
\draw[-{Stealth}, thick, popdark] (m) -- (c);
\draw[-{Stealth}, thick, popdark] (c) -- (i);
\node[poppurple, align=center, font=\bfseries\small] at (6.5, 1.8) {La Satzklammer :\\l'embrassade verbale};
\draw[poppurple, thick, -{Stealth[length=3mm]}] (5.5, 1.5) -- (m.north);
\draw[poppurple, thick, -{Stealth[length=3mm]}] (7.5, 1.5) -- (i.north);
\end{tikzpicture}
\end{center}
\legende{La construction modale : le modal conjugué en 2\textsuperscript{e} position, l'infinitif à la fin.}
\end{popfigure}

\begin{methode}[Construire une phrase avec un verbe de modalité]
\begin{enumerate}
\item Identifier le \textbf{sujet} (qui fait l'action ?) → le placer en position 1.
\item Conjuguer le \textbf{verbe de modalité} selon le sujet → le placer en \textbf{position 2}.
\item Placer les \textbf{compléments} (lieu, temps, manière, objet) → positions 3, 4...
\item Rejeter le \textbf{verbe lexical à l'infinitif} tout à la \textbf{fin de la phrase}.
\item \textbf{Dans une subordonnée (weil, obwohl...)} : placer l'infinitif \textbf{juste avant} le modal conjugué, tous les deux à la fin.
\end{enumerate}
\end{methode}

\begin{propriete}[Le modal dans la subordonnée (weil, obwohl, dass...)]
Dans une subordonnée, le verbe conjugué va à la fin. Si le verbe conjugué est un \textbf{modal}, c'est donc le modal qui occupe la dernière place, et l'infinitif se place juste \emph{avant} lui :

\all{Die Frau ist m\"ude, weil sie auf dem Feld arbeiten \textbf{muss}.} « La femme est fatiguée, parce qu'elle doit travailler aux champs. »

\all{Viele M\"adchen bleiben zu Hause, obwohl sie lernen \textbf{wollen}.} « Beaucoup de filles restent à la maison, bien qu'elles veuillent apprendre. »

Ordre en fin de subordonnée : \textbf{infinitif + modal conjugué}.
\end{propriete}

\begin{exemplebox}[Les modalités au service du thème]
\all{Sie kann lesen und schreiben.} « Elle sait lire et écrire. » (capacité)

\all{Frauen d\"urfen in Deutschland studieren.} « Les femmes ont le droit d'étudier en Allemagne. » (droit)

\all{Die B\"auerin muss fr\"uh aufstehen.} « La paysanne doit se lever tôt. » (obligation)

\all{Sie will Lehrerin werden.} « Elle veut devenir enseignante. » (volonté)

\all{Die Regierung soll den Frauen mehr helfen.} « Le gouvernement devrait davantage aider les femmes. » (devoir moral)
\end{exemplebox}

\begin{attention}[k\"onnen ou d\"urfen ?]
\all{k\"onnen} exprime la \textbf{capacité} (pouvoir, savoir faire) ; \all{d\"urfen} exprime la \textbf{permission} (avoir le droit). Ne les confondez pas :

\all{Sie \textbf{kann} nicht arbeiten.} = Elle ne \emph{peut} pas travailler (elle est malade, empêchée physiquement).

\all{Sie \textbf{darf} nicht arbeiten.} = Elle n'a \emph{pas le droit} de travailler (on le lui interdit).

Cette nuance est essentielle pour parler des droits des femmes : dans certains villages, une fille \all{kann} lesen (elle en est capable), mais sie \all{darf} nicht zur Schule gehen (on ne l'y autorise pas).
\end{attention}

\begin{savaistu}[m\"ochte : la politesse allemande]
\all{m\"ochte} (« je voudrais ») est le Konjunktiv II de \all{m\"ogen}. C'est la forme la plus courante pour exprimer poliment un souhait, beaucoup plus fréquente que \all{ich will}, qui peut sembler brutal. \all{Ich m\"ochte Lehrerin werden.} « Je voudrais devenir enseignante. » — \all{M\"ochten Sie einen Kaffee?} « Voudriez-vous un café ? » Au marché de Yaoundé comme à Berlin, on dit \all{Ich m\"ochte...} et jamais \all{Ich will...} pour demander quelque chose !
\end{savaistu}

\begin{savaistu}[Le récit au passé (Präteritum) dans le texte]
Dans le texte \guillemotleft Amina's Traum \guillemotright, vous rencontrerez deux formes du \emph{Präteritum} (passé narratif), temps utilisé à l'écrit pour raconter une histoire :
\begin{itemize}
\item \all{konnten} (de \all{k\"onnen}) : \all{Früher konnten die Mädchen hier nicht lernen.} « Autrefois, les filles ne pouvaient pas apprendre ici. »
\item \all{fand} (de \all{finden}) : \all{Auch Aminas Vater fand: Bildung ist wichtig...} « Le père d'Amina trouvait aussi : l'éducation est importante... »
\end{itemize}
Pour l'instant, reconnaissez-les ; leur formation sera étudiée ultérieurement. Au présent, ce serait : \all{können} / \all{findet}.
\end{savaistu}

\begin{textebox}[Aminas Traum — Le rêve d'Amina]
Amina lebt in einem kleinen Dorf in der Nähe von Bamenda. Sie ist vierzehn Jahre alt und geht gern zur Schule. Sie kann schon gut lesen und schreiben, und sie lernt auch Deutsch. „Ich m\"ochte Ärztin werden“, sagt sie oft. „Ich will den Menschen in meinem Dorf helfen.“

Aber das Leben im Dorf ist schwer. Jeden Morgen muss Amina vor der Schule Wasser holen und ihrer Mutter im Haushalt helfen. Am Nachmittag arbeitet sie mit den anderen Frauen auf dem Feld. Sie tr\"agt schwere Körbe und ist am Abend sehr m\"ude. Trotzdem liest sie jeden Abend ihre Schulbücher, weil sie lernen will.

Ihre Mutter ist stolz auf sie. „Meine Tochter darf die Schule besuchen“, sagt sie. „Früher konnten die Mädchen hier nicht lernen. Heute ist das anders.“ Auch Aminas Vater findet: Bildung ist wichtig für Mädchen und für Jungen.

In der Stadt ist die Situation der Frau oft besser: Frauen können dort studieren und einen Beruf lernen. Sie dürfen wählen und arbeiten. Aber auf dem Land müssen viele Frauen noch sehr hart arbeiten und verdienen wenig Geld.

Amina hat einen Traum: Sie m\"ochte später in Yaoundé Medizin studieren. „Ich kann es schaffen“, sagt sie, „weil ich stark bin und weil meine Familie mich unterstützt.“

\textbf{Glossaire :} das Dorf, die Dörfer (le village) — der Traum, die Träume (le rêve) — Wasser holen (aller chercher de l'eau) — der Korb, die Körbe (le panier) — stolz (fier, fière) — die Schule besuchen (fréquenter l'école) — früher (autrefois) — wählen (voter) — verdienen (gagner) — es schaffen (y arriver) — unterstützen (soutenir).

\textbf{Questions :} 1. Was m\"ochte Amina werden? 2. Was muss sie jeden Morgen machen? 3. Warum liest sie am Abend ihre Bücher? 4. Was können Frauen in der Stadt machen?
\end{textebox}

\begin{exoresolu}[Conjuguez au présent]
Conjuguez les verbes entre parenthèses au présent :

1. Die Frau (arbeiten) \_\_\_\_ viel auf dem Feld.

2. Sie (lesen) \_\_\_\_ jeden Abend ein Buch.

3. Die Mädchen (fahren) \_\_\_\_ mit dem Bus in die Stadt.

4. Du (können) \_\_\_\_ gut Deutsch sprechen.

5. Traduisez : « Elle ne peut pas étudier, parce qu'elle doit aider sa mère. »

\tcblower

\textbf{Solution détaillée :}

1. \all{Die Frau \textbf{arbeitet} viel auf dem Feld.} — Verbe faible ; radical en \all{-t} → e de liaison : \all{arbeitet}.

2. \all{Sie \textbf{liest} jeden Abend ein Buch.} — Verbe fort, changement e → ie à la 3\textsuperscript{e} personne du singulier.

3. \all{Die Mädchen \textbf{fahren} mit dem Bus in die Stadt.} — Sujet pluriel : pas de changement de voyelle (\all{fährt} serait pour \all{sie} au singulier).

4. \all{Du \textbf{kannst} gut Deutsch sprechen.} — Modal conjugué en 2\textsuperscript{e} position ; l'infinitif \all{sprechen} va à la fin.

5. \all{Sie \textbf{kann} nicht studieren, weil sie ihrer Mutter helfen \textbf{muss}.} — Dans la proposition principale, \all{kann} est en 2\textsuperscript{e} position et \all{studieren} à la fin ; dans la subordonnée en \all{weil}, l'ordre final est infinitif (\all{helfen}) + modal conjugué (\all{muss}). \all{ihrer Mutter} (datif après \all{helfen}).
\end{exoresolu}

\begin{attention}[Les trois erreurs classiques des francophones]
\begin{itemize}
\item \textbf{Oublier l'infinitif final} : ne dites jamais \all{Sie muss arbeiten viel}. Le français place l'infinitif juste après le modal ; l'allemand le renvoie à la fin : \all{Sie muss viel arbeiten}.
\item \textbf{Conjuguer le verbe lexical} : après un modal, le deuxième verbe reste à l'infinitif : \all{Sie kann lesen} (et non \all{Sie kann liest}).
\item \textbf{Oublier le changement de voyelle} : \all{er lesen} est impossible ; il faut \all{er liest}. Apprenez les verbes forts avec leur forme à la 3\textsuperscript{e} personne : \all{lesen – liest, fahren – fährt, sprechen – spricht, helfen – hilft, sehen – sieht, schlafen – schläft, tragen – trägt, essen – isst}.
\end{itemize}
\end{attention}

\begin{exercice}[À vous de conjuguer et de traduire]
1. Conjuguez : \all{sprechen} (er), \all{helfen} (du), \all{sehen} (sie, sing.), \all{arbeiten} (ihr), \all{müssen} (wir).

2. Traduisez en allemand : a) La femme veut apprendre l'allemand. b) Les filles ont le droit d'aller à l'école. c) Elle est fatiguée parce qu'elle doit beaucoup travailler. d) Je voudrais devenir enseignante.
\end{exercice}

\begin{corrige}[Exercice : À vous de conjuguer et de traduire]
1. \all{er spricht} ; \all{du hilfst} ; \all{sie sieht} ; \all{ihr arbeitet} ; \all{wir müssen}.

2. a) \all{Die Frau will Deutsch lernen.} b) \all{Die Mädchen dürfen zur Schule gehen.} c) \all{Sie ist müde, weil sie viel arbeiten muss.} d) \all{Ich möchte Lehrerin werden.}
\end{corrige}

\begin{aretenir}[L'essentiel de la section]
\begin{itemize}
\item Verbes faibles : radical invariable + \all{-e, -st, -t, -en, -t, -en} ; e de liaison après \all{-t/-d} (\all{du arbeitest}).
\item Verbes forts : voyelle modifiée à la 2\textsuperscript{e} et 3\textsuperscript{e} personne du singulier (\all{liest, fährt, spricht, hilft, sieht, schläft, trägt, isst}).
\item Modalités : \all{kann, muss, darf, will, soll, mag} ; \all{möchte} = « voudrais » (politesse).
\item Construction phrase principale : modal conjugué en 2\textsuperscript{e} position + infinitif à la fin.
\item Construction subordonnée (\all{weil, obwohl}) : infinitif + modal conjugué à la fin (\all{..., weil sie arbeiten muss}).
\item \all{können} = capacité ; \all{dürfen} = droit, permission.
\end{itemize}
\end{aretenir}

\coursec{Méthode du commentaire et production guidée}

Dans la section précédente, nous avons révisé le présent des verbes faibles et forts ainsi que les verbes de modalité. Ces outils grammaticaux ne sont pas une fin en soi : ils vont maintenant nous servir à \cle{produire}. Au baccalauréat, l'épreuve d'allemand demande de \cle{comprendre} un texte, de le \cle{commenter}, de le \cle{résumer} et de rédiger une courte \cle{production personnelle}. Cette dernière section vous donne une méthode simple, rigoureuse et directement applicable le jour de l'examen, en réemployant le texte support sur la situation de la femme, le lexique du thème et les connecteurs étudiés (\all{weil}, \all{denn}, \all{deshalb}, \all{obwohl}).

\courssub{La méthode du commentaire de texte au Bac}

\begin{methode}[Rédiger un commentaire de texte en trois étapes]
Un commentaire de texte allemand suit toujours la même architecture : \textbf{Einleitung} (introduction), \textbf{Hauptteil} (développement), \textbf{Schluss} (conclusion).
\begin{enumerate}
\item \textbf{Présenter le texte} (Einleitung) : indiquez le thème général avec la formule \all{Der Text handelt von...} (« Le texte traite de... »). Exemple : \all{Der Text handelt von der Situation der Frau in der Stadt und auf dem Land.} « Le texte traite de la situation de la femme en ville et au village. »
\item \textbf{Dégager 2 ou 3 idées en citant le texte} (Hauptteil) : utilisez \all{Im Text steht, dass...} (« Dans le texte, il est dit que... ») ou \all{Der Autor sagt, dass...} (« L'auteur dit que... »). Attention : après \all{dass}, le verbe conjugué va \textbf{en fin de phrase}. Exemple : \all{Im Text steht, dass viele Frauen auf dem Land auf dem Feld arbeiten.} « Dans le texte, il est dit que beaucoup de femmes au village travaillent aux champs. »
\item \textbf{Donner son avis personnel avec des connecteurs} (Schluss) : utilisez \all{Ich meine, dass... weil...} ou \all{Obwohl..., denke ich...}. Exemple : \all{Ich meine, dass die Gleichberechtigung wichtig ist, weil Frauen und Männer die gleichen Chancen haben müssen.} « Je pense que l'égalité des droits est importante, parce que les femmes et les hommes doivent avoir les mêmes chances. »
\end{enumerate}
\end{methode}

\begin{popfigure}
\begin{center}
\begin{tikzpicture}[
node distance=0.9cm,
etape/.style={rectangle, rounded corners, draw=popblue, fill=popblueL, thick, align=center, text width=6.2cm, minimum height=1.1cm},
fleche/.style={-{Stealth[length=3mm]}, thick, popdark}]
\node[etape, align=center] (e1){\textbf{1. Einleitung}\\ \all{Der Text handelt von...}\\ (texte + thème)};
\node[etape, below=of e1, fill=poporangeL, draw=poporange, align=center] (e2){\textbf{2. Hauptteil}\\ \all{Im Text steht, dass...}\\ (idées + arguments + exemples)};
\node[etape, below=of e2, fill=popgreenL, draw=popgreen, align=center] (e3){\textbf{3. Schluss}\\ \all{Ich meine, dass... weil...}\\ (avis personnel)};
\draw[fleche] (e1) -- (e2);
\draw[fleche] (e2) -- (e3);
\end{tikzpicture}
\end{center}
\legende{Organigramme du commentaire de texte : introduction, développement, conclusion.}
\end{popfigure}

\begin{aretenir}[Les trois règles d'or de la rédaction allemande]
\begin{itemize}
\item Le verbe conjugué est \textbf{toujours en 2e position} dans la phrase principale : \all{Viele Frauen arbeiten in der Stadt.} « Beaucoup de femmes travaillent en ville. »
\item Après \all{weil}, \all{obwohl}, \all{dass}, le verbe conjugué va \textbf{en fin de subordonnée} : \all{..., weil sie eine gute Schulbildung haben.} « ..., parce qu'elles ont une bonne formation scolaire. »
\item Après un verbe de modalité, l'infinitif va \textbf{en fin de phrase} : \all{Die Frauen müssen die gleichen Chancen haben.} « Les femmes doivent avoir les mêmes chances. »
\end{itemize}
\end{aretenir}

\courssub{Commentaire modèle rédigé}

\begin{textebox}[Commentaire modèle sur le texte support]
\all{Der Text handelt von der Situation der Frau in Stadt und Land. In der Stadt arbeiten viele Frauen in einem Beruf, weil sie eine gute Schulbildung haben. Obwohl sie berufstätig sind, machen sie oft noch den Haushalt. Auf dem Land arbeiten die Frauen auf dem Feld, deshalb haben sie wenig Zeit für eine Ausbildung. Ich meine, dass die Gleichberechtigung sehr wichtig ist, denn Frauen und Männer müssen die gleichen Chancen haben.}
\end{textebox}

\textbf{Traduction :} « Le texte traite de la situation de la femme en ville et au village. En ville, beaucoup de femmes exercent un métier, parce qu'elles ont une bonne formation scolaire. Bien qu'elles soient actives professionnellement, elles font souvent encore le ménage. Au village, les femmes travaillent aux champs, c'est pourquoi elles ont peu de temps pour une formation. Je pense que l'égalité des droits est très importante, car les femmes et les hommes doivent avoir les mêmes chances. »

\textbf{Analyse du modèle.} Observez comment ce commentaire de 8 lignes applique la méthode :

\begin{center}
\begin{tabularx}{\linewidth}{|X|X|X|}
\hline
\rowcolor{popblueL} Partie & Phrase allemande & Procédé \\
\hline
Einleitung & \all{Der Text handelt von der Situation der Frau...} & formule de présentation \\
\hline
Hauptteil & \all{..., weil sie eine gute Schulbildung haben.} & cause avec \all{weil} + verbe final \\
\hline
Hauptteil & \all{Obwohl sie berufstätig sind, machen sie...} & concession + verbe en 2e position (sujet en 3e) \\
\hline
Hauptteil & \all{..., deshalb haben sie wenig Zeit...} & conséquence + verbe en 2e position \\
\hline
Schluss & \all{Ich meine, dass..., denn...} & avis personnel + \all{denn} (verbe en 2e position) \\
\hline
\end{tabularx}
\end{center}

Remarquez le piège classique : \all{denn} est une \textbf{coordination} (le verbe reste en 2e position : \all{denn Frauen und Männer müssen...}), tandis que \all{weil} est une \textbf{subordination} (verbe en fin : \all{weil sie eine gute Schulbildung haben}). Le français « parce que / car » ne fait pas cette différence : l'allemand, si.

\courssub{Le résumé guidé}

\begin{definition}[Qu'est-ce qu'un résumé (die Zusammenfassung) ?]
Résumer, c'est reprendre \textbf{uniquement les idées principales} du texte, en 3 ou 4 phrases, \textbf{avec ses propres mots}, sans avis personnel et sans recopier de longues phrases du texte.
\end{definition}

\begin{methode}[Réussir un résumé en 4 étapes]
\begin{enumerate}
\item Soulignez dans le texte les 3 ou 4 idées principales (une par paragraphe en général).
\item Reformulez chaque idée avec un synonyme ou une structure différente : par exemple \all{arbeiten auf dem Feld} peut devenir \all{sind in der Landwirtschaft tätig} (« travailler dans l'agriculture »).
\item Reliez les phrases avec des connecteurs variés : \all{zuerst} (d'abord), \all{außerdem} (de plus), \all{deshalb} (c'est pourquoi), \all{obwohl} (bien que).
\item Relisez : pas d'avis personnel dans un résumé, pas de \all{ich meine} !
\end{enumerate}
\end{methode}

\begin{exemplebox}[Résumé guidé du texte support (4 phrases)]
\all{Der Text zeigt die Situation der Frau in Stadt und Land. In der Stadt haben die Frauen oft eine gute Schulbildung und einen Beruf, obwohl sie auch den Haushalt machen. Auf dem Land arbeiten die Frauen hart auf dem Feld, deshalb können sie nicht zur Schule gehen. Die Gleichberechtigung ist also noch ein großes Ziel.}

« Le texte montre la situation de la femme en ville et au village. En ville, les femmes ont souvent une bonne formation et un métier, bien qu'elles fassent aussi le ménage. Au village, les femmes travaillent dur aux champs, c'est pourquoi elles ne peuvent pas aller à l'école. L'égalité des droits reste donc un grand objectif. »
\end{exemplebox}

\courssub{Production écrite guidée : \all{Die Situation der Frau in meinem Land}}

Au baccalauréat, on vous demandera souvent une production d'environ 10 lignes sur un sujet proche du texte. Voici un plan fourni (comme à l'examen) et les phrases utiles pour le remplir.

\begin{popfigure}
\begin{center}
\begin{tikzpicture}[
node distance=0.7cm,
etape/.style={rectangle, rounded corners, draw=poppurple, fill=poppurpleL, thick, align=center, text width=5.6cm, minimum height=0.95cm},
fleche/.style={-{Stealth[length=3mm]}, thick, popdark}]
\node[etape, align=center] (p1){\textbf{1. Situation auf dem Land}\\ \all{Auf dem Land arbeiten die Frauen...}};
\node[etape, below=of p1, fill=popgoldL, draw=popgold, align=center] (p2){\textbf{2. Situation in der Stadt}\\ \all{In der Stadt können die Frauen...}};
\node[etape, below=of p2, fill=poporangeL, draw=poporange, align=center] (p3){\textbf{3. Vergleich}\\ \all{Obwohl es Unterschiede gibt,...}};
\node[etape, below=of p3, fill=popgreenL, draw=popgreen, align=center] (p4){\textbf{4. Meine Meinung}\\ \all{Ich finde es wichtig, dass...}};
\draw[fleche] (p1) -- (p2);
\draw[fleche] (p2) -- (p3);
\draw[fleche] (p3) -- (p4);
\end{tikzpicture}
\end{center}
\legende{Frise du plan de production : village, ville, comparaison, avis personnel.}
\end{popfigure}

\begin{exemplebox}[Phrases utiles pour la production (Redemittel)]
\begin{itemize}
\item \all{Meiner Meinung nach ist die Gleichberechtigung sehr wichtig.} « À mon avis, l'égalité des droits est très importante. »
\item \all{Ich finde es wichtig, dass alle Mädchen zur Schule gehen können.} « Je trouve important que toutes les filles puissent aller à l'école. »
\item \all{In meinem Dorf arbeiten die Frauen auf dem Feld.} « Dans mon village, les femmes travaillent aux champs. »
\item \all{In meiner Stadt haben viele Frauen einen Beruf.} « Dans ma ville, beaucoup de femmes ont un métier. »
\item \all{Obwohl die Frauen viel arbeiten, verdienen sie oft wenig Geld.} « Bien que les femmes travaillent beaucoup, elles gagnent souvent peu d'argent. »
\item \all{Die Regierung muss den Frauen mehr helfen.} « Le gouvernement doit aider davantage les femmes. »
\end{itemize}
\end{exemplebox}

\begin{exoresolu}[Production écrite : \all{Die Situation der Frau in meinem Land}]
Énoncé : Rédigez un texte d'environ 10 lignes en suivant le plan : 1) situation au village ; 2) situation en ville ; 3) comparaison ; 4) votre avis. Utilisez au moins trois connecteurs différents et un verbe de modalité.
\tcblower
Solution détaillée (modèle de réponse) :

\all{In meinem Land, Kamerun, ist die Situation der Frau sehr unterschiedlich. Auf dem Land arbeiten viele Frauen jeden Tag auf dem Feld, weil die Familien das Geld für das Essen brauchen. Die Mädchen müssen oft im Haushalt helfen, deshalb können sie nicht immer zur Schule gehen. In der Stadt, zum Beispiel in Yaoundé oder Douala, ist das Leben anders: viele Frauen haben eine gute Schulbildung und arbeiten als Lehrerin, Ärztin oder Verkäuferin. Obwohl sie einen Beruf haben, kümmern sie sich auch um die Kinder und den Haushalt. Ich finde es wichtig, dass alle Mädchen eine Ausbildung bekommen können, denn Bildung ist der Schlüssel für die Gleichberechtigung.}

Traduction : « Dans mon pays, le Cameroun, la situation de la femme est très différente. Au village, beaucoup de femmes travaillent chaque jour aux champs, parce que les familles ont besoin d'argent pour la nourriture. Les filles doivent souvent aider au ménage, c'est pourquoi elles ne peuvent pas toujours aller à l'école. En ville, par exemple à Yaoundé ou Douala, la vie est différente : beaucoup de femmes ont une bonne formation et travaillent comme institutrice, médecin ou vendeuse. Bien qu'elles aient un métier, elles s'occupent aussi des enfants et du ménage. Je trouve important que toutes les filles puissent recevoir une formation, car l'éducation est la clé de l'égalité des droits. »

Points de méthode : le texte suit le plan en 4 parties ; il contient les connecteurs \all{weil}, \all{deshalb}, \all{obwohl}, \all{denn} ; les modaux \all{müssen} et \all{können} renvoient bien l'infinitif en fin de phrase (\all{helfen}, \all{gehen}, \all{bekommen}).
\end{exoresolu}

\begin{attention}[Check-list avant de rendre votre copie]
Dans toute production écrite, vérifiez systématiquement :
\begin{enumerate}
\item \textbf{Verbe en 2e position} dans la phrase principale : écrivez \all{In der Stadt arbeiten viele Frauen.} et non \all{In der Stadt viele Frauen arbeiten.} (calque du français « En ville, beaucoup de femmes travaillent »).
\item \textbf{Verbe en fin} après \all{weil} et \all{obwohl} : \all{..., weil sie eine gute Schulbildung haben.} et non \all{weil sie haben eine gute Schulbildung}.
\item \textbf{Infinitif final} après un modal : \all{Die Mädchen müssen im Haushalt helfen.} et non \all{müssen helfen im Haushalt}.
\item \textbf{Articles et majuscules} : tous les noms prennent la majuscule (\all{die Frau}, \all{der Beruf}, \all{das Feld}) ; vérifiez le genre (der/die/das) de chaque nom.
\item \textbf{Orthographe} : Umlaute (\all{Mädchen}, \all{können}) et \all{ß}/\all{ss}.
\end{enumerate}
\end{attention}

\courssub{Production orale : présentation en une minute}

\begin{experience}[Mon exposé d'une minute]
Préparez en 5 minutes un exposé oral d'environ une minute sur le sujet : \all{Die Situation der Frau in meinem Dorf / in meiner Stadt} (« La situation de la femme dans mon village / ma ville »).
\begin{enumerate}
\item Notez 4 ou 5 mots-clés seulement (pas de phrases complètes !) : par exemple \all{Feldarbeit — Haushalt — Schule — Beruf — Meinung}.
\item Structurez oralement : \all{Ich spreche über...} (introduction) ; 2 ou 3 phrases avec \all{weil}, \all{deshalb}, \all{obwohl} (développement) ; \all{Meiner Meinung nach...} (conclusion).
\item Parlez lentement, articulez, prononcez le « w » allemand comme un « v » français (\all{die Frau} se prononce « di fraou », \all{weil} se prononce « vail »).
\item La classe écoute et note deux connecteurs entendus ; puis un camarade pose une question : \all{Warum ist die Schulbildung so wichtig?} « Pourquoi la formation scolaire est-elle si importante ? »
\end{enumerate}
\end{experience}

\begin{exemplebox}[Trame d'exposé oral]
\all{Ich spreche über die Situation der Frau in meinem Dorf. Die Frauen arbeiten sehr hart, weil sie auf dem Feld und im Haushalt arbeiten. Obwohl sie viel tun, haben sie oft kein eigenes Geld. Meiner Meinung nach müssen die Mädchen zur Schule gehen, denn Bildung verändert das Leben.}

« Je parle de la situation de la femme dans mon village. Les femmes travaillent très dur, parce qu'elles travaillent aux champs et au ménage. Bien qu'elles fassent beaucoup, elles n'ont souvent pas d'argent propre. À mon avis, les filles doivent aller à l'école, car l'éducation change la vie. »
\end{exemplebox}

\begin{savaistu}[Comment est notée la production écrite au Bac ?]
Au baccalauréat, la production écrite est évaluée selon quatre critères : le \textbf{contenu} (respect du sujet et du plan), le \textbf{vocabulaire} (richesse et exactitude du lexique du thème), la \textbf{grammaire} (ordre des mots, connecteurs, conjugaison) et l'\textbf{orthographe} (majuscules, Umlaute). Retenez ce conseil d'examinateur : une phrase simple mais correcte vaut toujours mieux qu'une phrase longue et fautive. Écrivez court, écrivez juste !
\end{savaistu}

\courssub{Grille d'auto-évaluation}

Après votre production, évaluez-vous honnêtement avec cette grille. Cochez chaque critère si votre texte le respecte.

\begin{center}
\begin{tabularx}{\linewidth}{|X|X|X|}
\hline
\rowcolor{popblueL} Critère & Question de contrôle & Oui / Non \\
\hline
Vocabulaire du thème & Ai-je employé au moins 5 mots du thème (\all{die Frau}, \all{die Gleichberechtigung}, \all{der Haushalt}, \all{der Beruf}, \all{das Feld}) ? & \\
\hline
Ordre des mots & Le verbe conjugué est-il en 2e position dans toutes mes phrases principales ? & \\
\hline
Subordonnées & Le verbe est-il en fin de phrase après \all{weil}, \all{obwohl}, \all{dass} ? & \\
\hline
Connecteurs & Ai-je utilisé au moins 3 connecteurs différents (\all{weil, denn, deshalb, obwohl}) ? & \\
\hline
Conjugaison & Les verbes sont-ils correctement conjugués au présent et l'infinitif est-il final après un modal ? & \\
\hline
Orthographe & Tous les noms ont-ils une majuscule ? Les Umlaute sont-ils présents ? & \\
\hline
Contenu & Ai-je suivi le plan et donné mon avis personnel à la fin ? & \\
\hline
\end{tabularx}
\end{center}

\begin{exercice}[À vous de produire]
\begin{enumerate}
\item Résumez le texte support de la leçon en 4 phrases, avec au moins deux connecteurs différents.
\item Rédigez une production de 10 lignes : \all{Die Situation der Frau in meinem Land}, en suivant la frise du plan (village, ville, comparaison, avis).
\item Préparez un exposé oral d'une minute sur la situation de la femme en ville au Cameroun, avec 5 mots-clés seulement.
\end{enumerate}
\end{exercice}

\begin{corrige}[Exercice — éléments de correction]
\begin{enumerate}
\item Résumé possible : \all{Der Text beschreibt die Situation der Frau in Stadt und Land. In der Stadt haben die Frauen oft einen Beruf, weil sie eine Schulbildung haben. Auf dem Land arbeiten sie auf dem Feld, deshalb fehlt ihnen die Zeit für eine Ausbildung. Obwohl es Fortschritte gibt, ist die Gleichberechtigung noch nicht überall Realität.} (« Le texte décrit la situation de la femme en ville et au village. En ville, les femmes ont souvent un métier parce qu'elles ont une formation. Au village, elles travaillent aux champs, c'est pourquoi le temps leur manque pour une formation. Bien qu'il y ait des progrès, l'égalité des droits n'est pas encore partout une réalité. »)
\item Production : utilisez le modèle de l'exercice résolu comme référence ; vérifiez votre texte avec la grille d'auto-évaluation (verbe en 2e position, verbe final après \all{weil}/\all{obwohl}, infinitif après modal, majuscules).
\item Exposé : la trame \all{Ich spreche über... / ...weil... / Obwohl... / Meiner Meinung nach...} garantit une structure complète en une minute.
\end{enumerate}
\end{corrige}

\begin{aretenir}[L'essentiel de la section]
\begin{itemize}
\item Un commentaire se construit en trois parties : \textbf{Einleitung} (\all{Der Text handelt von...}), \textbf{Hauptteil} (\all{Im Text steht, dass...}), \textbf{Schluss} (\all{Ich meine, dass... weil...}).
\item Un résumé reprend les idées principales en 3-4 phrases, sans avis personnel, avec des connecteurs variés.
\item Une production réussie suit un plan clair et respecte les trois règles d'or : verbe en 2e position, verbe final en subordonnée, infinitif final après un modal.
\item L'auto-évaluation systématique avant de rendre la copie fait gagner des points précieux au baccalauréat.
\end{itemize}
\end{aretenir}

\newpage

\coursec{Activité d'intégration}

\courssub{Objectifs de l'activité}

À la fin de cette activité d'intégration, l'élève sera capable de :

\begin{itemize}
\item réinvestir le \cle{vocabulaire thématique} de la leçon (\all{die Frau, die Gleichberechtigung, die Schulbildung, der Haushalt, der Beruf, die Arbeit auf dem Feld}) dans une production authentique ;
\item employer correctement les quatre \cle{connecteurs logiques} étudiés (\all{weil, denn, deshalb, obwohl}) en respectant la place du verbe ;
\item appliquer la règle fondamentale de l'\cle{ordre des mots} allemand : verbe conjugué en deuxième position dans la phrase principale, verbe conjugué en fin de subordonnée ;
\item conjuguer au présent les \cle{verbes de modalité} (\all{können, müssen, dürfen, wollen, sollen}) avec l'infinitif en fin de phrase ;
\item rédiger en groupe un court article de presse en allemand (10 à 12 lignes) et le présenter oralement en deux minutes ;
\item écouter, comparer et évaluer les productions des autres groupes selon des critères explicites.
\end{itemize}

\courssub{Situation}

Le journal bilingue de ton lycée, \emph{Der Löwe / Le Lion}, prépare son prochain numéro consacré à la société camerounaise d'aujourd'hui. La rédaction lance un appel à contributions aux classes de Terminale. Le thème retenu : la place de la femme camerounaise, en ville et au village. Ta classe de germanistes a décidé de participer. Par groupes de quatre, vous rédigerez un article en allemand ; le meilleur article de la classe sera publié dans le journal et affiché au CDI.

Voici l'appel à contributions tel qu'il paraît dans le journal :

\begin{textebox}[Appel à contributions — Der Löwe, Journal des Gymnasiums]
\textbf{Schreibt mit! Die Frau heute — in der Stadt und auf dem Land}

Liebe Schülerinnen und Schüler der Terminale,

unser Schülerjournal sucht Artikel zum Thema: Wie leben die Frauen in Kamerun heute? In Yaoundé oder Douala arbeiten viele Frauen als Lehrerinnen, Ärztinnen oder Geschäftsfrauen. Sie haben oft eine gute Schulbildung und einen Beruf. Auf dem Land aber müssen viele Frauen früh aufstehen, sie arbeiten auf dem Feld, kochen für die Familie und kümmern sich um die Kinder. Obwohl die Gleichberechtigung in der Verfassung steht, ist der Alltag oft noch ungerecht: Der Haushalt bleibt meistens die Arbeit der Frau.

Schreibt einen kurzen Artikel (10 bis 12 Zeilen): Vergleicht die Situation der Frau in der Stadt und auf dem Land. Was kann die Frau in der Stadt tun? Was muss sie auf dem Land tun? Was wollt ihr ändern?

Eure Redaktion
\end{textebox}

\begin{definition}[Glossaire de l'appel à contributions]
\begin{itemize}
\item \all{die Gleichberechtigung, -en} : l'égalité des droits
\item \all{die Verfassung, -en} : la constitution
\item \all{der Alltag} : la vie quotidienne
\item \all{ungerecht} : injuste
\item \all{sich kümmern um + Akkusativ} : s'occuper de
\item \all{vergleichen (vergleicht, verglich, hat verglichen)} : comparer
\item \all{ändern} : changer
\end{itemize}
\end{definition}

\courssub{Consignes}

\textbf{Phase 1 — Préparation en groupe (20 minutes)}

\begin{enumerate}
\item Lis attentivement l'appel à contributions avec ton groupe. Souligne dans le texte tous les mots du lexique de la leçon que tu reconnais et note-les avec leur article et leur pluriel.
\item Réponds d'abord oralement, en allemand, aux trois questions de la rédaction : \all{Was kann die Frau in der Stadt tun? Was muss sie auf dem Land tun? Was wollt ihr ändern?} Chaque membre du groupe propose au moins une phrase avec un verbe de modalité.
\item Fais un brainstorming : note au brouillon au moins 8 mots du lexique thématique que vous voulez utiliser, et répartissez les rôles (un secrétaire qui écrit, un correcteur grammatical, un chronométreur, un porte-parole).
\end{enumerate}

\textbf{Phase 2 — Rédaction de l'article (25 minutes)}

\begin{enumerate}
\setcounter{enumi}{3}
\item Rédigez un article de 10 à 12 lignes pour \emph{Der Löwe}. Votre article doit obligatoirement contenir :
\begin{itemize}
\item un \cle{titre} accrocheur en allemand ;
\item au moins \cle{8 mots du lexique} de la leçon ;
\item les \cle{4 connecteurs} \all{weil, denn, deshalb, obwohl} (au moins une fois chacun) ;
\item au moins \cle{3 verbes de modalité} différents au présent ;
\item une comparaison claire ville / village ;
\item une conclusion avec un souhait ou une proposition.
\end{itemize}
\item Relisez votre texte avec la grille de vérification : le verbe conjugué est-il en deuxième position dans chaque phrase principale ? Le verbe conjugué est-il à la fin après \all{weil} et \all{obwohl} ? L'infinitif est-il à la fin avec les verbes de modalité ? Tous les noms ont-ils une majuscule ?
\end{enumerate}

\textbf{Phase 3 — Présentation orale et vote (15 minutes)}

\begin{enumerate}
\setcounter{enumi}{5}
\item Le porte-parole de chaque groupe présente l'article en \cle{2 minutes} devant la classe. Les autres élèves écoutent et cochent sur une fiche les connecteurs et les verbes de modalité entendus.
\item La classe vote pour l'article le plus convaincant (clarté, respect des contraintes, richesse du vocabulaire, prononciation).
\item L'enseignant projette ou recopie un article au tableau et le corrige collectivement : la classe repère et justifie chaque position du verbe.
\end{enumerate}

\courssub{Ressources à mobiliser}

\begin{aretenir}[La place du verbe — rappel essentiel]
\begin{itemize}
\item Phrase principale : le verbe conjugué est en \cle{2e position}. \all{Die Frau arbeitet heute auf dem Feld.} « La femme travaille aujourd'hui aux champs. » Si un autre élément ouvre la phrase, le sujet passe après le verbe (inversion) : \all{Heute arbeitet die Frau auf dem Feld.}
\item Subordonnée avec \all{weil} et \all{obwohl} : le verbe conjugué va \cle{à la fin}. \all{Die Frau bleibt zu Hause, weil sie die Kinder betreut.} « La femme reste à la maison parce qu'elle garde les enfants. »
\item \all{denn} : coordination, le verbe reste en 2e position. \all{Sie bleibt zu Hause, denn sie betreut die Kinder.}
\item \all{deshalb} : adverbe de conséquence, il occupe la 1re position et provoque l'inversion. \all{Sie hat keinen Beruf, deshalb bleibt sie zu Hause.}
\item Verbe de modalité : le modal est conjugué en 2e position, l'infinitif va \cle{à la fin}. \all{Die Frauen wollen einen Beruf lernen.} « Les femmes veulent apprendre un métier. »
\end{itemize}
\end{aretenir}

\begin{center}
\begin{tabularx}{\linewidth}{|X|X|X|}\hline
\rowcolor{popblueL} Connecteur & Type & Place du verbe \\ \hline
\all{weil} (parce que) & subordination & verbe à la fin \\ \hline
\all{obwohl} (bien que) & subordination & verbe à la fin \\ \hline
\all{denn} (car) & coordination & verbe en 2e position \\ \hline
\all{deshalb} (c'est pourquoi) & adverbe & inversion : verbe en 2e position, sujet après \\ \hline
\end{tabularx}
\end{center}

\begin{center}
\begin{tabular}{|l|l|l|}\hline
\rowcolor{popblueL} Modal & Présent (ich/sie) & Sens \\ \hline
\all{können} & kann / kann & pouvoir \\ \hline
\all{müssen} & muss / muss & devoir \\ \hline
\all{dürfen} & darf / darf & avoir le droit de \\ \hline
\all{wollen} & will / will & vouloir \\ \hline
\all{sollen} & soll / soll & devoir (recommandation) \\ \hline
\end{tabular}
\end{center}

\begin{attention}[Pièges fréquents des francophones]
\begin{itemize}
\item Ne dis jamais \all{weil sie betreut die Kinder} : après \all{weil}, le verbe va à la fin : \all{weil sie die Kinder betreut}.
\item \all{deshalb} n'est pas une conjonction de subordination : il provoque l'inversion. Écris \all{deshalb arbeitet sie}, jamais \all{deshalb sie arbeitet}.
\item N'oublie pas la majuscule à tous les noms : \all{die Frau, der Beruf, das Feld, die Schule}.
\item Attention au faux ami : \all{der Beruf} signifie « le métier », pas « la berufe » ; « bekommen » signifie « recevoir », pas « devenir » (\all{werden}).
\end{itemize}
\end{attention}

\courssub{Proposition de production et corrigé commenté}

\begin{exoresolu}[Article modèle pour \emph{Der Löwe}]
Rédige un article de 10 à 12 lignes comparant la situation de la femme en ville et au village, avec au moins 8 mots du lexique, les 4 connecteurs et 3 verbes de modalité.
\tcblower
\textbf{Production modèle :}

\all{\textbf{Zwei Welten — die Frau in der Stadt und auf dem Land}}

\all{Die Frau in Kamerun lebt heute in zwei verschiedenen Welten. In der Stadt, zum Beispiel in Yaoundé oder Douala, hat sie oft eine gute Schulbildung und kann einen Beruf lernen: Sie arbeitet als Lehrerin, Ärztin oder Geschäftsfrau. Auf dem Land ist das Leben anders, denn viele Frauen müssen früh aufstehen und den ganzen Tag auf dem Feld arbeiten. Sie kochen, holen Wasser und kümmern sich um die Kinder, weil der Haushalt meistens ihre Arbeit ist. Obwohl die Gleichberechtigung in der Verfassung steht, verdienen viele Frauen auf dem Land kein eigenes Geld. Deshalb wollen viele junge Mädchen in die Stadt gehen und dort eine Schule besuchen. Die Frauen in den Dörfern sollen auch mehr Rechte bekommen, und die Männer müssen im Haushalt helfen. Wir meinen: Die Frau — in der Stadt und auf dem Land — soll dieselben Chancen haben!}

« Deux mondes — la femme en ville et au village. La femme au Cameroun vit aujourd'hui dans deux mondes différents. En ville, par exemple à Yaoundé ou Douala, elle a souvent une bonne scolarité et peut apprendre un métier : elle travaille comme institutrice, médecin ou femme d'affaires. Au village, la vie est différente, car beaucoup de femmes doivent se lever tôt et travailler toute la journée aux champs. Elles cuisinent, vont chercher l'eau et s'occupent des enfants, parce que le ménage est le plus souvent leur travail. Bien que l'égalité des droits soit dans la Constitution, beaucoup de femmes au village ne gagnent pas leur propre argent. C'est pourquoi beaucoup de jeunes filles veulent aller en ville et y fréquenter une école. Les femmes des villages devraient aussi obtenir plus de droits, et les hommes doivent aider au ménage. Nous pensons : la femme — en ville comme au village — doit avoir les mêmes chances ! »

\textbf{Commentaire des choix de langue :}
\begin{itemize}
\item \cle{Lexique} : 10 mots du thème sont réinvestis (\all{die Frau, die Schulbildung, der Beruf, das Feld, der Haushalt, die Gleichberechtigung, die Verfassung, die Schule, die Arbeit, die Stadt}).
\item \cle{Connecteurs} : \all{denn} (verbe en 2e position : \all{denn viele Frauen müssen...}), \all{weil} (verbe final : \all{weil der Haushalt meistens ihre Arbeit ist}), \all{obwohl} (verbe final : \all{Obwohl die Gleichberechtigung in der Verfassung steht...} — et inversion dans la principale qui suit : \all{verdienen viele Frauen}), \all{deshalb} (inversion : \all{Deshalb wollen viele junge Mädchen...}).
\item \cle{Modaux} : \all{kann ... lernen}, \all{müssen ... aufstehen und arbeiten}, \all{sollen ... bekommen}, \all{müssen ... helfen}, \all{soll ... haben} — cinq occurrences, infinitif toujours en fin de phrase.
\item \cle{Ordre des mots} : après \all{zum Beispiel in Yaoundé oder Douala}, le verbe reste en 2e position avec inversion du sujet (\all{hat sie oft...}).
\end{itemize}
\end{exoresolu}

\courssub{Bilan de l'activité}

Cette activité d'intégration a permis de mobiliser simultanément les quatre savoir-faire de la leçon : la compréhension d'un document authentique (l'appel à contributions), l'emploi du vocabulaire thématique de la femme et de la famille, l'application rigoureuse des règles d'ordre des mots avec les connecteurs et les verbes de modalité, et enfin la production écrite et orale dans une situation de communication réaliste. Le vote de la classe et la correction collective au tableau ont développé en outre l'écoute active et l'autoévaluation.

\begin{methode}[Grille d'autoévaluation à retenir]
\begin{enumerate}
\item Mon article a-t-il un titre et une conclusion ?
\item Ai-je utilisé au moins 8 mots du lexique avec le bon article ?
\item Ai-je placé le verbe conjugué en 2e position dans toutes mes phrases principales ?
\item Mes verbes sont-ils à la fin après \all{weil} et \all{obwohl} ?
\item Ai-je fait l'inversion après \all{deshalb} ?
\item Mes infinitifs sont-ils bien à la fin avec les verbes de modalité ?
\item Tous mes noms commencent-ils par une majuscule ?
\end{enumerate}
\end{methode}

Pour aller plus loin, l'article gagnant pourra être illustré, tapé à l'ordinateur et envoyé à la rédaction de \emph{Der Löwe}, puis comparé avec un article similaire sur la situation de la femme en Allemagne, en Autriche ou en Suisse — une ouverture interculturelle qui prépare l'épreuve d'expression écrite du baccalauréat.

\newpage

\coursec{Méthodes et exercices résolus}

\courssub{Méthode 1 : Comprendre un texte (global puis détaillé)}

\begin{methode}[Lire et comprendre un texte sur la situation de la femme]
\begin{enumerate}
\item \textbf{Lecture globale (Hauptidee).} Lis le texte une première fois sans dictionnaire. Repère : le thème (ici : \all{die Situation der Frau}), le lieu (\all{in der Stadt / auf dem Land}), le type de texte (article, témoignage, lettre) et l'idée principale.
\item \textbf{Repérage des mots-clés.} Souligne le vocabulaire du thème : \all{die Gleichberechtigung}, \all{die Schulbildung}, \all{der Haushalt}, \all{der Beruf}, \all{die Arbeit auf dem Feld}. Ces mots t'aident à structurer ta compréhension.
\item \textbf{Lecture détaillée.} Relis paragraphe par paragraphe. Pour chaque paragraphe, formule une phrase de résumé en allemand simple : \all{Im ersten Absatz geht es um ...} (« Dans le premier paragraphe, il s'agit de ... »).
\item \textbf{Exploitation des connecteurs.} Les connecteurs (\all{weil}, \all{denn}, \all{deshalb}, \all{obwohl}) signalent les causes, les conséquences et les oppositions : ils donnent les arguments du texte. Note-les.
\item \textbf{Réponse aux questions.} Réponds toujours par une \textbf{phrase complète} en allemand, en reprenant les mots de la question, jamais par un seul mot. Exemple : \all{Warum arbeiten viele Frauen auf dem Feld?} $\rightarrow$ \all{Viele Frauen arbeiten auf dem Feld, weil sie ihre Familien ernähren müssen.}
\end{enumerate}
\end{methode}

\begin{exoresolu}[Compréhension de texte — type Bac]
\textbf{Texte.} \all{In vielen Dörfern Kameruns arbeiten die Frauen jeden Tag auf dem Feld, weil sie ihre Kinder ernähren müssen. Obwohl sie sehr hart arbeiten, bekommen sie oft wenig Geld. In den Städten haben die Frauen dagegen mehr Möglichkeiten: Sie können zur Schule gehen und einen Beruf lernen. Trotzdem bleibt die Gleichberechtigung ein Ziel für die Zukunft.}

\textbf{Questions.} 1) \all{Wo arbeiten die Frauen auf dem Land?} 2) \all{Warum bekommen die Frauen im Dorf oft wenig Geld?} (question piège : le texte dit-il pourquoi ?) 3) \all{Welche Möglichkeiten haben die Frauen in der Stadt?} 4) \all{Was bleibt ein Ziel für die Zukunft?}
\tcblower
\textbf{Raisonnement.} On localise chaque information dans le texte, puis on rédige une phrase complète avec le verbe en 2e position.

1) L'information est dans la première phrase. \all{\textbf{Die Frauen arbeiten auf dem Land jeden Tag auf dem Feld.}} (« Les femmes travaillent chaque jour aux champs. ») — Verbe \all{arbeiten} en 2e position, complément de temps \all{jeden Tag} avant le complément de lieu \all{auf dem Feld}.

2) Attention : le texte dit \emph{que} les femmes gagnent peu (\all{obwohl sie sehr hart arbeiten}), mais il ne dit pas \emph{pourquoi}. La bonne réponse honnête est : \all{\textbf{Das steht nicht im Text. Der Text sagt nur, dass sie wenig Geld bekommen, obwohl sie hart arbeiten.}} (« Cela ne figure pas dans le texte. Le texte dit seulement qu'elles gagnent peu d'argent, bien qu'elles travaillent dur. ») — Au Bac, ne jamais inventer une information absente du texte.

3) \all{\textbf{In der Stadt können die Frauen zur Schule gehen und einen Beruf lernen.}} (« En ville, les femmes peuvent aller à l'école et apprendre un métier. ») — Inversion sujet-verbe après \all{In der Stadt} (le verbe conjugué \all{können} reste en 2e position, les infinitifs \all{gehen} et \all{lernen} vont en fin de phrase).

4) \all{\textbf{Die Gleichberechtigung bleibt ein Ziel für die Zukunft.}} (« L'égalité des droits reste un objectif pour l'avenir. »)

\textbf{Conclusion.} Une bonne réponse de compréhension = phrase complète + information exacte du texte + grammaire correcte (place du verbe, majuscules aux noms).
\end{exoresolu}

\courssub{Méthode 2 : Employer le vocabulaire thématique}

\begin{methode}[Apprendre et utiliser le lexique de la femme, de la famille et du travail]
\begin{enumerate}
\item \textbf{Apprends chaque nom avec son article et son pluriel} : \all{die Frau, -en} ; \all{die Gleichberechtigung, -en} ; \all{die Schulbildung} (surtout au singulier) ; \all{der Haushalt, -e} ; \all{der Beruf, -e} ; \all{die Arbeit, -en} ; \all{das Feld, -er} ; \all{das Dorf, -er} ; \all{die Stadt, -e} ; \all{die Familie, -n} ; \all{das Kind, -er} ; \all{das Recht, -e}.
\item \textbf{Apprends les familles de mots} : \all{arbeiten} (v.) $\rightarrow$ \all{die Arbeit} (n.) $\rightarrow$ \all{der Arbeiter / die Arbeiterin} (personne) ; \all{bilden} $\rightarrow$ \all{die Bildung} ; \all{gleich} (adj.) $\rightarrow$ \all{die Gleichberechtigung} ; \all{die Schule} $\rightarrow$ \all{die Schulbildung}.
\item \textbf{Mémorise les verbes du thème avec leur construction} : \all{arbeiten (auf dem Feld)} « travailler (aux champs) », \all{Geld verdienen} « gagner de l'argent », \all{sich kümmern um + Akk.} « s'occuper de », \all{die Kinder erziehen} « élever les enfants », \all{einen Beruf lernen} « apprendre un métier », \all{studieren} « faire des études (universitaires) ».
\item \textbf{Réutilise le lexique dans tes propres phrases} : un mot n'est acquis que s'il est employé. Construis des phrases sur le Cameroun : \all{Die Frauen auf dem Markt in Douala verkaufen Obst und Gemüse.} (« Les femmes au marché de Douala vendent des fruits et des légumes. »)
\item \textbf{Attention aux genres} : note toujours \all{der/die/das} ; au Bac, une faute d'article est une faute comptée.
\end{enumerate}
\end{methode}

\begin{exoresolu}[Vocabulaire — type Bac]
\textbf{Énoncé.} Complète les phrases avec le mot qui convient (avec le bon article et la bonne forme) : \all{Gleichberechtigung — Beruf — Haushalt — Schulbildung — Feld}.

1) \all{Meine Mutter arbeitet jeden Tag auf dem \_\_\_\_\_.} 2) \all{Die \_\_\_\_\_ zwischen Mann und Frau ist sehr wichtig.} 3) \all{Ohne \_\_\_\_\_ kann man keinen guten Beruf finden.} 4) \all{Mein Vater hilft oft im \_\_\_\_\_: Er kocht und wäscht.} 5) \all{Sie hat einen guten \_\_\_\_\_ gelernt: Sie ist Lehrerin.}
\tcblower
\textbf{Raisonnement.} On identifie le sens attendu, puis on vérifie l'article et le cas exigé par la préposition ou la fonction.

1) \all{auf dem \textbf{Feld}} — \all{das Feld} ; \all{auf} + datif (lieu où l'on est) : \all{auf dem Feld}. (« Ma mère travaille chaque jour aux champs. »)

2) \all{Die \textbf{Gleichberechtigung}} — \all{die Gleichberechtigung}, nominatif sujet. (« L'égalité entre l'homme et la femme est très importante. »)

3) \all{Ohne \textbf{Schulbildung}} — \all{die Schulbildung} ; \all{ohne} + accusatif ; le nom féminin garde la même forme à l'accusatif. (« Sans scolarisation, on ne peut pas trouver un bon métier. »)

4) \all{im \textbf{Haushalt}} — \all{der Haushalt} ; \all{in} + datif : \all{im Haushalt} (= \all{in dem}). (« Mon père aide souvent dans les tâches ménagères : il cuisine et fait la lessive. »)

5) \all{einen guten \textbf{Beruf}} — \all{der Beruf} ; accusatif masculin : \all{einen guten Beruf} (adjectif en \all{-en} après \all{einen}). (« Elle a appris un bon métier : elle est institutrice. »)

\textbf{Conclusion.} Toujours vérifier trois choses : le sens, l'article, le cas exigé par la préposition ou la fonction.
\end{exoresolu}

\courssub{Méthode 3 : Maîtriser l'ordre des mots}

\begin{methode}[La place du verbe : 2e position et subordonnées]
\begin{enumerate}
\item \textbf{Phrase principale : le verbe conjugué est TOUJOURS en 2e position.} La 1re position peut être le sujet (\all{Die Frau arbeitet.}) ou un complément (\all{Jeden Tag arbeitet die Frau.}) — dans ce cas, le sujet passe après le verbe (inversion).
\item \textbf{Subordonnée introduite par \all{weil}, \all{obwohl}, \all{dass} : le verbe conjugué va à la FIN.} \all{Die Frau arbeitet hart, weil sie ihre Kinder ernähren muss.} (« La femme travaille dur parce qu'elle doit nourrir ses enfants. »)
\item \textbf{\all{denn} ne change pas l'ordre} : c'est une coordination, le verbe reste en 2e position. \all{Sie arbeitet hart, denn sie muss ihre Kinder ernähren.} (« Elle travaille dur, car elle doit nourrir ses enfants. »)
\item \textbf{\all{deshalb} occupe la 1re position} et provoque l'inversion : \all{Sie ist müde, deshalb geht sie früh ins Bett.} (« Elle est fatiguée, c'est pourquoi elle se couche tôt. »)
\item \textbf{Réflexe de vérification} : après avoir écrit une phrase, compte les positions et souligne le verbe conjugué. Une seule question : « Mon verbe est-il en 2e position (principale) ou en fin (subordonnée) ? »
\end{enumerate}
\end{methode}

\begin{exoresolu}[Ordre des mots — transformation type Bac]
\textbf{Énoncé.} a) Mets les phrases dans le bon ordre. b) Relie les deux phrases avec \all{weil}.

a) 1) \all{arbeitet / auf dem Feld / die Frau / jeden Tag} 2) \all{in der Stadt / mehr Möglichkeiten / haben / die Frauen}

b) \all{Die Frauen im Dorf verdienen wenig Geld. + Sie bekommen keine Schulbildung.}
\tcblower
\textbf{Raisonnement.} a) On place d'abord le verbe conjugué en 2e position, puis on répartit sujet et compléments.

a) 1) \all{\textbf{Die Frau arbeitet jeden Tag auf dem Feld.}} — Ordre : sujet (1) + verbe (2) + complément de temps + complément de lieu. (« La femme travaille chaque jour aux champs. »)

a) 2) \all{\textbf{In der Stadt haben die Frauen mehr Möglichkeiten.}} — Le complément de lieu occupe la 1re position ; le verbe \all{haben} reste en 2e position, donc inversion du sujet. (« En ville, les femmes ont plus de possibilités. »)

b) On choisit la cause : pourquoi gagnent-elles peu ? Parce qu'elles n'ont pas de scolarisation. \all{weil} envoie le verbe conjugué à la fin : \all{\textbf{Die Frauen im Dorf verdienen wenig Geld, weil sie keine Schulbildung bekommen.}} (« Les femmes du village gagnent peu d'argent parce qu'elles ne reçoivent pas d'instruction. ») — Vérification : \all{bekommen} bien en dernière position de la subordonnée, virgule avant \all{weil}.

\textbf{Conclusion.} La place du verbe est le critère n° 1 de correction grammaticale au Bac : une phrase correcte au sens mais avec un verbe mal placé perd des points.
\end{exoresolu}

\courssub{Méthode 4 : Employer les connecteurs weil, denn, deshalb, obwohl}

\begin{methode}[Choisir et placer le bon connecteur]
\begin{enumerate}
\item \textbf{Identifier le rapport logique} : cause (\all{weil}, \all{denn}), conséquence (\all{deshalb}), opposition/concession (\all{obwohl}).
\item \textbf{\all{weil} + verbe final} : \all{Sie bleibt zu Hause, weil sie krank ist.} (« Elle reste à la maison parce qu'elle est malade. »)
\item \textbf{\all{denn} + verbe en 2e position} : \all{Sie bleibt zu Hause, denn sie ist krank.} — \all{denn} est plus soutenu, fréquent à l'écrit ; il ne se place jamais en début de phrase.
\item \textbf{\all{deshalb} + inversion} : \all{Sie ist krank, deshalb bleibt sie zu Hause.} (« Elle est malade, c'est pourquoi elle reste à la maison. ») — Le verbe conjugué vient juste après \all{deshalb}.
\item \textbf{\all{obwohl} + verbe final} : \all{Obwohl sie krank ist, geht sie zur Arbeit.} (« Bien qu'elle soit malade, elle va au travail. ») — Si la subordonnée est en tête, la principale commence par le verbe : \all{Obwohl..., \textbf{geht} sie...}.
\item \textbf{Traductions-pièges} : « parce que » = \all{weil/denn} ; « donc, c'est pourquoi » = \all{deshalb} ; « bien que » = \all{obwohl} (jamais un calque de « malgré que »).
\end{enumerate}
\end{methode}

\begin{center}
\begin{tabularx}{\linewidth}{|l|X|X|X|}
\hline
\rowcolor{popblueL} \textbf{Connecteur} & \textbf{Sens} & \textbf{Place du verbe} & \textbf{Exemple} \\
\hline
\all{weil} & parce que (cause) & verbe en fin & \all{Sie arbeitet, weil sie Geld braucht.} \\
\hline
\all{denn} & car (cause) & verbe en 2e position & \all{Sie arbeitet, denn sie braucht Geld.} \\
\hline
\all{deshalb} & c'est pourquoi (conséquence) & inversion & \all{Sie braucht Geld, deshalb arbeitet sie.} \\
\hline
\all{obwohl} & bien que (concession) & verbe en fin & \all{Sie arbeitet, obwohl sie müde ist.} \\
\hline
\end{tabularx}
\end{center}

\begin{exoresolu}[Connecteurs — type Bac]
\textbf{Énoncé.} Complète avec \all{weil}, \all{denn}, \all{deshalb} ou \all{obwohl}, puis traduis en français.

1) \all{Die Frauen arbeiten hart, \_\_\_\_\_ sie wenig Geld bekommen.} 2) \all{Viele Mädchen gehen nicht zur Schule, \_\_\_\_\_ ihre Eltern kein Geld haben.} 3) \all{Die Frau hat keinen Beruf gelernt, \_\_\_\_\_ kann sie keinen guten Job finden.} 4) \all{\_\_\_\_\_ die Frauen in der Stadt mehr Rechte haben, bleibt die Gleichberechtigung ein Problem.}
\tcblower
\textbf{Raisonnement.} On détermine le rapport logique, puis on vérifie la place du verbe.

1) Opposition (elles travaillent dur \emph{mais} gagnent peu) : \all{\textbf{obwohl}}. \all{Die Frauen arbeiten hart, obwohl sie wenig Geld bekommen.} — Verbe \all{bekommen} en fin. « Les femmes travaillent dur, bien qu'elles gagnent peu d'argent. »

2) Cause, et le verbe \all{haben} est en fin de phrase dans la phrase donnée $\rightarrow$ subordination : \all{\textbf{weil}}. \all{Viele Mädchen gehen nicht zur Schule, weil ihre Eltern kein Geld haben.} « Beaucoup de filles ne vont pas à l'école parce que leurs parents n'ont pas d'argent. » (Avec \all{denn}, il aurait fallu : \all{..., denn ihre Eltern haben kein Geld.} — verbe en 2e position.)

3) Conséquence : \all{\textbf{deshalb}}. \all{Die Frau hat keinen Beruf gelernt, deshalb kann sie keinen guten Job finden.} — Inversion : le verbe \all{kann} précède le sujet \all{sie}. « La femme n'a pas appris de métier, c'est pourquoi elle ne peut pas trouver un bon emploi. »

4) Concession en tête de phrase : \all{\textbf{Obwohl}}. \all{Obwohl die Frauen in der Stadt mehr Rechte haben, bleibt die Gleichberechtigung ein Problem.} — Subordonnée en tête $\rightarrow$ la principale commence par le verbe \all{bleibt}. « Bien que les femmes en ville aient plus de droits, l'égalité reste un problème. »

\textbf{Conclusion.} Connecteur choisi = sens juste + place du verbe juste. Les deux sont indissociables.
\end{exoresolu}

\courssub{Méthode 5 : Conjuguer au présent et employer les verbes de modalité}

\begin{methode}[Présent et verbes de modalité]
\begin{enumerate}
\item \textbf{Verbes faibles} : radical + \all{-e, -st, -t, -en, -t, -en}. \all{arbeiten} : \all{ich arbeite, du arbeitest, er arbeitet, wir arbeiten, ihr arbeitet, sie arbeiten} (attention : \all{-est/-et} après un radical en \all{-t}).
\item \textbf{Verbes forts} : changement de voyelle à la 2e et 3e personne du singulier : \all{geben $\rightarrow$ du gibst, er gibt} ; \all{lesen $\rightarrow$ du liest, er liest} ; \all{fahren $\rightarrow$ du fährst, er fährt}. Le pluriel garde le radical de l'infinitif.
\item \textbf{Verbes de modalité} : \all{können} (pouvoir), \all{müssen} (devoir), \all{wollen} (vouloir), \all{dürfen} (avoir le droit), \all{sollen} (devoir — recommandation), \all{mögen/möchte} (aimer / vouloir poliment).
\item \textbf{Construction} : modal conjugué en 2e position + infinitif à la FIN : \all{Die Frau \textbf{muss} jeden Tag auf dem Feld \textbf{arbeiten}.} Dans une subordonnée, le modal conjugué va à la fin : \all{..., weil sie arbeiten \textbf{muss}.}
\item \textbf{Singulier irrégulier des modaux} : la 1re et la 3e personne du singulier sont IDENTIQUES (pas de \all{-t} à \all{er}) : \all{ich kann, er kann} ; \all{ich muss, er muss} ; \all{ich will, er will} ; \all{ich darf, er darf} ; \all{ich soll, er soll}.
\end{enumerate}
\end{methode}

\begin{propriete}[Conjugaison des verbes de modalité au présent]
\begin{center}
\begin{tabular}{|l|l|l|l|l|}
\hline
\rowcolor{popblueL} & \textbf{können} & \textbf{müssen} & \textbf{wollen} & \textbf{dürfen} \\
\hline
\all{ich} & \all{kann} & \all{muss} & \all{will} & \all{darf} \\
\hline
\all{du} & \all{kannst} & \all{musst} & \all{willst} & \all{darfst} \\
\hline
\all{er/sie/es} & \all{kann} & \all{muss} & \all{will} & \all{darf} \\
\hline
\all{wir} & \all{können} & \all{müssen} & \all{wollen} & \all{dürfen} \\
\hline
\all{ihr} & \all{könnt} & \all{müsst} & \all{wollt} & \all{dürft} \\
\hline
\all{sie/Sie} & \all{können} & \all{müssen} & \all{wollen} & \all{dürfen} \\
\hline
\end{tabular}
\end{center}
Remarque : au singulier, la voyelle change (\all{ö $\rightarrow$ a}, \all{ü $\rightarrow$ u}) et il n'y a pas de \all{-t} à la 3e personne : \all{er kann} (jamais \all{*er kannt}).
\end{propriete}

\begin{exoresolu}[Conjugaison — thème type Bac]
\textbf{Énoncé.} Traduis en allemand : 1) « Les femmes doivent travailler dur. » 2) « Elle lit un livre parce qu'elle veut apprendre. » 3) « Les filles ne peuvent pas aller à l'école. » 4) « Mon père donne de l'argent à ma mère. »
\tcblower
\textbf{Raisonnement.} On identifie le verbe (faible, fort, modal), on conjugue, on place l'infinitif à la fin s'il y a un modal.

1) \all{müssen} (modal) + \all{arbeiten} (infinitif final) : \all{\textbf{Die Frauen müssen hart arbeiten.}} — \all{müssen} au pluriel = forme de l'infinitif. (« Les femmes doivent travailler dur. »)

2) \all{lesen} fort : \all{er/sie liest} ; subordonnée avec \all{wollen} : le modal conjugué va à la fin. \all{\textbf{Sie liest ein Buch, weil sie lernen will.}} — Attention : \all{will} (et non \all{*willt}) à la 3e personne du singulier. (« Elle lit un livre parce qu'elle veut apprendre. »)

3) \all{können} + négation : \all{\textbf{Die Mädchen können nicht zur Schule gehen.}} — \all{gehen} infinitif en fin de phrase, \all{nicht} avant le complément de direction. (« Les filles ne peuvent pas aller à l'école. »)

4) \all{geben} fort : \all{er gibt} ; la personne qui reçoit est au datif : \all{meiner Mutter}. \all{\textbf{Mein Vater gibt meiner Mutter Geld.}} (« Mon père donne de l'argent à ma mère. ») — Piège : \all{die Mutter} $\rightarrow$ datif \all{der Mutter} ; avec le possessif : \all{meiner Mutter}.

\textbf{Conclusion.} Trois réflexes : voyelle du verbe fort au singulier, infinitif en fin avec un modal, cas correct (datif avec \all{geben}).
\end{exoresolu}

\courssub{Méthode 6 : Rédiger une courte production guidée}

\begin{methode}[Résumer et produire un court paragraphe]
\begin{enumerate}
\item \textbf{Résumé (Zusammenfassung)} : reprends l'idée principale + 2 ou 3 idées secondaires, avec TES mots, au présent, en 3 à 5 phrases. Commence par : \all{Der Text handelt von ...} (« Le texte traite de ... ») ou \all{Im Text geht es um ...} (« Dans le texte, il s'agit de ... »).
\item \textbf{Plan d'un court paragraphe d'opinion} : phrase d'introduction (le thème) $\rightarrow$ 2 arguments avec connecteurs (\all{weil}, \all{denn}, \all{deshalb}, \all{obwohl}) $\rightarrow$ exemple (ville/village, Cameroun) $\rightarrow$ conclusion personnelle : \all{Ich meine, dass ...} / \all{Meiner Meinung nach ...} (« À mon avis ... »).
\item \textbf{Structures à retenir} : \all{Die Frauen spielen eine wichtige Rolle in der Familie.} (« Les femmes jouent un rôle important dans la famille. ») ; \all{Die Gleichberechtigung ist wichtig für die Zukunft.} (« L'égalité des droits est importante pour l'avenir. ») ; \all{Man muss den Mädchen eine gute Schulbildung geben.} (« Il faut donner aux filles une bonne scolarisation. »)
\item \textbf{Relecture obligatoire} : verbe en 2e position ? verbe final dans les subordonnées ? majuscules aux noms ? articles corrects ?
\end{enumerate}
\end{methode}

\begin{exoresolu}[Production écrite — type Bac]
\textbf{Énoncé.} \all{Schreiben Sie einen kurzen Text (5--6 Sätze): „Die Situation der Frau in meinem Land“. Benutzen Sie mindestens zwei Konnektoren (weil, denn, deshalb, obwohl) und ein Modalverb.} (« Rédigez un court texte (5-6 phrases) : “La situation de la femme dans mon pays”. Utilisez au moins deux connecteurs et un verbe de modalité. »)
\tcblower
\textbf{Raisonnement.} Plan : introduction (rôle des femmes) $\rightarrow$ argument 1 (travail au village, \all{weil}) $\rightarrow$ argument 2 (concession, \all{obwohl}) $\rightarrow$ conséquence (\all{deshalb}) $\rightarrow$ conclusion avec modal (\all{müssen} / \all{man muss}).

\textbf{Production modèle :}

\all{\textbf{Die Frauen spielen in Kamerun eine sehr wichtige Rolle. Auf dem Land arbeiten sie jeden Tag auf dem Feld, weil sie ihre Familien ernähren müssen. Obwohl sie sehr hart arbeiten, haben sie oft wenig Geld und wenig Zeit. In der Stadt haben die Frauen mehr Möglichkeiten, deshalb können viele Mädchen dort zur Schule gehen. Ich meine, dass die Gleichberechtigung sehr wichtig ist. Man muss den Frauen mehr Rechte und eine gute Schulbildung geben.}}

Traduction : « Les femmes jouent un rôle très important au Cameroun. À la campagne, elles travaillent chaque jour aux champs parce qu'elles doivent nourrir leurs familles. Bien qu'elles travaillent très dur, elles ont souvent peu d'argent et peu de temps. En ville, les femmes ont plus de possibilités, c'est pourquoi beaucoup de filles peuvent y aller à l'école. Je pense que l'égalité des droits est très importante. Il faut donner aux femmes plus de droits et une bonne instruction. »

\textbf{Vérification.} \all{weil} + verbe final (\all{ernähren müssen}) ; \all{obwohl} en tête + principale avec verbe en 1re position (\all{haben sie}) ; \all{deshalb} + inversion (\all{können viele Mädchen}) ; modal \all{muss} + infinitif \all{geben} en fin ; majuscules à tous les noms (\all{Rolle, Land, Möglichkeiten, Gleichberechtigung, Rechte, Schulbildung}).

\textbf{Conclusion.} Une production réussie = plan simple + connecteurs variés et bien placés + relecture systématique.
\end{exoresolu}

\begin{attention}[Les 5 erreurs qui coûtent des points au Bac]
\begin{enumerate}
\item \textbf{Verbe mal placé.} Écrire \all{*Weil sie ihre Kinder ernähren muss, die Frau arbeitet hart} ou oublier l'inversion après un complément en tête (\all{*In der Stadt die Frauen haben...}). Règle d'or : verbe conjugué en 2e position dans la principale, en FIN dans la subordonnée. Correct : \all{Weil sie ihre Kinder ernähren muss, arbeitet die Frau hart.} et \all{In der Stadt haben die Frauen mehr Möglichkeiten.}
\item \textbf{Oubli de la majuscule aux noms.} \all{*die frau, *das dorf, *die gleichberechtigung} sont FAUX : tout nom allemand prend une majuscule : \all{die Frau, das Dorf, die Gleichberechtigung}.
\item \textbf{Confusion des connecteurs.} \all{denn} ne renvoie PAS le verbe à la fin (\all{*denn sie kein Geld hat} est faux : \all{denn sie hat kein Geld}) ; \all{deshalb} exige l'inversion (\all{deshalb kann sie...}, jamais \all{*deshalb sie kann...}).
\item \textbf{Faux amis et calques du français.} \all{bekommen} = « recevoir » et non « devenir » (\all{werden}) ; « apprendre un métier » = \all{einen Beruf lernen} ; « s'occuper de » = \all{sich kümmern um} ; ne traduis jamais mot à mot « la situation de la femme » par \all{*die Situation von der Frau} : dis \all{die Situation der Frau} (génitif).
\item \textbf{Articles et cas négligés.} \all{der Haushalt} (masculin !), \all{das Feld} (neutre : \all{auf dem Feld}), \all{die Stadt} (féminin : \all{in der Stadt}). Après \all{geben}, la personne est au datif : \all{Er gibt der Frau Geld.} (« Il donne de l'argent à la femme. ») Apprends chaque nom avec son article et chaque verbe avec sa construction.
\end{enumerate}
\end{attention}

\newpage

\coursec{Exercices}

\courssub{Je vérifie mes connaissances}

\begin{exercice}[QCM : la grammaire de la leçon]
Entoure la bonne réponse. Une seule réponse est correcte.
\begin{enumerate}
\item \all{Die Frau arbeitet viel, \dots\ sie auch den Haushalt macht.}\\
a) \all{weil} \quad b) \all{denn} \quad c) \all{obwohl} \quad d) \all{deshalb}
\item \all{Mama bleibt zu Hause, \dots\ sie krank \dots\ .} (le verbe conjugué doit être en fin de phrase)\\
a) \all{denn} \quad b) \all{weil} \quad c) \all{deshalb} \quad d) \all{und}
\item \all{Sie hat kein Geld, \dots\ kann sie nicht auf den Markt gehen.}\\
a) \all{weil} \quad b) \all{obwohl} \quad c) \all{deshalb} \quad d) \all{denn}
\item Dans la phrase \all{Obwohl sie müde ist, arbeitet sie auf dem Feld.}, le verbe conjugué de la subordonnée se trouve :\\
a) en deuxième position \quad b) en troisième position \quad c) en fin de phrase \quad d) au début de la phrase
\item Le pluriel de \all{die Frau} est :\\
a) \all{die Fraus} \quad b) \all{die Frauen} \quad c) \all{die Fraue} \quad d) \all{die Fräulein}
\end{enumerate}
\end{exercice}

\begin{exercice}[Vrai ou faux ? Justifie ta réponse]
Réponds par \all{richtig} ou \all{falsch} et justifie en une phrase en français.
\begin{enumerate}
\item \all{Denn} est une conjonction de subordination : le verbe conjugué va à la fin de la phrase.
\item \all{Deshalb} est un adverbe de liaison : le verbe reste en deuxième position et le sujet se place après le verbe.
\item \all{Obwohl} exprime une cause.
\item Au présent, \all{können} se conjugue : \all{ich kann, du kannst, er kann}.
\item Le verbe fort \all{lesen} fait \all{er lest} à la troisième personne du singulier.
\end{enumerate}
\end{exercice}

\begin{exercice}[Vocabulaire : la femme, la famille, le travail]
Complète chaque phrase avec le mot qui convient, choisi dans la liste : \all{die Gleichberechtigung -- der Haushalt -- die Schulbildung -- das Feld -- der Beruf}.
\begin{enumerate}
\item Meine Mutter arbeitet jeden Tag auf dem \dots\ ; sie pflanzt Maniok und Mais.
\item Viele Mädchen im Dorf haben keine \dots\ , weil ihre Eltern kein Geld für die Schule haben.
\item Meine Schwester hat einen guten \dots\ : sie ist Lehrerin in Yaoundé.
\item Vater und Mutter teilen jetzt den \dots\ : er kocht, sie wäscht.
\item \dots\ bedeutet: Männer und Frauen haben die gleichen Rechte.
\end{enumerate}
\end{exercice}

\begin{exercice}[Conjugaison : le présent]
Conjugue les verbes entre parenthèses au présent.
\begin{enumerate}
\item Amina (arbeiten) \dots\ jeden Tag von 6 bis 17 Uhr.
\item Sie (nehmen) \dots\ das Gemüse und (tragen) \dots\ es zum Markt.
\item Ihre Tochter (helfen) \dots\ ihr am Wochenende.
\item Amina (können) \dots\ nicht lesen, aber ihre Tochter (wollen) \dots\ Ärztin werden.
\item Die Frauen im Dorf (müssen) \dots\ viel arbeiten, aber sie (dürfen) \dots\ nicht über das Geld entscheiden.
\item Mein Vater (sprechen) \dots\ oft über die Gleichberechtigung.
\end{enumerate}
\end{exercice}

\courssub{Je m'entraîne}

\begin{exercice}[Remets les mots dans l'ordre]
Forme des phrases correctes. Attention à la place du verbe (2e position dans la principale, fin de phrase dans la subordonnée).
\begin{enumerate}
\item arbeitet / die Frau / jeden Tag / auf dem Feld
\item sie / weil / müde / ist / geht / früh / ins Bett / sie
\item können / viele Mädchen / nicht / zur Schule / gehen / weil / ihre Eltern / arm / sind
\item obwohl / viel / arbeitet / sie / verdient / wenig Geld / sie
\item deshalb / will / meine Schwester / studieren / in Deutschland
\end{enumerate}
\end{exercice}

\begin{exercice}[Relie les phrases avec le bon connecteur]
Relie les deux phrases avec \all{weil}, \all{denn}, \all{deshalb} ou \all{obwohl} selon la relation logique indiquée entre parenthèses. Attention à la place du verbe.
\begin{enumerate}
\item Fatou lernt Deutsch. Sie will in Deutschland studieren. (cause, verbe à la fin)
\item Die Frau arbeitet den ganzen Tag. Sie ist nie müde. (opposition)
\item Es regnet stark. Die Frauen gehen nicht auf das Feld. (conséquence)
\item Meine Mutter kann nicht lesen. Sie ist nie zur Schule gegangen. (cause, verbe en 2e position)
\item Die Mädchen haben wenig Zeit. Sie müssen im Haushalt helfen. (cause, verbe à la fin)
\end{enumerate}
\end{exercice}

\begin{exercice}[Texte à trous : la vie de Maman Rose]
Complète le texte avec les verbes au présent : \all{aufstehen -- kochen -- verkaufen -- können -- müssen -- helfen -- verdienen -- wollen}.

\begin{textebox}[Maman Rose, eine Frau aus dem Dorf]
Maman Rose \dots\ (1) jeden Morgen um fünf Uhr \dots\ . Zuerst \dots\ (2) sie das Frühstück für die Familie. Dann geht sie auf den Markt und \dots\ (3) Gemüse und Früchte. Sie \dots\ (4) nicht viel Geld, aber sie \dots\ (5) ihre Kinder ernähren. Ihre Tochter Sandrine \dots\ (6) ihr am Wochenende. Maman Rose \dots\ (7) nicht lesen und schreiben, aber Sandrine \dots\ (8) später Lehrerin werden.
\end{textebox}
\end{exercice}

\begin{exercice}[Appariements : mots et définitions]
Relie chaque mot allemand à sa définition française.

\begin{center}
\begin{tabularx}{\linewidth}{|X|X|}
\hline
\rowcolor{popblueL} \textbf{Deutsch} & \textbf{Français} \\
\hline
1. \all{die Gleichberechtigung} & a) le travail des champs \\
\hline
2. \all{der Haushalt} & b) l'égalité des droits entre hommes et femmes \\
\hline
3. \all{die Arbeit auf dem Feld} & c) la scolarisation \\
\hline
4. \all{die Schulbildung} & d) les tâches ménagères \\
\hline
5. \all{der Beruf} & e) le métier, la profession \\
\hline
6. \all{die Ehe} & f) le mariage \\
\hline
\end{tabularx}
\end{center}
\end{exercice}

\courssub{Je me prépare au Bac}

\begin{exercice}[Compréhension écrite type Bac]
Lis le texte suivant, puis réponds aux questions.

\begin{textebox}[Eine starke Frau aus Kamerun]
Aminatou lebt in einem Dorf in der Nähe von Garoua. Sie ist 42 Jahre alt und hat sechs Kinder. Jeden Morgen steht sie um fünf Uhr auf, weil sie vor der Hitze auf dem Feld arbeiten muss. Dort pflanzt sie Mais und Erdnüsse. Obwohl sie sehr hart arbeitet, verdient sie wenig Geld, denn die Preise auf dem Markt sind niedrig. Aminatou kann nicht lesen, weil ihre Eltern sie nicht zur Schule geschickt haben. Deshalb will sie, dass ihre Töchter eine gute Schulbildung bekommen. Ihre älteste Tochter Mariam besucht jetzt das Gymnasium in Garoua. Sie will Ärztin werden, obwohl das Studium teuer ist. Aminatou verkauft deshalb jeden Samstag Gemüse auf dem Markt, um das Schulgeld zu bezahlen. Am Abend hilft sie Mariam auch im Haushalt, damit das Mädchen Zeit zum Lernen hat. Die Frauen im Dorf sagen: „Aminatou ist ein Beispiel für uns alle.“
\end{textebox}

\textbf{A. Réponds par \all{richtig} ou \all{falsch} et justifie ta réponse par une phrase du texte citée en allemand.}
\begin{enumerate}
\item Aminatou arbeitet nur am Nachmittag.
\item Aminatou verdient viel Geld auf dem Markt.
\item Aminatou ist nie zur Schule gegangen.
\item Mariam will Lehrerin werden.
\item Aminatou verkauft Gemüse, um die Ausbildung ihrer Tochter zu finanzieren.
\end{enumerate}

\textbf{B. Réponds en allemand par des phrases complètes.}
\begin{enumerate}
\item Warum steht Aminatou so früh auf?
\item Was will Aminatou für ihre Töchter?
\item Warum hilft Aminatou ihrer Tochter im Haushalt?
\end{enumerate}

\textbf{C. Questions de langue.}
\begin{enumerate}
\item Souligne dans le texte deux subordonnées introduites par \all{weil} et explique la place du verbe.
\item Relève dans le texte un verbe de modalité conjugué et conjugue-le au présent à toutes les personnes.
\item Trouve dans le texte le pluriel de \all{die Tochter} et donne l'article et le pluriel de : \all{das Dorf}, \all{der Markt}.
\end{enumerate}
\end{exercice}

\begin{exercice}[Thème et version type Bac]
\textbf{A. Thème.} Traduis en allemand les phrases suivantes. Utilise les connecteurs et les verbes de modalité de la leçon.
\begin{enumerate}
\item Ma mère travaille beaucoup parce qu'elle doit nourrir la famille.
\item Bien qu'elle soit fatiguée, elle va au marché tous les jours.
\item Les femmes au village ne peuvent pas aller à l'école, c'est pourquoi elles veulent une bonne éducation pour leurs filles.
\item Mon père aide ma mère dans le ménage, car il croit à l'égalité des droits.
\item Ma sœur veut devenir médecin, mais elle doit d'abord étudier beaucoup.
\end{enumerate}

\textbf{B. Version.} Traduis en français le paragraphe suivant.

\begin{textebox}[Zum Übersetzen]
Viele Frauen in der Stadt haben heute einen Beruf, weil sie Geld für ihre Familie verdienen müssen. Obwohl sie den ganzen Tag im Büro arbeiten, müssen sie am Abend auch den Haushalt machen. Die Männer sollen mehr helfen, denn die Gleichberechtigung beginnt zu Hause. Manche Frauen wollen deshalb nicht mehr heiraten, weil sie frei leben wollen.
\end{textebox}
\end{exercice}

\begin{exercice}[Production écrite type Bac]
Rédige un texte en allemand d'environ \textbf{10 lignes} (80 à 100 mots) sur le sujet suivant :

\begin{center}
\textbf{\all{Die Rolle der Frau in meiner Familie}}
\end{center}

Respecte le plan imposé :
\begin{enumerate}
\item \textbf{Einleitung :} présente la femme de ta famille dont tu parles (ta mère, ta grand-mère, ta tante...) : qui est-elle, où vit-elle ?
\item \textbf{Hauptteil :} décris son travail (à la maison, au champ, au marché, dans un bureau...) et explique pourquoi elle est importante pour la famille. Utilise au moins deux connecteurs parmi \all{weil}, \all{denn}, \all{deshalb}, \all{obwohl}.
\item \textbf{Schluss :} donne ton opinion : la situation de la femme doit-elle changer ? Utilise au moins un verbe de modalité (\all{müssen}, \all{sollen}, \all{können}, \all{wollen}, \all{dürfen}).
\end{enumerate}

\emph{Attention : verbe en 2e position dans les phrases principales, verbe conjugué à la fin des subordonnées, majuscule à tous les noms.}
\end{exercice}



\newpage

\coursec{Corrigés des exercices}

\coursec{Corrigés des exercices — Leçon AL01 : La situation de la femme}

\courssub{Exercice 1 — QCM : la grammaire de la leçon}
\begin{corrige}[Exercice 1]
\begin{enumerate}
\item \textbf{c) \all{obwohl}} : \all{Die Frau arbeitet viel, obwohl sie auch den Haushalt macht.} « La femme travaille beaucoup, bien qu'elle fasse aussi le ménage. » La relation est une \cle{opposition} (concession) : \all{obwohl} est une \cle{conjonction de subordination}, le verbe conjugué (\all{macht}) va en \cle{fin de subordonnée}.
\item \textbf{b) \all{weil}} : \all{Mama bleibt zu Hause, weil sie krank ist.} « Maman reste à la maison parce qu'elle est malade. » L'énoncé exige le verbe en fin de phrase : seule une conjonction de subordination convient. \all{denn} aurait imposé l'ordre \all{denn sie ist krank} (verbe en 2e position).
\item \textbf{c) \all{deshalb}} : \all{Sie hat kein Geld, deshalb kann sie nicht auf den Markt gehen.} « Elle n'a pas d'argent, c'est pourquoi elle ne peut pas aller au marché. » \all{deshalb} est un \cle{adverbe de liaison} (conjonctif) : il occupe la 1re position, le verbe reste en 2e position et le sujet suit le verbe (\cle{inversion}).
\item \textbf{c) en fin de phrase} : dans toute subordonnée introduite par une conjonction de subordination (\all{weil}, \all{obwohl}, \all{dass}...), le verbe conjugué se place en \cle{dernière position} : \all{Obwohl sie müde \textbf{ist}, arbeitet sie auf dem Feld.}
\item \textbf{b) \all{die Frauen}} : \all{die Frau, die Frauen}. Pluriel irrégulier en \all{-en}, à apprendre par cœur. \all{das Fräulein} (« mademoiselle ») est un autre mot, aujourd'hui vieilli et à éviter.
\end{enumerate}
\end{corrige}

\courssub{Exercice 2 — Vrai ou faux ? Justifie ta réponse}
\begin{corrige}[Exercice 2]
\begin{enumerate}
\item \all{Falsch}. \all{denn} (« car ») est une conjonction de \cle{coordination} : le verbe conjugué reste en 2e position. Exemple : \all{Sie bleibt zu Hause, denn sie ist krank.}
\item \all{Richtig}. \all{deshalb} occupe la 1re position et provoque l'\cle{inversion} du sujet : \all{Sie ist krank, deshalb \textbf{bleibt sie} zu Hause.}
\item \all{Falsch}. \all{obwohl} (« bien que ») exprime une \cle{opposition} (concession), pas une cause. La cause s'exprime avec \all{weil} (subordination) ou \all{denn} (coordination).
\item \all{Richtig}. \all{können} (présent) : \all{ich kann, du kannst, er/sie/es kann, wir können, ihr könnt, sie/Sie können}. La 1re et la 3e personne du singulier sont identiques, sans Umlaut.
\item \all{Falsch}. \all{lesen} est un verbe fort à changement de voyelle \all{e} $\rightarrow$ \all{ie} : \all{er \textbf{liest}} (comme \all{du liest}). La forme \all{lest} correspond à la 2e personne du pluriel : \all{ihr lest}.
\end{enumerate}
\end{corrige}

\courssub{Exercice 3 — Vocabulaire : la femme, la famille, le travail}
\begin{corrige}[Exercice 3]
\begin{enumerate}
\item \all{auf dem \textbf{Feld}} : « au champ » (\all{das Feld, die Felder} ; \all{auf dem Feld} = datif après \all{auf} pour indiquer le lieu où l'on est).
\item \all{keine \textbf{Schulbildung}} : « pas de scolarisation » (\all{die Schulbildung, -en} ; accusatif féminin après \all{haben}, d'où \all{keine}).
\item \all{einen guten \textbf{Beruf}} : « un bon métier » (\all{der Beruf, -e} ; accusatif masculin : \all{einen guten}).
\item \all{den \textbf{Haushalt}} : « le ménage, les tâches ménagères » (\all{der Haushalt, -e} ; accusatif masculin après \all{teilen}).
\item \all{\textbf{Die Gleichberechtigung} bedeutet ...} : « L'égalité des droits signifie ... » (\all{die Gleichberechtigung, -en} ; sujet au nominatif, majuscule en début de phrase).
\end{enumerate}

\textbf{Récapitulatif du lexique (articles, pluriels, traduction) :}
\begin{center}
\begin{tabular}{|l|l|l|l|}
\hline
\rowcolor{popblueL} \textbf{Mot} & \textbf{Article / Pluriel} & \textbf{Traduction} & \textbf{Exemple / Cas} \\ \hline
\all{Frau} & \all{die Frau, die Frauen} & la femme & \all{die Frau} (nom./akk.) \\ \hline
\all{Gleichberechtigung} & \all{die Gleichberechtigung, -en} & l'égalité des droits & \all{die Gleichberechtigung} (nom.) \\ \hline
\all{Schulbildung} & \all{die Schulbildung, -en} & la scolarisation & \all{keine Schulbildung} (akk.) \\ \hline
\all{Beruf} & \all{der Beruf, -e} & le métier & \all{einen guten Beruf} (akk.) \\ \hline
\all{Haushalt} & \all{der Haushalt, -e} & le ménage & \all{den Haushalt} (akk.) \\ \hline
\all{Arbeit} & \all{die Arbeit, -en} & le travail & \all{auf dem Feld arbeiten} \\ \hline
\all{Ehe} & \all{die Ehe, -n} & le mariage & \all{in die Ehe eingehen} \\ \hline
\end{tabular}
\end{center}
\end{corrige}

\courssub{Exercice 4 — Conjugaison : le présent (verbes faibles, forts, modaux)}
\begin{corrige}[Exercice 4]
\begin{enumerate}
\item Amina \all{\textbf{arbeitet}} jeden Tag von 6 bis 17 Uhr. (verbe faible : radical \all{arbeit-} + \all{-et} à la 3e pers. sg., radical en \all{-t})
\item Sie \all{\textbf{nimmt}} das Gemüse und \all{\textbf{trägt}} es zum Markt. (verbes forts : \all{nehmen} $\rightarrow$ \all{nimmt} [e$\rightarrow$ie] ; \all{tragen} $\rightarrow$ \all{trägt} [a$\rightarrow$ä])
\item Ihre Tochter \all{\textbf{hilft}} ihr am Wochenende. (\all{helfen} $\rightarrow$ \all{hilft} [e$\rightarrow$i] ; \cle{datif} après \all{helfen} : \all{sie hilft \textbf{ihr}})
\item Amina \all{\textbf{kann}} nicht lesen, aber ihre Tochter \all{\textbf{will}} Ärztin werden. (verbes de modalité : \all{ich/er kann}, \all{ich/er will} ; l'infinitif \all{lesen}/\all{werden} en fin de phrase)
\item Die Frauen im Dorf \all{\textbf{müssen}} viel arbeiten, aber sie \all{\textbf{dürfen}} nicht über das Geld entscheiden. (au pluriel : forme identique à l'infinitif)
\item Mein Vater \all{\textbf{spricht}} oft über die Gleichberechtigung. (\all{sprechen} $\rightarrow$ \all{spricht} [e$\rightarrow$i])
\end{enumerate}

\textbf{Tableau de conjugaison (présent) — verbes de la leçon :}
\begin{center}
\begin{tabular}{|l|l|l|l|l|l|l|}
\hline
\rowcolor{popblueL} \textbf{Verbe} & \textbf{ich} & \textbf{du} & \textbf{er/sie/es} & \textbf{wir} & \textbf{ihr} & \textbf{sie/Sie} \\ \hline
\all{arbeiten} (faible) & arbeite & arbeitest & \textbf{arbeitet} & arbeiten & arbeitet & arbeiten \\ \hline
\all{nehmen} (fort, e$\rightarrow$ie) & nehme & nimmst & \textbf{nimmt} & nehmen & nehmt & nehmen \\ \hline
\all{tragen} (fort, a$\rightarrow$ä) & trage & trägst & \textbf{trägt} & tragen & tragt & tragen \\ \hline
\all{helfen} (fort, e$\rightarrow$i) & helfe & hilfst & \textbf{hilft} & helfen & helft & helfen \\ \hline
\all{können} (modal) & kann & kannst & \textbf{kann} & können & könnt & können \\ \hline
\all{wollen} (modal) & will & willst & \textbf{will} & wollen & wollt & wollen \\ \hline
\all{müssen} (modal) & muss & musst & \textbf{muss} & müssen & müsst & müssen \\ \hline
\all{dürfen} (modal) & darf & darfst & \textbf{darf} & dürfen & dürft & dürfen \\ \hline
\all{sprechen} (fort, e$\rightarrow$i) & spreche & sprichst & \textbf{spricht} & sprechen & sprecht & sprechen \\ \hline
\end{tabular}
\end{center}
\end{corrige}

\courssub{Exercice 5 — Remets les mots dans l'ordre (syntaxe)}
\begin{corrige}[Exercice 5]
\begin{enumerate}
\item \all{Die Frau arbeitet jeden Tag auf dem Feld.} « La femme travaille tous les jours au champ. » (Sujet — Verbe en 2e position — Compléments : temps, lieu)
\item \all{Weil sie müde ist, geht sie früh ins Bett.} « Parce qu'elle est fatiguée, elle se couche tôt. » (Subordonnée en tête : verbe \all{ist} à la fin ; principale qui suit : verbe \all{geht} en 1re position, sujet \all{sie} en 2e)
\item \all{Viele Mädchen können nicht zur Schule gehen, weil ihre Eltern arm sind.} « Beaucoup de filles ne peuvent pas aller à l'école parce que leurs parents sont pauvres. » (Principale : verbe \all{können} 2e pos. ; subordonnée \all{weil} : verbe \all{sind} à la fin)
\item \all{Obwohl sie viel arbeitet, verdient sie wenig Geld.} « Bien qu'elle travaille beaucoup, elle gagne peu d'argent. » (On accepte aussi : \all{Sie verdient wenig Geld, obwohl sie viel arbeitet.})
\item \all{Deshalb will meine Schwester in Deutschland studieren.} « C'est pourquoi ma sœur veut étudier en Allemagne. » (\all{deshalb} 1re pos., verbe \all{will} 2e, sujet \all{meine Schwester} 3e, infinitif \all{studieren} à la fin)
\end{enumerate}
\end{corrige}

\courssub{Exercice 6 — Relie les phrases avec le bon connecteur}
\begin{corrige}[Exercice 6]
\begin{enumerate}
\item \all{Fatou lernt Deutsch, weil sie in Deutschland studieren will.} « Fatou apprend l'allemand parce qu'elle veut étudier en Allemagne. » (\all{weil} : verbe conjugué \all{will} à la fin, après l'infinitif \all{studieren} — structure double infinitif)
\item \all{Die Frau arbeitet den ganzen Tag, obwohl sie nie müde ist.} « La femme travaille toute la journée, bien qu'elle ne soit jamais fatiguée. » (On accepte aussi : \all{Obwohl die Frau den ganzen Tag arbeitet, ist sie nie müde.})
\item \all{Es regnet stark, deshalb gehen die Frauen nicht auf das Feld.} « Il pleut fort, c'est pourquoi les femmes ne vont pas au champ. » (\all{deshalb} + inversion : \all{gehen die Frauen})
\item \all{Meine Mutter kann nicht lesen, denn sie ist nie zur Schule gegangen.} « Ma mère ne sait pas lire, car elle n'est jamais allée à l'école. » (\all{denn} : coordination, verbe \all{ist} en 2e position)
\item \all{Die Mädchen haben wenig Zeit, weil sie im Haushalt helfen müssen.} « Les filles ont peu de temps parce qu'elles doivent aider au ménage. » (\all{weil} : à la fin, l'infinitif \all{helfen} précède le verbe conjugué \all{müssen})
\end{enumerate}
\end{corrige}

\courssub{Exercice 7 — Texte à trous : la vie de Maman Rose}
\begin{corrige}[Exercice 7]
\begin{enumerate}
\item \all{Maman Rose \textbf{steht} jeden Morgen um fünf Uhr \textbf{auf}.} (verbe à particule séparable \all{aufstehen} : \all{steht} en 2e position, \all{auf} en fin de phrase)
\item \all{Zuerst \textbf{kocht} sie das Frühstück für die Familie.} (\all{zuerst} en 1re position $\rightarrow$ \cle{inversion} du sujet)
\item \all{Dann geht sie auf den Markt und \textbf{verkauft} Gemüse und Früchte.}
\item \all{Sie \textbf{verdient} nicht viel Geld, ...} « Elle ne gagne pas beaucoup d'argent. »
\item \all{... aber sie \textbf{muss} ihre Kinder ernähren.} « ... mais elle doit nourrir ses enfants. » (\all{müssen} : \all{er/sie muss}, sans \all{-t} et sans Umlaut)
\item \all{Ihre Tochter Sandrine \textbf{hilft} ihr am Wochenende.} (\all{helfen} + \cle{datif} : \all{ihr})
\item \all{Maman Rose \textbf{kann} nicht lesen und schreiben, ...}
\item \all{... aber Sandrine \textbf{will} später Lehrerin werden.} « ... mais Sandrine veut devenir institutrice plus tard. »
\end{enumerate}

\textbf{Rappel : verbe à particule séparable \all{aufstehen} (présent) :}
\begin{center}
\begin{tabular}{|l|l|}
\hline
\rowcolor{popblueL} \textbf{Personne} & \textbf{Forme} \\ \hline
ich & \all{stehe ... auf} \\ \hline
du & \all{stehst ... auf} \\ \hline
er/sie/es & \all{steht ... auf} \\ \hline
wir & \all{stehen ... auf} \\ \hline
ihr & \all{steht ... auf} \\ \hline
sie/Sie & \all{stehen ... auf} \\ \hline
\end{tabular}
\end{center}
\end{corrige}

\courssub{Exercice 8 — Appariements : mots et définitions}
\begin{corrige}[Exercice 8]
\begin{enumerate}
\item \all{die Frau} — \textbf{b)} la femme (adulte)
\item \all{das Mädchen} — \textbf{d)} la fille (jeune)
\item \all{der Haushalt} — \textbf{a)} le ménage, les tâches domestiques
\item \all{der Beruf} — \textbf{c)} le métier, la profession
\item \all{die Gleichberechtigung} — \textbf{e)} l'égalité des droits
\item \all{die Schulbildung} — \textbf{f)} la scolarisation, l'instruction scolaire
\end{enumerate}

\textbf{À retenir absolument (articles, pluriels, genre) :}
\begin{center}
\begin{tabular}{|l|l|l|}
\hline
\rowcolor{popblueL} \textbf{Nom} & \textbf{Article / Pluriel} & \textbf{Traduction} \\ \hline
\all{Gleichberechtigung} & \all{die Gleichberechtigung, -en} & l'égalité des droits \\ \hline
\all{Haushalt} & \all{der Haushalt, -e} & le ménage \\ \hline
\all{Arbeit} & \all{die Arbeit, -en} & le travail \\ \hline
\all{Schulbildung} & \all{die Schulbildung, -en} & la scolarisation \\ \hline
\all{Beruf} & \all{der Beruf, -e} & le métier \\ \hline
\all{Ehe} & \all{die Ehe, -n} & le mariage \\ \hline
\end{tabular}
\end{center}
\end{corrige}

\courssub{Exercice 9 — Compréhension écrite type Bac}
\begin{corrige}[Exercice 9]
\textbf{A. Richtig oder falsch ?} (5 $\times$ 2 pts = 10 pts : 1 pt pour le jugement, 1 pt pour la citation justificative exacte)
\begin{enumerate}
\item \all{Falsch}. Le texte dit : \all{„Jeden Morgen steht sie um fünf Uhr auf, weil sie vor der Hitze auf dem Feld arbeiten muss.“} Elle travaille donc dès le matin, pas seulement l'après-midi.
\item \all{Falsch}. \all{„Obwohl sie sehr hart arbeitet, verdient sie wenig Geld.“}
\item \all{Richtig}. \all{„Aminatou kann nicht lesen, weil ihre Eltern sie nicht zur Schule geschickt haben.“}
\item \all{Falsch}. \all{„Sie will Ärztin werden.“} Mariam veut devenir médecin, pas institutrice.
\item \all{Richtig}. \all{„Aminatou verkauft deshalb jeden Samstag Gemüse auf dem Markt, um das Schulgeld zu bezahlen.“}
\end{enumerate}

\textbf{B. Réponses en allemand (modèles).} (3 $\times$ 2 pts = 6 pts)
\begin{enumerate}
\item \all{Aminatou steht so früh auf, weil sie vor der Hitze auf dem Feld arbeiten muss.}
\item \all{Sie will, dass ihre Töchter eine gute Schulbildung bekommen.} (\all{dass} + verbe à la fin)
\item \all{Sie hilft ihrer Tochter im Haushalt, damit das Mädchen Zeit zum Lernen \textbf{habe}.} (\all{damit} + \cle{Konjunktiv I} \all{habe} ; l'indicatif \all{hat} est toléré au Bac mais le subjonctif est la norme)
\end{enumerate}

\textbf{C. Questions de langue.} (4 pts)
\begin{enumerate}
\item \all{..., weil sie vor der Hitze auf dem Feld arbeiten \textbf{muss}.} et \all{..., weil ihre Eltern sie nicht zur Schule geschickt \textbf{haben}.} Dans une subordonnée introduite par \all{weil}, le verbe conjugué se place en \textbf{dernière position}, après l'infinitif (\all{muss}) ou le participe passé (\all{haben}) s'il y en a un.
\item Exemple : \all{muss} (de \all{müssen}). Présent : 
\begin{center}
\begin{tabular}{|l|l|l|l|l|l|l|}
\hline
\rowcolor{popblueL} ich & du & er/sie/es & wir & ihr & sie/Sie \\ \hline
\all{muss} & \all{musst} & \all{muss} & \all{müssen} & \all{müsst} & \all{müssen} \\ \hline
\end{tabular}
\end{center}
(Autres modaux dans le texte : \all{kann}, \all{will}.)
\item Pluriels dans le texte : \all{die \textbf{Töchter}} (\all{die Tochter, die Töchter}) ; \all{das Dorf, die \textbf{Dörfer}} ; \all{der Markt, die \textbf{Märkte}}.
\end{enumerate}

\textbf{Points de vigilance} : ne pas traduire mot à mot ; citer la phrase allemande exacte (guillemets „…“) pour justifier ; vérifier l'accord sujet--verbe dans les réponses rédigées.
\end{corrige}

\courssub{Exercice 10 — Thème et version type Bac}
\begin{corrige}[Exercice 10]
\textbf{A. Thème (corrigé).} (5 $\times$ 2 pts = 10 pts : 1 pt pour le connecteur et la place du verbe, 1 pt pour le vocabulaire et la conjugaison)
\begin{enumerate}
\item \all{Meine Mutter arbeitet viel, weil sie die Familie ernähren muss.} (verbe conjugué \all{muss} à la fin, après l'infinitif \all{ernähren})
\item \all{Obwohl sie müde ist, geht sie jeden Tag auf den Markt.} (subordonnée en tête : la principale commence par le verbe \all{geht})
\item \all{Die Frauen im Dorf können nicht zur Schule gehen, deshalb wollen sie eine gute Schulbildung für ihre Töchter.} (\all{deshalb} en position 1 de la principale $\rightarrow$ inversion : \all{wollen sie})
\item \all{Mein Vater hilft meiner Mutter im Haushalt, denn er glaubt an die Gleichberechtigung.} (\all{helfen} + \cle{datif} : \all{meiner Mutter} ; \all{glauben an} + \cle{accusatif} ; \all{denn} : verbe en 2e position \all{glaubt})
\item \all{Meine Schwester will Ärztin werden, aber sie muss zuerst viel lernen.} (« devenir » = \all{werden} ; « étudier beaucoup » $\rightarrow$ \all{viel lernen})
\end{enumerate}

\textbf{B. Version (corrigé).} (10 pts)

\all{„Heute üben viele Frauen in der Stadt einen Beruf aus, weil sie Geld für ihre Familie verdienen müssen. Obwohl sie den ganzen Tag im Büro arbeiten, müssen sie abends auch noch den Haushalt machen. Die Männer sollten mehr helfen, denn die Gleichberechtigung fängt zu Hause an. Deshalb wollen manche Frauen nicht mehr heiraten, weil sie frei leben wollen.“}

\textbf{Points de vigilance} : 
\begin{itemize}
\item \all{der Beruf} = le métier (faux ami : \all{die Berufe} au pluriel, pas « la berufe ») ; \all{einen Beruf ausüben} = exercer un métier.
\item \all{sollen} + infinitif = « devraient » (recommandation morale).
\item \all{heiraten} = se marier (verbe transitif, pas de \cle{sich}).
\item \all{manche} = certain(e)s (indéfini, décliné comme \all{dieser} : \all{manche Frauen}).
\item \all{frei leben} = vivre librement.
\end{itemize}
\end{corrige}

\courssub{Exercice 11 — Production écrite type Bac}
\begin{corrige}[Exercice 11]
\textbf{Barème indicatif} (20 pts) :
\begin{itemize}
\item Respect du plan en trois parties (introduction, développement, conclusion) : 4 pts
\item Emploi d'au moins deux connecteurs corrects (\all{weil}, \all{obwohl}, \all{denn}, \all{deshalb}) : 4 pts
\item Emploi d'au moins un verbe de modalité (\all{sollen}, \all{müssen}, \all{können}, \all{wollen}) : 2 pts
\item Ordre des mots et place du verbe (principal, subordonnée, inversion) : 4 pts
\item Vocabulaire du thème (femme, travail, ménage, éducation, droits) : 3 pts
\item Orthographe, majuscules (noms), Umlaute (ä, ö, ü), ß : 3 pts
\end{itemize}

\textbf{Modèle de rédaction} (environ 90 mots — à adapter avec tes propres idées) :

\begin{textebox}[Modèle de production écrite]
\all{Meine Mutter heißt Jeanne und lebt mit uns in Douala. Sie ist 45 Jahre alt und Verkäuferin auf dem Markt. Jeden Morgen steht sie um fünf Uhr auf, weil sie früh auf dem Markt sein muss. Dort verkauft sie Gemüse und Früchte. Obwohl sie den ganzen Tag hart arbeitet, findet sie am Abend noch Zeit für die Familie, denn sie kocht und hilft uns bei den Hausaufgaben. Meine Mutter ist sehr wichtig für uns, weil sie Geld für die Schule verdient. Deshalb respektieren wir sie alle. Ich finde: Die Männer sollen mehr im Haushalt helfen, und die Frauen müssen die gleichen Rechte wie die Männer haben. Dann kann unsere Familie noch glücklicher sein.}
\end{textebox}

\textbf{Analyse du modèle (points de contrôle pour l'auto-évaluation) :}
\begin{itemize}
\item \textbf{Introduction} : Qui ? (\all{Meine Mutter Jeanne}), Où ? (\all{Douala}), Métier ? (\all{Verkäuferin}).
\item \textbf{Connecteurs variés et syntaxe} :
      \begin{itemize}
      \item \all{weil ... muss} (verbe à la fin, double infinitif)
      \item \all{Obwohl ... arbeitet, findet sie ...} (subordonnée en tête $\rightarrow$ verbe principal en 1re position)
      \item \all{denn ... kocht} (coordination, verbe 2e position)
      \item \all{weil ... verdient} (subordonnée, verbe à la fin)
      \item \all{Deshalb respektieren wir ...} (adverbe de liaison, inversion)
      \end{itemize}
\item \textbf{Verbes de modalité en conclusion} : \all{sollen helfen}, \all{müssen haben}, \all{kann ... sein} (infinitifs en fin de phrase).
\item \textbf{Majuscules} : \all{Mutter, Markt, Gemüse, Früchte, Familie, Abend, Hausaufgaben, Schule, Geld, Männer, Haushalt, Frauen, Rechte, Familie}.
\item \textbf{Nombre de mots} : $\approx$ 95 mots (dans la fourchette 80--100).
\end{itemize}
\end{corrige}

\newpage

\coursec{Fiche bilan}

\begin{popfigure}
\begin{tikzpicture}[mindmap, grow cyclic, every node/.style={concept, text=white, font=\small},
  root concept/.append style={concept color=popblue, font=\bfseries\small, minimum size=2.6cm},
  level 1/.append style={level distance=3.4cm, sibling angle=60},
  level 1 concept/.append style={minimum size=2.1cm, font=\footnotesize}]
\node[root concept, align=center]{Die Situation der Frau\\ (Stadt und Dorf)}
  child[concept color=poporange]{ node[align=center]{Texte\\ verstehen} }
  child[concept color=popgreen]{ node[align=center]{Wortschatz\\ die Frau, der Beruf,\\ der Haushalt} }
  child[concept color=poppurple]{ node[align=center]{Wortstellung\\ Verb an 2. Stelle\\ Nebensatz: Verb am Ende} }
  child[concept color=poppink]{ node[align=center]{Konnektoren\\ weil, denn,\\ deshalb, obwohl} }
  child[concept color=popgold]{ node[align=center]{Konjugation\\ Präsens +\\ Modalverben} }
  child[concept color=popblue!70]{ node[align=center]{Produktion\\ Fragen, Zusammen-\\ fassung, kurzer Text} };
\end{tikzpicture}
\legende{Carte mentale de la leçon : la situation de la femme en ville et au village}
\end{popfigure}

\begin{aretenir}[L'essentiel de la leçon]
\textbf{1. Lexique clé (à connaître avec article et pluriel)}

\begin{center}
\begin{tabularx}{\linewidth}{|X|X|}\hline
\rowcolor{popblueL} \textbf{Deutsch} & \textbf{Français} \\ \hline
die Frau, -en & la femme \\ \hline
die Gleichberechtigung, -en & l'égalité des droits \\ \hline
die Schulbildung, -en & l'instruction, la scolarisation \\ \hline
der Haushalt, -e & le ménage, les tâches domestiques \\ \hline
der Beruf, -e & le métier, la profession \\ \hline
die Arbeit auf dem Feld & le travail aux champs \\ \hline
das Dorf, „-er & le village \\ \hline
die Stadt, „-e & la ville \\ \hline
das Mädchen, - & la jeune fille \\ \hline
die Mutter, „- & la mère \\ \hline
verdienen (verdient, verdiente, hat verdient) & gagner (de l'argent) \\ \hline
sich kümmern um + Akk. & s'occuper de \\ \hline
aufstehen (steht auf, stand auf, ist aufgestanden) & se lever \\ \hline
\end{tabularx}
\end{center}

\textbf{2. Grammaire à connaître par cœur}

\begin{itemize}
\item \cle{Verbe en 2\up{e} position} dans la phrase principale : \all{Die Frau \textbf{arbeitet} heute auf dem Feld.} « La femme travaille aujourd'hui aux champs. » / \all{Heute \textbf{arbeitet} die Frau auf dem Feld.} « Aujourd'hui, la femme travaille aux champs. » (quand un complément ouvre la phrase, le sujet passe après le verbe ; le verbe ne recule jamais en 3\up{e} position).
\item \cle{Subordonnée} : le verbe conjugué va \textbf{à la fin} : \all{Die Mädchen bleiben zu Hause, \textbf{weil} sie der Mutter im Haushalt \textbf{helfen müssen}.} « Les filles restent à la maison parce qu'elles doivent aider la mère dans les tâches domestiques. » (avec un modal, l'infinitif précède le verbe conjugué : \all{müssen} se place après \all{helfen}).
\item \cle{weil} + subordonnée (verbe final) ; \cle{denn} + principale (verbe en 2\up{e} position, pas d'inversion) ; \cle{deshalb} + inversion (verbe puis sujet) ; \cle{obwohl} + subordonnée (concession, verbe final).
\item \cle{Présent} : verbes faibles (\all{machen : ich mache, du machst, er macht}) ; verbes forts à changement de voyelle (\all{lesen : er liest ; fahren : du fährst ; sprechen : er spricht}).
\item \cle{Modalverben} au présent : \all{ich kann, du kannst, er kann ; ich muss, du musst, er muss ; ich will, du willst, er will ; ich darf, du darfst, er darf ; ich soll, du sollst, er soll ; ich mag, du magst, er mag / ich möchte} + infinitif \textbf{à la fin} : \all{Die Frau \textbf{muss} früh aufstehen und Wasser \textbf{holen}.} « La femme doit se lever tôt et aller chercher de l'eau. »
\end{itemize}

\textbf{3. Tableau des verbes de modalité au présent}

\begin{center}
\begin{tabular}{|l|l|l|l|l|}\hline
\rowcolor{popblueL} \textbf{Personne} & \textbf{können} & \textbf{müssen} & \textbf{wollen} & \textbf{dürfen} \\ \hline
ich & kann & muss & will & darf \\ \hline
du & kannst & musst & willst & darfst \\ \hline
er/sie/es & kann & muss & will & darf \\ \hline
wir & können & müssen & wollen & dürfen \\ \hline
ihr & könnt & müsst & wollt & dürft \\ \hline
sie/Sie & können & müssen & wollen & dürfen \\ \hline
\end{tabular}
\end{center}

\textbf{4. Méthodes express}

\begin{itemize}
\item \cle{Compréhension} : 1) lire le titre et deviner le thème ; 2) lecture globale (qui ? où ? quel problème ?) ; 3) lecture détaillée avec soulignage des mots-clés ; 4) répondre par des phrases complètes en allemand.
\item \cle{Résumé} : reprendre les idées principales avec \all{zuerst, dann, außerdem, schließlich}, sans recopier le texte.
\item \cle{Production} : 5 à 8 phrases simples, un connecteur par phrase, vérifier la place du verbe avant de rendre.
\end{itemize}
\end{aretenir}

\begin{attention}[Checklist avant le Bac]
\begin{itemize}
\item[$\square$] Je connais le vocabulaire du thème avec l'article et le pluriel de chaque nom.
\item[$\square$] Je mets une majuscule à tous les noms allemands (\all{die Frau, das Dorf, der Beruf}).
\item[$\square$] Je place toujours le verbe conjugué en 2\up{e} position dans la phrase principale.
\item[$\square$] Je place le verbe conjugué à la fin après \all{weil} et \all{obwohl}.
\item[$\square$] Je fais l'inversion sujet-verbe après \all{deshalb} et après un complément en tête de phrase.
\item[$\square$] Je distingue \all{weil} (subordonnée) et \all{denn} (principale, pas d'inversion).
\item[$\square$] Je conjugue sans erreur les verbes forts au présent (\all{er liest, du fährst, er spricht}).
\item[$\square$] Je place l'infinitif à la fin avec les verbes de modalité (\all{Sie muss Wasser holen.}).
\item[$\square$] Je réponds aux questions de compréhension par des phrases complètes en allemand.
\item[$\square$] Je rédige un court paragraphe avec au moins trois connecteurs différents.
\item[$\square$] Je me relis : place du verbe, accords, majuscules, Umlaute.
\end{itemize}
\end{attention}

\findecours

\end{document}
